
\SGARefTarget{sga3:xxii:chapter:xxii}\chapter*{Exposé XXII\\
Reductive Groups: Splittings, Subgroups, Quotient Groups}
\addcontentsline{toc}{chapter}{Exposé XXII. Reductive Groups:
Splittings, Subgroups, Quotient Groups}
\markboth{EXPOSÉ XXII. REDUCTIVE GROUPS}
  {EXPOSÉ XXII. REDUCTIVE GROUPS}

\begin{center}
  \textit{by Michel Demazure}
\end{center}

\begingroup
\renewcommand{\thefootnote}{0}
\addtocounter{footnote}{-1}
\endgroup
This Exposé consists of two parts. The first (\S\S1--\SGARefLink{sga3:xxii:subsection:5-5}{5.5}) gathers the
technical results needed for the proofs of the uniqueness and
existence theorems. The second (\S5.6 to the end) will not be used in
that proof; the end of \S5 will be used in particular in Exposé XXVI,
which is devoted to parabolic subgroups; \S6 establishes, in the
scheme-theoretic setting, the classical results concerning the
derived group of a reductive group.

\SGARefTarget{sga3:xxii:section:1}\section*{1. Roots and Coroots.
Split Groups and Root Data}
\addcontentsline{toc}{section}{1. Roots and Coroots. Split Groups and
Root Data}

\SGARefTarget{sga3:xxii:thm:1.1}\paragraph{\SGARefLink{sga3:xxii:thm:1.1}{Theorem 1.1}.}
\textit{Let \(S\) be a scheme, \(G\) a reductive \(S\)-group,
\(T\) a maximal torus of \(G\), and \(\alpha\) a root of \(G\)
relative to \(T\).}

\begin{enumerate}
\item[\SGAInlineRefTarget{sga3:xxii:thm:1.1:item:i}(i)] \textit{There exists a unique \(T\)-equivariant
homomorphism of group schemes}
\[
  \exp_{\alpha}:W(\mathfrak g^{\alpha})\longrightarrow G
\]
\textit{which induces on Lie algebras the canonical morphism
\(\mathfrak g^{\alpha}\to\mathfrak g\). This homomorphism is a closed
immersion. The corresponding homomorphism}
\[
  T\mathbin{\cdot_{\alpha}}W(\mathfrak g^{\alpha})
  \longrightarrow G
\]
\textit{is likewise a closed immersion.}

\textit{If
\(p_{\alpha}:\mathbb G_{a,S}\to G\) is a monomorphism normalized by
\(T\), with multiplier \(\alpha\), there exists a unique
\(X_{\alpha}\in\Gamma(S,\mathfrak g^{\alpha})^{\times}\) such that
\(p_{\alpha}(x)=\exp_{\alpha}(xX_{\alpha})\); one has
\(\Lie(p_{\alpha})(1)=X_{\alpha}\), and the preceding two formulas
establish a bijective correspondence between
\(\Gamma(S,\mathfrak g^{\alpha})^{\times}\) and the set of
monomorphisms \(\mathbb G_{a,S}\to G\) normalized by \(T\), with
multiplier \(\alpha\).}

\item[(ii)] \textit{There exists a unique duality, denoted
\((X,Y)\longmapsto XY\),}
\[
  \mathfrak g^{\alpha}\otimes_{\mathcal O_S}\mathfrak g^{-\alpha}
  \xrightarrow{\;\sim\;}\mathcal O_S,
\]
\textit{and a unique homomorphism of group schemes}
\[
  \alpha^*:\mathbb G_{m,S}\longrightarrow T,
\]
\textit{such that formula \textup{(F)} of \SGARefLink{sga3:xx:2.1}{Exposé XX, 2.1} holds. One
has}
\[
  \alpha\circ\alpha^*=2,
  \qquad
  (-\alpha)^*=-\alpha^*,
\]
\textit{and \(\alpha^*\) is given by the formula of \SGARefLink{sga3:xx:2.7}{Exposé XX, 2.7}.}
\end{enumerate}

Indeed, a homomorphism normalized by \(T\), with multiplier \(\alpha\),
necessarily factors through the closed subgroup
\[
  Z_\alpha=\underline{\operatorname{Centr}}_G(T_\alpha)
\]
of \(G\) (cf. \SGARefLink{sga3:xix:paragraph:3-9}{Exposé XIX, 3.9}). Now \((Z_\alpha,T,\alpha)\) is an
elementary \(S\)-system (\SGARefLink{sga3:xx:1.4}{Exposé XX, 1.4}), and one is reduced to the
results of Exposé XX (\SGARefLink{sga3:xxii:def:1.5}{1.5}, \SGARefLink{sga3:xxii:prop:2.1}{2.1}, and \SGARefLink{sga3:xxii:section:5.9}{5.9}).

\SGARefTarget{sga3:xxii:rem:1.2}\paragraph{\SGARefLink{sga3:xxii:rem:1.2}{Remark 1.2}.}
Part \SGARefLink{sga3:xxii:thm:1.1:item:i}{(i)} of \SGARefLink{sga3:xxii:thm:1.1}{Theorem 1.1} remains valid if one assumes only that
\(\alpha\) is a character of \(T\), nontrivial on every fiber. Indeed,
there is then a decomposition
\[
  S=S'\amalg S'',
\]
such that \(\alpha|_{S'}\) is a root of \(G_{S'}\) relative to
\(T_{S'}\), and
\(\mathfrak g^\alpha|_{S''}=0\). If \(S=S'\), one is reduced to \SGARefLink{sga3:xxii:thm:1.1}{1.1};
if \(S=S''\), the result is trivial; the general case follows at once.

\SGARefTarget{sga3:xxii:not:1.3}\paragraph{\SGARefLink{sga3:xxii:not:1.3}{Notations 1.3}.}
As in Exposé XX, let \(U_\alpha\) denote the image of
\(W(\mathfrak g^\alpha)\); it is a closed subgroup of \(G\),
canonically endowed with a vector-group structure. We call it the
\emph{vector group associated with the root \(\alpha\)}. We say that
\(\alpha^*\) is the \emph{coroot associated with \(\alpha\)}. Sections
\[
  X_\alpha\in\Gamma(S,\mathfrak g^\alpha)
  \quad\text{and}\quad
  X_{-\alpha}\in\Gamma(S,\mathfrak g^{-\alpha})
\]
are said to be \emph{paired} if \(X_\alpha X_{-\alpha}=1\). Then
\(X_\alpha\in\Gamma(S,\mathfrak g^\alpha)^\times\), and likewise for
\(X_{-\alpha}\). The corresponding homomorphisms \(p_\alpha\) and
\(p_{-\alpha}\) are contragredient to one another in the sense of
\SGARefLink{sga3:xx:1.5}{Exposé XX, 1.5}, and one has
\[
  p_\alpha(x)p_{-\alpha}(y)
  =
  p_{-\alpha}\left(\frac{y}{1+xy}\right)
  \alpha^*(1+xy)
  p_\alpha\left(\frac{x}{1+xy}\right).
\]

\SGARefTarget{sga3:xxii:prop:1.4}\paragraph{\SGARefLink{sga3:xxii:prop:1.4}{Proposition 1.4}.}
\textit{Under the hypotheses of \SGARefLink{sga3:xxii:thm:1.1}{1.1}, let
\(w\in\underline{\operatorname{Norm}}_G(T)(S)\). Then}
\[
  \beta=\alpha\circ\operatorname{int}(w)^{-1}
  :T\longrightarrow\mathbb G_{m,S}
\]
\textit{is a root of \(G\) relative to \(T\),
\(\beta^*=\operatorname{int}(w)\circ\alpha^*\) is the corresponding
coroot, and the following diagram is commutative:}

\begin{center}
  \begin{tikzcd}[column sep=5.2em,row sep=3.2em]
  W(\mathfrak g^\alpha)
    \arrow[r,"\exp_\alpha"]
    \arrow[d,"\operatorname{Ad}(w)"']
  &
  G
    \arrow[d,"\operatorname{int}(w)"]
  \\
  W(\mathfrak g^\beta)
    \arrow[r,"\exp_\beta"]
  &
  G.
\end{tikzcd}

\end{center}

Trivial: transport of structure.

\SGARefTarget{sga3:xxii:def:1.5}\paragraph{\SGARefLink{sga3:xxii:def:1.5}{Definitions 1.5}.}
\textup{(a)} Under the hypotheses of \SGARefLink{sga3:xxii:thm:1.1}{1.1}, let \(s_\alpha\) denote the
automorphism of \(T\) defined by
\[
  s_\alpha(t)=t\mathbin{\cdot}\alpha^*(\alpha(t))^{-1}.
\]
Let \((\, ,\,)\) denote the canonical pairing
\[
  \underline{\operatorname{Hom}}_{S\text{-gr}}
    (\mathbb G_{m,S},T)
  \times
  \underline{\operatorname{Hom}}_{S\text{-gr}}
    (T,\mathbb G_{m,S})
  \longrightarrow
  \underline{\operatorname{Hom}}_{S\text{-gr}}
    (\mathbb G_{m,S},\mathbb G_{m,S})
  =\mathbb Z_S.
\]
Then \(s_\alpha\) acts on
\(\underline{\operatorname{Hom}}_{S\text{-gr}}
(T,\mathbb G_{m,S})\), respectively on
\(\underline{\operatorname{Hom}}_{S\text{-gr}}
(\mathbb G_{m,S},T)\), by the following formulas, where \(\chi\),
respectively \(u\), denotes an arbitrary section of the corresponding
scheme:
\[
  s_\alpha(\chi)=\chi-(\alpha^*,\chi)\alpha,
  \qquad
  s_\alpha(u)=u-(u,\alpha)\alpha^*.
\]
One has
\[
  s_\alpha\circ s_\alpha=\operatorname{id},
  \qquad
  s_{-\alpha}=s_\alpha.
\]

\textup{\SGAInlineRefTarget{sga3:xxii:def:1.5:item:b}(b)} If
\(X\in\Gamma(S,\mathfrak g^\alpha)^\times\), then the inner
automorphism \(w_\alpha(X)\) of \(T\), defined by
\[
  w_\alpha(X)
  =
  \exp_\alpha(X)\exp_{-\alpha}(-X^{-1})\exp_\alpha(X)
\]
(cf. \SGARefLink{sga3:xx:3.1}{Exposé XX, 3.1}), coincides with \(s_\alpha\) (loc. cit.).
\SGARefLink{sga3:xxii:prop:1.4}{Proposition 1.4} therefore gives the following.

\SGARefTarget{sga3:xxii:cor:1.6}\paragraph{\SGARefLink{sga3:xxii:cor:1.6}{Corollary 1.6}.}
\textit{Let \(S\) be a scheme, \(G\) a reductive \(S\)-group, \(T\)
a maximal torus of \(G\), and \(\alpha,\beta\) two roots of \(G\)
relative to \(T\). Then}
\[
  s_\alpha(\beta)=\beta-(\alpha^*,\beta)\alpha
\]
\textit{is a root of \(G\) relative to \(T\), and its corresponding
coroot is}
\[
  s_\alpha(\beta)^*
  =
  s_\alpha(\beta^*)
  =
  \beta^*-(\beta^*,\alpha)\alpha^*.
\]

\SGARefTarget{sga3:xxii:cor:1.7}\paragraph{\SGARefLink{sga3:xxii:cor:1.7}{Corollary 1.7}.}
\textit{Under the preceding conditions,
\(\alpha^*=\beta^*\) implies \(\alpha=\beta\).}

Indeed, if \(\alpha^*=\beta^*\), one has (cf. \SGARefLink{sga3:xx:1.4}{Exposé XXI, 1.4})
\[
  s_\beta(\alpha)=\alpha-2\beta,
  \qquad
  s_\alpha(\beta)=\beta-2\alpha,
\]
and hence immediately
\[
  (s_\beta s_\alpha)^n(\alpha)
  =
  \alpha+2n(\beta-\alpha).
\]
If \(\beta\ne\alpha\), there exists \(s\in S\) such that
\(\alpha_s\ne\beta_s\). The preceding formula would then give
infinitely many distinct roots of \(G_{\bar s}\) relative to
\(T_{\bar s}\), which is impossible.

\SGARefTarget{sga3:xxii:def:1.8.0}\paragraph{\SGARefLink{sga3:xxii:def:1.8.0}{Definitions 1.8.0}.}
If \(u:\mathbb G_{m,S}\to T\) is a homomorphism of groups, we say
that \(u\) is a \emph{coroot} of \(G\) relative to \(T\) if there
exists a root \(\alpha\) of \(G\) relative to \(T\) such that
\(\alpha^*=u\). Consider the functor \(\mathcal R^*\) of coroots of
\(G\) relative to \(T\), defined by
\[
  \mathcal R^*(S')
  =
  \{\text{coroots of \(G_{S'}\) relative to \(T_{S'}\)}\}.
\]
If \(\mathcal R\) is the functor of roots of \(G\) relative to \(T\)
(\SGARefLink{sga3:xix:proposition:3:8}{Exposé XIX, 3.8}), there is a canonical morphism
\(\mathcal R\to\mathcal R^*\). By \SGARefLink{sga3:xxii:cor:1.7}{1.7} and \SGARefLink{sga3:xix:proposition:3-8}{Exposé XIX, 3.8}, one has:

\SGARefTarget{sga3:xxii:cor:1.8}\paragraph{\SGARefLink{sga3:xxii:cor:1.8}{Corollary 1.8}.}
\textit{The canonical morphism
\(\mathcal R\to\mathcal R^*\) is an isomorphism. In particular,
\(\mathcal R^*\) is represented by a twisted finite constant
\(S\)-scheme that is an open and closed subscheme of}
\[
  \underline{\operatorname{Hom}}_{S\text{-gr}}
    (\mathbb G_{m,S},T).
\]

This leads to the following definition.

\SGARefTarget{sga3:xxii:def:1.9}\paragraph{\SGARefLink{sga3:xxii:def:1.9}{Definition 1.9}.}
\textit{Let \(S\) be a scheme and \(T\) an \(S\)-torus. A
\emph{twisted root datum} in \(T\) consists of}
\begin{enumerate}
\item[(i)] \textit{a finite subscheme \(\mathcal R\) of}
  \[
    \underline{\operatorname{Hom}}_{S\text{-gr}}
      (T,\mathbb G_{m,S});
  \]
\item[(ii)] \textit{a finite subscheme \(\mathcal R^*\) of}
  \[
    \underline{\operatorname{Hom}}_{S\text{-gr}}
      (\mathbb G_{m,S},T);
  \]
\item[\SGAInlineRefTarget{sga3:xxii:def:1.9:item:iii}(iii)] \textit{an isomorphism
  \(\mathcal R\xrightarrow{\sim}\mathcal R^*\), denoted
  \(\alpha\mapsto\alpha^*\),}
\end{enumerate}
\textit{satisfying the following conditions:}
\begin{description}
\item[\textup{(DR 1)}]
  \textit{For every \(S'\to S\) and every
  \(\alpha\in\mathcal R(S')\), one has
  \(\alpha\circ\alpha^*=2\).}
\item[\textup{(DR 2)}]
  \textit{For every \(S'\to S\) and all
  \(\alpha,\beta\in\mathcal R(S')\), one has}
  \[
    \alpha-(\beta^*,\alpha)\beta\in\mathcal R(S'),
    \qquad
    \alpha^*-(\alpha^*,\beta)\beta^*\in\mathcal R^*(S').
  \]
\end{description}
\textit{Moreover, the twisted root datum is said to be
\emph{reduced} if, for every nonempty \(S'\to S\) and every
\(\alpha\in\mathcal R(S')\), one has
\(2\alpha\notin\mathcal R(S')\).}

\SGARefTarget{sga3:xxii:prop:1.10}\paragraph{\SGARefLink{sga3:xxii:prop:1.10}{Proposition 1.10}.}
\textit{Let \(S\) be a scheme, \(G\) a reductive \(S\)-group, \(T\)
a maximal torus of \(G\), and let \(\mathcal R\), respectively
\(\mathcal R^*\), be the scheme of roots, respectively coroots, of
\(G\) relative to \(T\). Then \((\mathcal R,\mathcal R^*)\) is a
reduced twisted root datum in \(T\).}

It remains only to verify that this twisted root datum is reduced.
This was done in \SGARefLink{sga3:xix:paragraph:3:10}{Exposé XIX, 3.10}.

\SGARefTarget{sga3:xxii:par:1.11}\paragraph{\SGARefLink{sga3:xxii:par:1.11}{1.11}.}
Let \(T=\mathrm D_S(M)\) be a \emph{trivialized} torus. If \(M^*\)
denotes the abelian group dual to \(M\), there are canonical
isomorphisms (cf. \SGARefLink{sga3:viii:corollary:1:5}{Exposé VIII, 1.5})
\[
  \underline{\operatorname{Hom}}_{S\text{-gr}}
    (T,\mathbb G_{m,S})
  \xrightarrow{\sim} M_S,
  \qquad
  \underline{\operatorname{Hom}}_{S\text{-gr}}
    (\mathbb G_{m,S},T)
  \xrightarrow{\sim} M^*_S,
\]
and hence group isomorphisms
\[
  \operatorname{Hom}_{S\text{-gr}}(T,\mathbb G_{m,S})
  \xrightarrow{\sim}
  \operatorname{Hom}_{\mathrm{loc.const.}}(S,M),
\]
\[
  \operatorname{Hom}_{S\text{-gr}}(\mathbb G_{m,S},T)
  \xrightarrow{\sim}
  \operatorname{Hom}_{\mathrm{loc.const.}}(S,M^*).
\]
A character of \(T\), respectively a group homomorphism
\(\mathbb G_{m,S}\to T\), is said to be \emph{constant} relative to
the given trivialization if the preceding isomorphism carries it to
a constant map from \(S\) to \(M\), respectively to \(M^*\).

\SGARefTarget{sga3:xxii:par:1.12}\paragraph{\SGARefLink{sga3:xxii:par:1.12}{1.12}.}
With the same notation, let \((M,M^*,R,R^*)\) be a root datum
(Exposé XXI). Then \((R_S,R^*_S)\) is a twisted root datum in \(T\).
Conversely, if \((\mathcal R,\mathcal R^*)\) is a twisted root datum
in a torus \(T\), a \emph{splitting} of this root datum means the
data of an ordinary root datum \((M,M^*,R,R^*)\) and an isomorphism
\[
  T\simeq\mathrm D_S(M)
\]
that carries \((\mathcal R,\mathcal R^*)\) to \((R_S,R^*_S)\).

\SGARefTarget{sga3:xxii:def:1.13}\paragraph{\SGARefLink{sga3:xxii:def:1.13}{Definition 1.13}.}
\textit{Let \(S\) be a scheme, \(G\) a reductive \(S\)-group, and
\(T\) a maximal torus of \(G\). A \emph{splitting of \(G\) relative
to \(T\)} consists of}
\begin{enumerate}
\item[\SGAInlineRefTarget{sga3:xxii:def:1.13:item:i}(i)] \textit{an abelian group \(M\) and an isomorphism
  \(T\simeq\mathrm D_S(M)\);}
\item[\SGAInlineRefTarget{sga3:xxii:def:1.13:item:ii}(ii)] \textit{a system of roots \(R\) of \(G\) relative to
  \(T\) \textup{(\SGARefLink{sga3:xix:definition:3-6}{Exposé XIX, 3.6})},}
\end{enumerate}
\textit{satisfying the following two conditions:}
\begin{description}
\item[\textup{(D 1)}]
  \textit{\(S\) is nonempty, and the roots \(\alpha\in R\)
  \textup{(respectively, the corresponding coroots)} identify with
  constant functions from \(S\) to \(M\) \textup{(respectively, to
  \(M^*\))}.}
\item[\textup{(D 2)}]
  \textit{The \(\mathfrak g^\alpha\), \(\alpha\in R\), are free
  \(\mathcal O_S\)-modules.}
\end{description}

We say that \(G\) is \emph{splittable relative to \(T\)} if a
splitting of \(G\) relative to \(T\) exists. A \emph{splitting of
\(G\)} consists of a maximal torus \(T\) of \(G\) together with a
splitting of \(G\) relative to \(T\). We say that \(G\) is
\emph{splittable} if it admits a splitting. A \emph{split
\(S\)-group} is a reductive \(S\)-group equipped with a splitting;
we denote it by a symbol of the form \((G,T,M,R)\), or simply by
\(G\) when there is no risk of confusion.

Condition \textup{(D 1)} canonically identifies \(R\), respectively
\(R^*\), with a subset of \(M\), respectively of \(M^*\).

\SGARefTarget{sga3:xxii:prop:1.14}\paragraph{\SGARefLink{sga3:xxii:prop:1.14}{Proposition 1.14}.}
\textit{Let \(S\) be a nonempty scheme and let \((G,T,M,R)\) be a
split \(S\)-group. Then}
\[
  \mathcal R(G,T,M,R)=(M,M^*,R,R^*)
\]
\textit{is a reduced root datum \textup{(\SGARefLink{sga3:xxi:subsection:1-1}{Exposé XXI, 1.1} and
\SGARefLink{sga3:XXI:2.1.3}{2.1.3})}; it is a splitting of the twisted root datum of \SGARefLink{sga3:xxii:prop:1.10}{1.10}.}

This follows trivially from \SGARefLink{sga3:xxii:prop:1.10}{1.10} and \SGARefLink{sga3:xix:lemma:3:7}{Exposé XIX, 3.7}.

For brevity we shall sometimes write
\(\mathcal R(G,T,M,R)=\mathcal R(G)\). We shall systematically use
the notation \(V\), \(V(\mathcal R)\), \(W,\ldots\) of Exposé XXI.

\SGARefTarget{sga3:xxii:rem:1.15}\paragraph{\SGARefLink{sga3:xxii:rem:1.15}{Remark 1.15}.}
\textup{(a)} If \(S\) is nonempty and connected, respectively if
\(\operatorname{Pic}(S)=0\), condition \textup{(D 1)}, respectively
condition \textup{(D 2)}, is automatically satisfied.

\SGAInlineRefTarget{sga3:xxii:rem:1.15:item:b}\textup{\SGARefLink{sga3:xxii:def:1.5:item:b}{(b)}} If \((G,T,M,R)\) is a split \(S\)-group, then for every
nonempty \(S'\to S\), \((G_{S'},T_{S'},M,R)\) is a split
\(S'\)-group and
\[
  \mathcal R(G,T,M,R)
  =
  \mathcal R(G_{S'},T_{S'},M,R).
\]

\SGARefTarget{sga3:xxii:par:1.16}\paragraph{\SGARefLink{sga3:xxii:par:1.16}{1.16}.}
Let \(T=\mathrm D_S(M)\) be a trivialized torus. The Lie algebra
\(\mathfrak t\) of \(T\) identifies canonically \textup{(\SGARefLink{sga3:ii:corollary:5-1-1}{Exposé II,
5.1.1})} with
\[
  \mathfrak t\simeq M^*\otimes\mathcal O_S.
\]
For every homomorphism of groups \(u:T\to\mathbb G_{m,S}\),
\(\Lie(u)\) is a linear form
\[
  \Lie(u):
  \mathfrak t\longrightarrow
  \mathcal O_S=\Lie(\mathbb G_{m,S}/S).
\]
In particular, if \(u\) is defined by an element \(\alpha\in M\),
then \(\Lie(u)\) is the linear form \(\bar\alpha\) on
\(M^*\otimes\mathcal O_S\) defined by
\[
  \bar\alpha(m\otimes x)=(m,\alpha)x.
\]

Symmetrically, for every homomorphism of groups
\(h:\mathbb G_{m,S}\to T\), \(\Lie(h)\) is an
\(\mathcal O_S\)-linear morphism
\[
  \mathcal O_S=\Lie(\mathbb G_{m,S}/S)
  \longrightarrow\mathfrak t,
\]
canonically determined by the section
\[
  H=\Lie(h)(1)\in\Gamma(S,\mathfrak t).
\]
In particular, if \(h\) is defined by an element \(m\in M^*\), then
\[
  H=\Lie(h)(1)=m\otimes 1.
\]
Comparing the two definitions, we find in particular
\[
  \bar\alpha(H)=(h,\alpha)\cdot 1
  \in\Gamma(S,\mathcal O_S).
\]

\SGARefTarget{sga3:xxii:par:1.17}\paragraph{\SGARefLink{sga3:xxii:par:1.17}{1.17}.}
These definitions apply in particular when \(T\) is the maximal
torus of a split group. Every root \(\alpha\in R\) defines an
\emph{infinitesimal root}
\[
  \bar\alpha\in
  \operatorname{Hom}_{\mathcal O_S}(\mathfrak t,\mathcal O_S),
  \qquad
  \bar\alpha(m\otimes x)=(m,\alpha)x.
\]
Every coroot \(\alpha^*\), \(\alpha\in R\), defines an
\emph{infinitesimal coroot}
\[
  H_\alpha\in\Gamma(S,\mathfrak t),
  \qquad
  H_\alpha=\alpha^*\otimes 1.
\]
For \(\alpha,\beta\in R\), we have
\[
  \bar\alpha(H_\beta)=(\beta^*,\alpha)\cdot 1,
\]
and in particular
\[
  \bar\alpha(H_\alpha)=2.
\]
Consequently, if \(2\) is invertible on \(S\), then
\(\bar\alpha\) and \(H_\alpha\) are nonzero on every fiber.

\SGARefTarget{sga3:xxii:section:2}\section*{2. Existence of a
Splitting. Type of a Reductive Group}
\addcontentsline{toc}{section}{2. Existence of a Splitting. Type of
a Reductive Group}

\SGARefTarget{sga3:xxii:prop:2.1}\paragraph{\SGARefLink{sga3:xxii:prop:2.1}{Proposition 2.1}.}
\textit{Let \(S\) be a scheme, \(G\) a reductive \(S\)-group, and
\(T\) a maximal torus of \(G\). Suppose that \(T\) is split. Then
\(G\) is locally splittable relative to \(T\): for every
\(s_0\in S\), there exists an open neighborhood \(U\) of \(s_0\)
such that the \(U\)-group \(G_U\) is splittable relative to
\(T_U\).}

Indeed, write \(T\simeq\mathrm D_S(M)\) and
\[
  \mathfrak g=\coprod_{m\in M}\mathfrak g^m.
\]
Let
\[
  R=\{m\in M\mid m\ne0,\ \mathfrak g^m(s_0)\ne0\}.
\]
After restricting \(S\) and replacing it by an open neighborhood
of \(s_0\), we may suppose that the \(\mathfrak g^\alpha\),
\(\alpha\in R\), are free and that the \(\mathfrak g^m\), for
\(m\ne0\) and \(m\notin R\), vanish. We then have
\[
  \mathfrak g
  =
  \mathfrak t\oplus
  \coprod_{\alpha\in R}\mathfrak g^\alpha,
\]
where the \(\mathfrak g^\alpha\) are free of rank \(1\). It follows
that \(R\) is a system of roots of \(G\) relative to \(T\)
\textup{(\SGARefLink{sga3:xix:definition:3-6}{Exposé XIX, 3.6})}. The coroots \(\alpha^*\) corresponding
to the \(\alpha\in R\) then identify with locally constant
functions on \(S\) with values in \(M^*\). Shrinking \(S\) once
more, we may suppose that they are constant, and the proof is
complete.

Observe that the proof immediately gives:

\SGARefTarget{sga3:xxii:prop:2.2}\paragraph{\SGARefLink{sga3:xxii:prop:2.2}{Proposition 2.2}.}
\textit{Let \(S\) be a nonempty connected scheme such that
\(\operatorname{Pic}(S)=0\), for example
\(\operatorname{Spec}(\mathbb Z)\) or a local scheme
\textup{(in particular, the spectrum of a field)}. If \(G\) is a
reductive \(S\)-group possessing a split maximal torus \(T\), then
\(G\) is splittable relative to \(T\).}

We deduce immediately from \SGARefLink{sga3:xxii:prop:2.1}{2.1} and from the fact that a reductive
group locally possesses maximal tori for the étale topology
\textup{(\SGARefLink{sga3:xix:theorem:2:5}{Exposé XIX, 2.5})}:

\SGARefTarget{sga3:xxii:cor:2.3}\paragraph{\SGARefLink{sga3:xxii:cor:2.3}{Corollary 2.3}.}
\textit{Let \(S\) be a scheme and \(G\) a reductive \(S\)-group
\textup{(respectively, let \(T\) also be a maximal torus of \(G\))}.
Then \(G\) is locally splittable
\textup{(respectively, locally splittable relative to \(T\))}
for the étale topology on \(S\).}

\SGARefTarget{sga3:xxii:cor:2.4}\paragraph{\SGARefLink{sga3:xxii:cor:2.4}{Corollary 2.4}.}
\textit{Let \(k\) be a field and \(G\) a reductive \(k\)-group.
There exists a finite separable extension \(K/k\) such that \(G_K\)
is splittable.}

\SGARefTarget{sga3:xxii:rem:2.5}\paragraph{\SGARefLink{sga3:xxii:rem:2.5}{Remark 2.5}.}
Using \SGARefLink{sga3:xxii:prop:2.1}{2.1} and \SGARefLink{sga3:xix:remark:2:9}{Exposé XIX, Remark 2.9}, one immediately proves the
following result: let \(G=(G,T,M,R)\) be a split \(S\)-group. There
exists an open covering \((U_i)\) of \(S\) such that each split group
\(G_{U_i}\) is obtained by base change from a split group over a
noetherian ring \textup{(and in fact over a finitely generated
\(\mathbb Z\)-algebra)}. We shall moreover prove that every split
group over \(S\) is already obtained from a split
\(\mathbb Z\)-group \textup{(Exposé XXV)}.

\SGARefTarget{sga3:xxii:par:2.6}\paragraph{\SGARefLink{sga3:xxii:par:2.6}{2.6}.}
Let \(k\) be an algebraically closed field and \(G\) a reductive
\(k\)-group. We know \textup{(for example by \SGARefLink{sga3:xxii:cor:2.4}{2.4})} that splittings of
\(G\) exist. Let \((G,T,M,R)\) and \((G,T',M',R')\) be two splittings
of \(G\). Then the root data
\[
  \mathcal R(G,T,M,R)
  \qquad\text{and}\qquad
  \mathcal R(G,T',M',R')
\]
are isomorphic.

Indeed, we first see that it is enough to reduce to the case
\(T=T'\): there exists \(g\in G(k)\) such that
\(T'=\operatorname{int}(g)T\), and one readily verifies that
transporting a splitting by an automorphism of \(G\) produces a root
datum isomorphic to the original one. But since
\(S=\operatorname{Spec}(k)\) is connected, the isomorphism
\[
  \mathrm D_k(M)\xrightarrow{\sim}T
  \xrightarrow{\sim}\mathrm D_k(M')
\]
comes from a unique isomorphism \(M\simeq M'\). For the same reason,
there exists at most one root system of \(G\) relative to \(T\).

\SGARefTarget{sga3:xxii:def:2.6.1}\paragraph{\SGARefLink{sga3:xxii:def:2.6.1}{Definition 2.6.1}.}
If \(G\) is a reductive \(k\)-group, where \(k\) is algebraically
closed, the \emph{type of \(G\)} is the isomorphism class of the root
datum defined by any splitting of \(G\). If \(G\) is a torus of type
\(M\) in the sense of \SGARefLink{sga3:ix:definition:1:4}{Exposé IX, 1.4}, then the type of \(G\) as a
reductive group is given by the trivial root datum
\[
  (M,M^*,\varnothing,\varnothing).
\]

By \SGARefLink{sga3:xxii:rem:1.15:item:b}{1.15(b)}, the type is invariant under an
\textup{(algebraically closed)} extension of the base field.

\SGARefTarget{sga3:xxii:def:2.7}\paragraph{\SGARefLink{sga3:xxii:def:2.7}{Definition 2.7}.}
If \(G\) is a reductive \(S\)-group and \(s\in S\), the \emph{type of
\(G\) at \(s\)} is the type of the reductive
\(\bar s\)-group \(G_{\bar s}\).

For every \(S'\to S\) and every \(s'\in S'\) mapping to \(s\in S\),
the type of \(G_{S'}\) at \(s'\) is equal to the type of \(G\) at
\(s\).

If \(G\) is splittable and \((G,T,M,R)\) is a splitting of \(G\),
then the type of \(G\) at \(s\) is the isomorphism class of
\(\mathcal R(G,T,M,R)\), by \SGARefLink{sga3:xxii:rem:1.15:item:b}{1.15(b)}. The following
then follows immediately from \SGARefLink{sga3:xxii:cor:2.3}{2.3}:

\SGARefTarget{sga3:xxii:prop:2.8}\paragraph{\SGARefLink{sga3:xxii:prop:2.8}{Proposition 2.8}.}
\textit{Let \(G\) be a reductive \(S\)-group, where
\(S\ne\varnothing\). The function}
\[
  s\longmapsto\text{\textit{the type of \(G\) at \(s\)}}
\]
\textit{is locally constant on \(S\). In particular, there exists a
partition of \(S\) into nonempty open subschemes on each of which
\(G\) has constant type. More precisely, let \(E\) be the set of
types of the fibers of \(G\), and, for every \(t\in E\), let \(S_t\)
be the set of points \(s\in S\) where \(G\) has type \(t\). Then
\((S_t)_{t\in E}\) is a partition of \(S\), and each \(S_t\) is open
and closed \textup{(and nonempty)}.}

\SGARefTarget{sga3:xxii:section:3}\section*{3. The Weyl Group}
\addcontentsline{toc}{section}{3. The Weyl Group}

\SGARefTarget{sga3:xxii:par:3.1}\paragraph{\SGARefLink{sga3:xxii:par:3.1}{3.1}.}
Let \(S\) be a scheme, \(G\) a reductive \(S\)-group, and \(T\) a
maximal torus of \(G\). Then
\[
  W_G(T)=\underline{\operatorname{Norm}}_G(T)/T
\]
is a finite étale \(S\)-group \textup{(\SGARefLink{sga3:xix:theorem:2-5}{Exposé XIX, 2.5})}. The
morphism \(n\mapsto\operatorname{int}(n)\) induces, upon passage to
the quotient, a canonical monomorphism \textup{(which is moreover an
open immersion)}
\[
  W_G(T)\longrightarrow
  \underline{\operatorname{Aut}}_{S\text{-gr.}}(T).
\]

\SGARefTarget{sga3:xxii:par:3.2}\paragraph{\SGARefLink{sga3:xxii:par:3.2}{3.2}.}
Suppose now that \(G\) is splittable relative to \(T\). Choose a
splitting, say \((G,T,M,R)\). We then have a canonical isomorphism
\textup{(\SGARefLink{sga3:viii:corollary:1:5}{Exposé VIII, 1.5})}
\[
  \underline{\operatorname{Aut}}_{S\text{-gr.}}(T)
  \simeq
  \bigl(\operatorname{Aut}_{\mathrm{gr.}}(M)\bigr)_S.
\]
In particular, if \(W\) is the Weyl group of the root datum
\(\mathcal R(G)\) \textup{(\SGARefLink{sga3:XXI:1.1.8}{Exposé XXI, 1.1.8})}, we have a
monomorphism
\[
  W_S\longrightarrow
  \underline{\operatorname{Aut}}_{S\text{-gr.}}(T).
\]

\SGARefTarget{sga3:xxii:par:3.3}\paragraph{\SGARefLink{sga3:xxii:par:3.3}{3.3}.}
For every root \(\alpha\in R\), the reflection \(s_\alpha\in W\)
acts on \(M\) by
\[
  s_\alpha(x)=x-(\alpha^*,x)\alpha,
\]
and hence acts on \(T\), by the preceding morphism, according to
\[
  s_\alpha(t)=t\cdot\alpha^*\bigl(\alpha(t)\bigr)^{-1}.
\]
On the other hand, since \(\mathfrak g^\alpha\) is assumed to be
free, there exists
\[
  X\in\Gamma(S,\mathfrak g^\alpha)^\times.
\]
Consider
\[
  w_\alpha(X)\in
  \underline{\operatorname{Norm}}_G(T)(S)
  \qquad\textup{(Exposé XX, 3.1)}.
\]
We have \textup{(loc.\ cit.)}
\[
  \operatorname{int}\bigl(w_\alpha(X)\bigr)(t)=s_\alpha(t).
\]
Since \(W\) is generated by the \(s_\alpha\), \(\alpha\in R\), it
follows from the preceding remarks that, if \(W\) and
\(\underline{\operatorname{Norm}}_G(T)(S)/T(S)\) are regarded as
groups of automorphisms of \(T\), then
\[
  W
  \subset
  \underline{\operatorname{Norm}}_G(T)(S)/T(S)
  \subset
  W_G(T)(S).
\]

By definition of the constant group \(W_S\) associated with \(W\)
\textup{(cf.\ \SGARefLink{sga3:i:subsection:1:8}{Exposé I, 1.8})}, we therefore have a commutative
diagram:

\begin{center}
  \begin{tikzcd}[column sep=5.0em,row sep=3.0em]
  W_S
    \arrow[rr]
    \arrow[dr]
  &&
  W_G(T)
    \arrow[dl]
  \\
  &
  \underline{\operatorname{Aut}}_{S\text{-gr.}}(T).
  &
\end{tikzcd}

\end{center}

\SGARefTarget{sga3:xxii:prop:3.4}\paragraph{\SGARefLink{sga3:xxii:prop:3.4}{Proposition 3.4}.}
\textit{Let \(S\) be a scheme, let \((G,T,M,R)\) be a split
\(S\)-group, and let \(W\) be the Weyl group of the root datum
\(\mathcal R(G)\). Then the canonical monomorphism}
\[
  W_S\longrightarrow
  W_G(T)=\underline{\operatorname{Norm}}_G(T)/T
\]
\textit{is an isomorphism.}

Indeed, these are étale groups over \(S\); it therefore suffices to
verify that, for every \(s\in S\), the morphism
\[
  W_S(\bar s)\longrightarrow W_G(T)(\bar s)
\]
is an isomorphism.
The latter assertion follows, for example, from \emph{Bible},
\S11.3, Theorem 2.

\SGARefTarget{sga3:xxii:rem:3.5}\paragraph{\SGARefLink{sga3:xxii:rem:3.5}{Remark 3.5}.}
Using \SGARefLink{sga3:xxii:cor:2.3}{2.3}, the preceding proposition gives a new proof of the fact
that the Weyl group of a maximal torus of a reductive \(S\)-group
\(G\) is finite over \(S\) \textup{(Exposé XIX, 2.5(ii))}.

\SGARefTarget{sga3:xxii:par:3.6}\paragraph{\SGARefLink{sga3:xxii:par:3.6}{3.6}.}
Under the hypotheses of \SGARefLink{sga3:xxii:par:3.1}{3.1}, for every \(w\in W_G(T)(S)\), let
\(N_w\)
denote the fiber product in the following diagram:

\begin{center}
  \begin{tikzcd}[column sep=5.2em,row sep=3.2em]
  N_w
    \arrow[r]
    \arrow[d]
  &
  \underline{\operatorname{Norm}}_G(T)
    \arrow[d]
  \\
  S
    \arrow[r,"w"]
  &
  W_G(T).
\end{tikzcd}

\end{center}

It is an open and closed subscheme of
\(\underline{\operatorname{Norm}}_G(T)\), and is a principal
homogeneous bundle under \(T\) on the left \textup{(respectively,
on the right)} for the law
\[
  (t,q)\longmapsto tq
  \qquad
  \textup{(respectively, \((q,t)\longmapsto qt\))}.
\]
If \(n\in N_w(S)\), then
\[
  N_{ww'}=n\mathbin{\cdot}N_{w'},
  \qquad
  N_{w'w}=N_{w'}\mathbin{\cdot}n.
\]

\SGARefTarget{sga3:xxii:par:3.7}\paragraph{\SGARefLink{sga3:xxii:par:3.7}{3.7}.}
In particular, if \(\alpha\) is a root of \(G\) with respect to
\(T\), then \(N_{s_\alpha}\) is none other than what was denoted by
\(N^\times\) in \SGARefLink{sga3:xx:3.0}{Exposé XX, 3.0}. If \(\mathfrak g^\alpha\) is free
over \(S\), we therefore have
\[
  N_{s_\alpha}(S)\ne\varnothing.
\]
By \SGARefLink{sga3:xxii:prop:3.4}{3.4} and condition \textup{(D\,2)} of the splitting, we deduce:

\SGARefTarget{sga3:xxii:cor:3.8}\paragraph{\SGARefLink{sga3:xxii:cor:3.8}{Corollary 3.8}.}
\textit{Under the hypotheses of \SGARefLink{sga3:xxii:prop:3.4}{3.4}, the morphism}
\[
  \underline{\operatorname{Norm}}_G(T)(S)
  \longrightarrow
  W_G(T)(S)
  =
  \operatorname{Hom}_{\mathrm{loc.const.}}(S,W)
\]
\textit{is surjective. In particular, for every \(w\in W\), there
exists \(n_w\in\underline{\operatorname{Norm}}_G(T)(S)\) such that}
\[
  \operatorname{int}(n_w)|_T=w.
\]

\SGARefTarget{sga3:xxii:section:4}
\section*{4. Homomorphisms of Split Groups}
\addcontentsline{toc}{section}{4. Homomorphisms of Split Groups}

\SGARefTarget{sga3:xxii:subsection:4.1}
\subsection*{4.1. The ``Big Cell''}
\addcontentsline{toc}{subsection}{\SGARefLink{sga3:xxii:subsection:4.1}{4.1}. The ``Big Cell''}

\SGARefTarget{sga3:xxii:par:4.1.1}\paragraph{\SGARefLink{sga3:xxii:par:4.1.1}{4.1.1}.}
Let \((G,T,M,R)\) be a split reductive \(S\)-group. Choose a system
of positive roots \(R_+\) of the root datum \(\mathcal R(G)\)
\textup{(\SGARefLink{sga3:XXI:3.2.1}{Exposé XXI, 3.2.1})}. Put \(R_-=-R_+\).

Choose a total ordering on \(R_+\) \textup{(respectively, \(R_-\))}
and consider the morphism induced by multiplication in \(G\):
\[
  u:
  \left(\prod_{\alpha\in R_-}U_\alpha\right)
  \times_S T
  \times_S
  \left(\prod_{\alpha\in R_+}U_\alpha\right)
  \longrightarrow G.
\]
This is an open immersion. Indeed, since both sides are flat and of
finite presentation over \(S\), it suffices to verify this on every
geometric fiber \textup{(SGA 1, I, 5.7 and VIII, 5.4)}. We are thus
reduced to the case in which \(S\) is the spectrum of an
algebraically closed field. But by \emph{Bible}, \S13.4,
Corollary 2 to Theorem 3, \(u\) is radicial and dominant. Since the
tangent map of \(u\) at the origin is an isomorphism
\textup{(definition of a root system)}, \(u\) is birational. As
\(G\) is normal, Zariski's ``Main Theorem''
\textup{(EGA III\textsubscript{1}, 4.4.9)} applies, and \(u\) is an
open immersion.

Let us show that the image \(\Omega\) of this open immersion is
independent of the ordering chosen on \(R_+\)
\textup{(respectively, \(R_-\))}. Since this amounts to comparing
open subsets of \(G\), we are reduced to proving that they have the
same geometric points, and may therefore again suppose that \(S\)
is the spectrum of an algebraically closed field. The assertion is
then precisely \emph{Bible}, \S13, Proposition 1(c) and
Theorem 1(a).

We have therefore proved:

\SGARefTarget{sga3:xxii:prop:4.1.2}\paragraph{\SGARefLink{sga3:xxii:prop:4.1.2}{Proposition 4.1.2}.}
\textit{Let \((G,T,M,R)\) be a split \(S\)-group, and let \(R_+\)
be a system of positive roots of \(R\). There exists an open subset
\(\Omega_{R_+}\) of \(G\) such that, for every total ordering on
\(R_+\) \textup{(respectively, \(R_-\))}, the morphism induced by
multiplication in \(G\)}
\[
  \left(\prod_{\alpha\in R_-}U_\alpha\right)
  \times_S T
  \times_S
  \left(\prod_{\alpha\in R_+}U_\alpha\right)
  \longrightarrow G
\]
\textit{is an open immersion with image \(\Omega_{R_+}\).}

\SGARefTarget{sga3:xxii:rem:4.1.3}\paragraph{\SGARefLink{sga3:xxii:rem:4.1.3}{Remark 4.1.3}.}
\SGARefLink{sga3:xxii:prop:4.1.2}{Proposition 4.1.2} may be restated as follows: for every
\(\alpha\in R\), choose an isomorphism of vector groups
\[
  p_\alpha:\mathbb G_{a,S}\xrightarrow{\;\sim\;}U_\alpha
  \qquad\textup{(cf.\ 1.19)}.
\]
Then the morphism \textup{(where
\(N=\operatorname{Card}(R_+)=\operatorname{Card}(R_-)\))}
\[
  \mathbb G_{a,S}^{N}
  \times_S T
  \times_S
  \mathbb G_{a,S}^{N}
  \longrightarrow G
\]
defined on points by
\[
  \bigl((x_\alpha)_{\alpha\in R_-},t,
        (x_\alpha)_{\alpha\in R_+}\bigr)
  \longmapsto
  \prod_{\alpha\in R_-}p_\alpha(x_\alpha)\mathbin{\cdot}t
  \mathbin{\cdot}\prod_{\alpha\in R_+}p_\alpha(x_\alpha)
\]
is an open immersion whose image depends only on \(R_+\), and not
on the choice of the \(p_\alpha\) or of the orderings on \(R_+\)
and \(R_-\).

\SGARefTarget{sga3:xxii:notation:4.1.4}\paragraph{\SGARefLink{sga3:xxii:notation:4.1.4}{Notation 4.1.4}.}
We write
\[
  \Omega_{R_+}
  =
  \prod_{\alpha\in R_-}U_\alpha
  \mathbin{\cdot}T\mathbin{\cdot}
  \prod_{\alpha\in R_+}U_\alpha.
\]

\SGARefTarget{sga3:xxii:prop:4.1.5}\paragraph{\SGARefLink{sga3:xxii:prop:4.1.5}{Proposition 4.1.5}.}
\textit{The scheme \(\Omega_{R_+}\) is of finite presentation over
\(S\) \textup{(hence quasi-compact in \(G\))} and is universally
schematically dense in \(G\) relative to \(S\)
\textup{(cf.\ \SGARefLink{sga3:xviii:section:1}{Exposé XVIII, 1})}.}

The first assertion is trivial. Thus,\footnote{Editor's note: We
have expanded the original in what follows.}
\(\Omega_{R_+}\) is flat and of finite presentation over \(S\), and
contains the identity section; it therefore meets every fiber of
\(G\) in a nonempty, hence dense, open subset. The second assertion
thus follows from Exposé XVIII, 1.3.

\SGARefTarget{sga3:xxii:cor:4.1.6}\paragraph{\SGARefLink{sga3:xxii:cor:4.1.6}{Corollary 4.1.6}.}
\textit{Let \((G,T,M,R)\) be a split reductive \(S\)-group. Then}
\[
  \underline{\operatorname{Centr}}(G)
  =
  \bigcap_{\alpha\in R}\operatorname{Ker}(\alpha).
\]
\textit{Consequently,
\(\underline{\operatorname{Centr}}(G)\) is representable by a
closed diagonalizable subgroup of \(G\).}
\footnote{Editor's note: We have added the preceding sentence.}

The second assertion follows immediately from the first. To prove
the first, one may invoke \SGARefLink{sga3:xii:4.8}{Exposé XII, 4.8} and \SGARefLink{sga3:xii:4.11}{4.11}; alternatively,
one may argue directly as follows.\footnote{Editor's note: We have
expanded the original in what follows.}
Let \(S'\to S\). If \(t\in T(S')\) and \(\alpha(t)=1\) for every
\(\alpha\in R\), then \(\operatorname{int}(t)\) induces the identity
on \(T_{S'}\) and on every \((U_\alpha)_{S'}\), \(\alpha\in R\),
and therefore also on \((\Omega_{R_+})_{S'}\). By schematic density
it consequently induces the identity on \(G_{S'}\), whence
\[
  t\in\underline{\operatorname{Centr}}(G)(S').
\]
Conversely, since
\[
  \underline{\operatorname{Centr}}_{G_{S'}}(T_{S'})
  =
  T_{S'}
  \qquad\textup{(cf.\ Exposé XIX, 2.8)},
\]
if \(g\in G(S')\) centralizes \(T_{S'}\) and the
\((U_\alpha)_{S'}\), then \(g\) is a section of \(T_{S'}\) on
which every \(\alpha\in R\) vanishes.

\SGARefTarget{sga3:xxii:cor:4.1.7}\paragraph{\SGARefLink{sga3:xxii:cor:4.1.7}{Corollary 4.1.7}.}
\textit{Let \(S\) be a scheme and \(G\) a reductive \(S\)-group.
Then the center of \(G\) is representable by a closed subgroup of
\(G\) of multiplicative type. It is also ``the intersection of the
maximal tori of \(G\)'' in the following sense: for every
\(S'\to S\), \(\underline{\operatorname{Centr}}(G)(S')\) is the set
of \(g\in G(S')\) such that, for every \(S''\to S'\), the inverse
image of \(g\) in \(G(S'')\) belongs to \(T(S'')\) for every maximal
torus \(T\) of \(G_{S''}\).}

In view of \SGARefLink{sga3:xxii:cor:2.3}{2.3}, the first assertion follows from \SGARefLink{sga3:xxii:cor:4.1.6}{4.1.6} by
descent.
Let us prove the second assertion. Let \(H\) denote ``the
intersection of the maximal tori of \(G\)'' in the preceding sense.
Clearly,\footnote{Editor's note: We have expanded the original in
what follows.}
\[
  \underline{\operatorname{Centr}}(G)\subset H.
\]
By descent, it is therefore enough to prove
\(\underline{\operatorname{Centr}}(G)=H\) when \(G\) is split.
Since \(H\) is contained in the intersection of the maximal tori of
\(G\) in the usual sense, this follows from the following remark:
if \((G,T,M,R)\) is a splitting, \(\alpha\in R\), and
\(X\in\Gamma(S,\mathfrak g^\alpha)^\times\), then
\[
  \operatorname{int}\bigl(\exp_\alpha(X)\bigr)(T)\cap T
  =
  \operatorname{Ker}(\alpha),
\]
as a trivial calculation shows. See also \SGARefLink{sga3:xii:8.6}{Exposé XII, 8.6} and \SGARefLink{sga3:xii:8.8}{8.8}
for a more general statement.

\SGARefTarget{sga3:xxii:rem:4.1.8}\paragraph{\SGARefLink{sga3:xxii:rem:4.1.8}{Remark 4.1.8}.}
In what follows, in the split case we shall systematically identify
\(T\) with \(D_S(M)\). Then
\[
  \underline{\operatorname{Centr}}(G)
  =
  D_S\bigl(M/\Gamma_0(R)\bigr),
\]
where \(\Gamma_0(R)\) is the subgroup of \(M\) generated by \(R\)
\textup{(cf.\ \SGARefLink{sga3:XXI:1.1.6}{Exposé XXI, 1.1.6})}.
\par\noindent
If
\(\{\alpha_1,\alpha_2,\ldots,\alpha_n\}\) is a system of simple
roots of \(R\), then one immediately has
\textup{(cf.\ \SGARefLink{sga3:xx:1.19}{Exposé XX, 1.19})}
\[
  \underline{\operatorname{Centr}}(G)
  =
  \bigcap_i\operatorname{Ker}(\alpha_i)
  =
  \bigcap_i\underline{\operatorname{Centr}}(Z_{\alpha_i}).
\]

\SGARefTarget{sga3:xxii:prop:4.1.9}\paragraph{\SGARefLink{sga3:xxii:prop:4.1.9}{Proposition 4.1.9}.}
\textit{Let \(S\) be a scheme, let \((G,T,M,R)\) be a split
\(S\)-group, let \(Q\) be an \(S\)-torus, let \(\alpha_0\) be a
character of \(Q\), let \(\mathcal L\) be an invertible
\(\mathcal O_S\)-module, and let}
\[
  f:Q\longrightarrow T,
  \qquad
  p:W(\mathcal L)\longrightarrow G
\]
\textit{be group homomorphisms satisfying the following
set-theoretic relation:}
\[
  p\bigl(\alpha_0(q)x\bigr)
  =
  \operatorname{int}\bigl(f(q)\bigr)\mathbin{\cdot}p(x)
\]
\textit{for every \(q\in Q(S')\), \(x\in W(\mathcal L)(S')\), and
\(S'\to S\). Suppose that \(f\) separates the elements of \(R\) in
the following sense: if \(\alpha,\alpha'\in R\) and
\(m,m'\in\mathbb Z\), then}
\[
  m\alpha\circ f=m'\alpha'\circ f
  \quad\Longrightarrow\quad
  m\alpha=m'\alpha'.
\]
\textit{Finally, let \(s\in S\) be such that
\((\alpha_0)_{\bar s}\ne e\) and \(p_{\bar s}\ne e\).}

\textit{There then exist an open subset \(U\) of \(S\) containing
\(s\), an integer \(q>0\) such that \(x\mapsto x^q\) is an
endomorphism of \(\mathbb G_{a,U}\), a root \(\alpha\in R\), and an
isomorphism of \(\mathcal O_U\)-modules}
\[
  h:(\mathcal L|_U)^{\otimes q}
  \xrightarrow{\;\sim\;}
  \mathfrak g^\alpha|_U
\]
\textit{such that}
\[
  \textup{(i)}\quad
  (\alpha\circ f)|_U=(q\alpha_0)|_U,
\]
\[
  \textup{(ii)}\quad
  p(X)=\exp_\alpha\bigl(h(X^q)\bigr)
  \quad\textup{for every \(X\in W(\mathcal L)(S')\), \(S'\to U\).}
\]
\textit{Moreover, once \(U\) has been chosen, \(q\), \(\alpha\),
and \(h\) are uniquely determined.}

After shrinking \(S\), we may suppose that \(\alpha_0\) is nonzero
on every fiber of \(S\). Choose a system of positive roots \(R_+\)
of \(R\), and put
\[
  V=p^{-1}(\Omega_{R_+}).
\]
This is an open subset of \(W(\mathcal L)\) containing the zero
section and stable under multiplication by every
\(\alpha_0(q)\), where \(q\in Q(S')\) and \(S'\to S\). Since
\(\alpha_0\) is nontrivial on every fiber, it follows immediately
that \(V=W(\mathcal L)\), and hence that \(p\) factors through
\(\Omega_{R_+}\). Choose arbitrary orderings of \(R_+\) and
\(R_-\); every product below is understood to be taken in the
chosen ordering. There are therefore unique morphisms
\[
  a_\alpha:W(\mathcal L)\longrightarrow U_\alpha,
  \qquad \alpha\in R,
  \qquad
  b:W(\mathcal L)\longrightarrow T
\]
such that
\[
  p(x)
  =
  \prod_{\alpha\in R_-}a_\alpha(x)\mathbin{\cdot}b(x)
  \mathbin{\cdot}\prod_{\alpha\in R_+}a_\alpha(x).
\]
Writing out the covariance condition under \(Q\), one immediately
obtains
\[
  a_\alpha\bigl(\alpha_0(q)x\bigr)
  =
  \alpha\bigl(f(q)\bigr)a_\alpha(x),
  \qquad \alpha\in R,
\]
\[
  b\bigl(\alpha_0(q)x\bigr)=b(x)
\]
for every \(x\in W(\mathcal L)(S')\), \(q\in Q(S')\), and
\(S'\to S\). The second condition immediately gives \(b=e\).

Now let \(\alpha\in R\) be such that
\((a_\alpha)_{\bar s}\ne e\); such an \(\alpha\) exists because
\(p_{\bar s}\ne e\) by hypothesis. Applying \SGARefLink{sga3:xix:proposition:4:12:item:a}{Exposé XIX, 4.12(a)},
we find an integer \(n>0\) such that
\[
  (\alpha\circ f)_{\bar s}=(n\alpha_0)_{\bar s}.
\]
After shrinking \(S\), we may suppose that
\(\alpha\circ f=n\alpha_0\) (\SGARefLink{sga3:ix:corollary:5:3}{Exposé IX, 5.3}). But then, for every
\(\alpha'\in R\) with \(\alpha'\ne\alpha\), one has
\[
  (\alpha'\circ f)_{\bar s}\ne m\alpha_0
  \qquad\textup{for every integer \(m>0\),}
\]
by the hypothesis on \(f\) and the fact that the only roots
proportional to \(\alpha\) are \(\alpha\) and \(-\alpha\).
Applying \SGARefLink{sga3:xix:proposition:4:12:item:a}{Exposé XIX, 4.12(a)}, now to \(a_{\alpha'}\), we find that
\(a_{\alpha'}\) vanishes in a neighborhood of \(s\).
Since \(R\) is finite, after shrinking \(S\) once more we may
suppose that \(a_{\alpha'}=0\) for every
\(\alpha'\in R\), \(\alpha'\ne\alpha\). Thus \(p=a_\alpha\), and
we may apply \SGARefLink{sga3:xix:proposition:4-12:item:b}{Exposé XIX, 4.12(b)}, followed by \SGARefLink{sga3:xix:proposition:4:12:item:c}{4.12(c)}, which gives
the asserted result. The uniqueness assertions are clear.

\SGARefTarget{sga3:xxii:rem:4.1.10}\paragraph{\SGARefLink{sga3:xxii:rem:4.1.10}{Remark 4.1.10}.}
The condition imposed on \(f\) in \SGARefLink{sga3:xxii:prop:4.1.9}{4.1.9} is satisfied in particular
when \(f\) is surjective \textup{(that is, faithfully flat)}.

\SGARefTarget{sga3:xxii:prop:4.1.11}\paragraph{\SGARefLink{sga3:xxii:prop:4.1.11}{Proposition 4.1.11}.}
\textit{Let \((G,T,M,R)\) be a split \(S\)-group, let \(R_+\) be a
system of positive roots of \(R\), and let \(\Omega_{R_+}\) be the
corresponding ``big cell.''}

\textit{\textup{\SGARefLink{sga3:xxii:def:1.13:item:i}{(i)}} Let \(H\) be a group-valued \(S\)-functor
that is separated
for the fppf topology. If
\(f,g:G\rightrightarrows H\) are two group homomorphisms that
coincide on \(\Omega_{R_+}\), then \(f=g\).}

\SGAInlineRefTarget{sga3:xxii:prop:4.1.11:item:ii}\textit{\textup{\SGARefLink{sga3:xxii:def:1.13:item:ii}{(ii)}} Let \(H\) be an \(S\)-sheaf of groups for the
fppf topology, and let \(f:\Omega_{R_+}\to H\) be an \(S\)-morphism
with the following property: for every \(S'\to S\) and every
\(x,y\in\Omega_{R_+}(S')\) such that
\(xy\in\Omega_{R_+}(S')\), one has
\(f(xy)=f(x)f(y)\). Then there exists a group homomorphism
\(\bar f:G\to H\), unique by \textup{\SGARefLink{sga3:xxii:def:1.13:item:i}{(i)}}, extending \(f\).}

Indeed, by \SGARefLink{sga3:xxii:prop:4.1.5}{4.1.5}, \textup{\SGARefLink{sga3:xxii:def:1.13:item:i}{(i)}} \textup{(respectively
\textup{\SGARefLink{sga3:xxii:def:1.13:item:ii}{(ii)}})} follows immediately from \SGARefLink{sga3:xviii:cor:2.2}{Exposé XVIII, 2.2}
\textup{(respectively \SGARefLink{sga3:xxii:cor:2.3}{2.3} and \SGARefLink{sga3:xxii:cor:2.4}{2.4})}.

\SGARefTarget{sga3:xxii:rem:4.1.12}\paragraph{\SGARefLink{sga3:xxii:rem:4.1.12}{Remark 4.1.12}.}
If \(\alpha\in R_+\), then
\[
  \tag{\(\dagger\)}\label{sga3:xxii:subsection:4-1:equation:dagger}
  \Omega_{R_+}\cap Z_\alpha
  =
  U_{-\alpha}\mathbin{\cdot}T\mathbin{\cdot}U_\alpha.
\]
Indeed,\footnote{Editor's note: We have expanded the original in
what follows.} for every \(S'\to S\), if
\[
  g=
  \prod_{\beta\in R_-}p_\beta(x_\beta)
  \mathbin{\cdot}t\mathbin{\cdot}
  \prod_{\beta\in R_+}p_\beta(x_\beta)
\]
is an element of \(\Omega_{R_+}(S')\), and if
\(t'\in T_\alpha(S')\), then
\[
  t'gt'^{-1}
  =
  \prod_{\beta\in R_-}p_\beta\bigl(\beta(t')x_\beta\bigr)
  \mathbin{\cdot}t\mathbin{\cdot}
  \prod_{\beta\in R_+}p_\beta\bigl(\beta(t')x_\beta\bigr).
\]
Since \(\alpha\) and \(-\alpha\) are the only elements of \(R\)
whose restriction to \(T_\alpha\) is \(1\), it follows that \(g\)
belongs to \(Z_\alpha=\operatorname{Centr}(T_\alpha)\) if and only
if \(x_\beta=0\) for \(\beta\ne\pm\alpha\).

From \((\dagger)\) and \SGARefLink{sga3:xx:2.1}{Exposé XX, 2.1}, one deduces that if
\(X\in\Gamma(S,\mathfrak g^\alpha)\) and
\(Y\in\Gamma(S,\mathfrak g^{-\alpha})\), then
\[
  \exp_\alpha(X)\exp_{-\alpha}(Y)\in\Omega_{R_+}(S)
  \quad\Longleftrightarrow\quad
  1+XY\ \text{is invertible}.
\]

\SGARefTarget{sga3:xxii:section:4.2}
\subsection*{\SGAInlineRefTarget{sga3:xxii:subsection:4-2}4.2. Morphisms of Split Groups}
\addcontentsline{toc}{subsection}{\SGARefLink{sga3:xxii:subsection:4-2}{4.2}. Morphisms of Split Groups}

\SGARefTarget{sga3:xxii:def:4.2.1}\paragraph{\SGARefLink{sga3:xxii:def:4.2.1}{Definition 4.2.1}.}
Let \(S\) be a nonempty scheme, and let
\((G,T,M,R)\) and \((G',T',M',R')\) be two split \(S\)-groups.
An \(S\)-group homomorphism \(f:G\to G'\) is said to be
\emph{compatible with the splittings}, or to define a
\emph{morphism of split groups}, if the restriction of \(f\) to
\(T\) factors as a morphism \(f_T:T\to T'\) of the form
\(f_T=D_S(h)\), where \(h:M'\to M\) is a group homomorphism
satisfying the following condition: there exist a bijection
\(d:R\xrightarrow{\sim}R'\) and, for
every \(\alpha\in R\), an integer \(q(\alpha)>0\) such that
\(x\mapsto x^{q(\alpha)}\) is an endomorphism of
\(\mathbb G_{a,S}\) and
\[
  h\bigl(d(\alpha)\bigr)=q(\alpha)\alpha,
  \qquad
  {}^t h(\alpha^*)=q(\alpha)d(\alpha)^*.
\]

\SGARefTarget{sga3:xxii:rem:4.2.2}\paragraph{\SGARefLink{sga3:xxii:rem:4.2.2}{Remark 4.2.2}.}
It is immediate that \(h\), \(d\), and the \(q(\alpha)\), for
\(\alpha\in R\), are uniquely determined by \(f\). We write
\(h=\mathscr R(f)\). The \(q(\alpha)\) are the radicial exponents
of \(f\) \textup{(or of \(h\))}.

Let \(p\) be the prime number, if one exists, that vanishes on
\(S\); put \(p=1\) if no prime number vanishes on \(S\). Then
\(\mathscr R(f)\) is a \(p\)-morphism of reduced root data in the
sense of \SGARefLink{sga3:XXI:6.8}{Exposé XXI, 6.8}. We have thus defined a functor
\(\mathscr R\) from the category of split \(S\)-groups to the
category of reduced root data \textup{(equipped with
\(p\)-morphisms)}.

\SGARefTarget{sga3:xxii:prop:4.2.3}\paragraph{\SGARefLink{sga3:xxii:prop:4.2.3}{Proposition 4.2.3}.}
\textit{Under the hypotheses of \SGARefLink{sga3:xxii:def:4.2.1}{4.2.1}, the following properties
hold:}

\textit{\textup{\SGARefLink{sga3:xxii:def:1.13:item:i}{(i)}} For every \(\alpha\in R\), there exists a
unique isomorphism of \(\mathcal O_S\)-modules}
\[
  f_\alpha:
  (\mathfrak g^\alpha)^{\otimes q(\alpha)}
  \xrightarrow{\;\sim\;}
  \mathfrak g'^{\,d(\alpha)}
\]
\textit{such that}
\[
  f\bigl(\exp_\alpha(X)\bigr)
  =
  \exp_{d(\alpha)}
  \bigl(f_\alpha(X^{q(\alpha)})\bigr)
\]
\textit{for every \(X\in W(\mathfrak g^\alpha)(S')\) and every
\(S'\to S\).}

\textit{\textup{\SGARefLink{sga3:xxii:def:1.13:item:ii}{(ii)}} For every \(\alpha\in R\), one has
\(q(-\alpha)=q(\alpha)\), and \(f_\alpha\) and \(f_{-\alpha}\) are
contragredient to one another.}

\textit{\textup{\SGARefLink{sga3:xxii:def:1.9:item:iii}{(iii)}} For every \(\alpha\in R\), every
\(Z\in W(\mathfrak g^\alpha)^*(S')\), and every \(S'\to S\), one
has}
\[
  f\bigl(w_\alpha(Z)\bigr)
  =
  w_{d(\alpha)}\bigl(Z^{q(\alpha)}\bigr).
\]

By hypothesis, the following diagram is commutative:

\SGARefTarget{sga3:xxii:diagram:4.2.3}
\begin{center}
  \begin{tikzcd}[column sep=4.6em,row sep=3.2em]
  \mathbb G_{m,S}
    \arrow[r,"\alpha^*"]
    \arrow[d,"q(\alpha)"']
  &
  T
    \arrow[r,"\alpha"]
    \arrow[d,"f_T"']
  &
  \mathbb G_{m,S}
    \arrow[d,"q(\alpha)"]
  \\
  \mathbb G_{m,S}
    \arrow[r,"d(\alpha)^*"]
  &
  T'
    \arrow[r,"d(\alpha)"]
  &
  \mathbb G_{m,S}
\end{tikzcd}

\end{center}

It follows that \(f\) maps \(\operatorname{Ker}(\alpha)\) into
\(\operatorname{Ker}(d(\alpha))\), hence \(T_\alpha\) into
\(T'_{d(\alpha)}\), and hence \(Z_\alpha\) into
\(Z'_{d(\alpha)}\). It remains only to apply \SGARefLink{sga3:xx:3.10}{Exposé XX, 3.10} and
\SGARefLink{sga3:xx:3.11}{3.11}, to the groups \(Z_\alpha\) and \(Z'_{d(\alpha)}\).

\begingroup
\emergencystretch=2em
\SGARefTarget{sga3:xxii:prop:4.2.4}\paragraph{\SGARefLink{sga3:xxii:prop:4.2.4}{Proposition 4.2.4}.}
\textit{The morphism \(f\) induces a morphism \(f_N\) from
\(\underline{\operatorname{Norm}}_G(T)\) to
\(\underline{\operatorname{Norm}}_{G'}(T')\), hence a morphism
\(f_W\) from \(W_G(T)\) to \(W_{G'}(T')\); the latter is an
isomorphism. More precisely, let}
\[
  \bar d:W\bigl(\mathscr R(G)\bigr)=W
  \longrightarrow
  W'=W\bigl(\mathscr R(G')\bigr)
\]
\textit{be the isomorphism extending
\(s_\alpha\mapsto s_{d(\alpha)}\) \textup{(\SGARefLink{sga3:XXI:6.8.4}{Exposé XXI, 6.8.4})}.
Then there is a commutative diagram of isomorphisms:}
\endgroup

\SGARefTarget{sga3:xxii:diagram:4.2.4}
\begin{center}
  \begin{tikzcd}[column sep=5.2em,row sep=3.2em]
  W_G(T)
    \arrow[r,"f_W","\sim"']
  &
  W_{G'}(T')
  \\
  W_S
    \arrow[u,"\sim"']
    \arrow[r,"\bar d_S","\sim"']
  &
  W'_S.
    \arrow[u,"\sim"']
\end{tikzcd}

\end{center}

This follows immediately from \SGARefLink{sga3:xxii:prop:3.4}{3.4}, \SGARefLink{sga3:XXI:6.8.4}{Exposé XXI, 6.8.4}, and
\textup{\SGARefLink{sga3:xxii:def:1.9:item:iii}{(iii)}} above.

\SGARefTarget{sga3:xxii:rem:4.2.5}\paragraph{\SGARefLink{sga3:xxii:rem:4.2.5}{Remark 4.2.5}.}
With the notation of \SGARefLink{sga3:xxii:prop:4.2.3}{4.2.3}, the restriction of \(f\) to
\(\Omega_{R_+}\), for a system of positive roots \(R_+\), can be
written explicitly: it maps \(\Omega_{R_+}\) into
\(\Omega'_{d(R_+)}\) \textup{(\(d(R_+)\) is a system of positive
roots of \(R'\) by \SGARefLink{sga3:XXI:6.8.7}{Exposé XXI, 6.8.7})}, and is given by the
set-theoretic formula
\[
\begin{aligned}
  f\Biggl(
    \prod_{\alpha\in R_-}\exp_\alpha(X_\alpha)
    \mathbin{\cdot}t\mathbin{\cdot}
    \prod_{\alpha\in R_+}\exp_\alpha(X_\alpha)
  \Biggr)
  &=
  \prod_{\alpha\in R_-}
    \exp_{d(\alpha)}
    \Bigl(f_\alpha(X_\alpha^{q(\alpha)})\Bigr)
  \mathbin{\cdot}f_T(t)\mathbin{\cdot}\\
  &\qquad
  \prod_{\alpha\in R_+}
    \exp_{d(\alpha)}
    \Bigl(f_\alpha(X_\alpha^{q(\alpha)})\Bigr).
\end{aligned}
\]

\SGARefTarget{sga3:xxii:prop:4.2.6}\paragraph{\SGARefLink{sga3:xxii:prop:4.2.6}{Proposition 4.2.6}.}
\textit{\textup{\SGARefLink{sga3:xxii:def:1.13:item:i}{(i)}} The morphism \(f\) is surjective
\textup{(\(=\) faithfully flat in the present case; cf.\
Exposé \(\mathrm{VI_B}\), 3.11)} if and only if \(f_T\) is
surjective.}

\textit{\textup{\SGAInlineRefTarget{sga3:xxii:prop:4.2.6:item:ii}(ii)} One has
\(\operatorname{Ker}(f)\subset\Omega_{R_+}\).}

We prove \textup{\SGARefLink{sga3:xxii:def:1.13:item:i}{(i)}}. If \(f\) is surjective, then
\(f_T(T)=f(T)\) is a maximal torus of \(G'\). Indeed, \(f(T)\) is
a subtorus of a maximal torus \(T'\) \textup{(\SGARefLink{sga3:ix:theorem:6:8}{Exposé IX, 6.8})}; to
verify that \(f(T)=T'\), one reduces to the case of an algebraically
closed field, where this is \emph{Bible}, \S7.3, Theorem 3\textup{(a)}.

If \(f_T\) is surjective, the preceding formula shows that \(f\)
induces a surjection from \(\Omega=\Omega_{R_+}\) onto
\(\Omega'=\Omega'_{d(R_+)}\). Since the fibers
of \(G'\) are connected, it follows \textup{(cf.\ Exposé
\(\mathrm{VI_A}\), 0.5)} that \(f\) is surjective.

We prove \textup{\SGARefLink{sga3:xxii:prop:4.2.6:item:ii}{(ii)}}, and for this purpose admit a result that will
be proved below \textup{(\SGARefLink{sga3:xxii:thm:5.7.4}{5.7.4})}. For each \(w\in W\), choose an
\(n_w\in\underline{\operatorname{Norm}}_G(T)(S)\) representing it.
Then the open subschemes \(n_w\Omega\), for \(w\in W\), cover \(G\).
It is therefore enough to prove that
\[
  \operatorname{Ker}(f)\cap n_w\Omega\ne\varnothing
  \quad\Longrightarrow\quad w=1.
\]
If \(x\in\Omega(S')\), where \(S'\to S\), and \(f(n_wx)=1\), then
\(f(x)=f(n_w)^{-1}\). But \(f(x)\in\Omega'(S')\), whereas
\[
  f(n_w)^{-1}\in
  \underline{\operatorname{Norm}}_{G'}(T')(S').
\]
By \SGARefLink{sga3:xxii:prop:4.2.4}{4.2.4}, we are reduced to proving the following.

\SGARefTarget{sga3:xxii:lem:4.2.7}\paragraph{\SGARefLink{sga3:xxii:lem:4.2.7}{Lemma 4.2.7}.}
\textit{Under the hypotheses of \SGARefLink{sga3:xxii:prop:4.1.2}{4.1.2}, one has}
\[
  \Omega\cap\underline{\operatorname{Norm}}_G(T)=T.
\]

Let
\[
  x=
  \prod_{\alpha\in R_-}p_\alpha(x_\alpha)
  \mathbin{\cdot}t\mathbin{\cdot}
  \prod_{\alpha\in R_+}p_\alpha(x_\alpha)
  =
  vtu
  \in\Omega(S').
\]
If \(x\) normalizes \(T_{S'}\), then, for every \(t'\in T(S')\),
\[
  xt'x^{-1}=t''\in T(S'),
\]
that is, \(xt'=t''x\), which may be written
\[
  v(tt')(t'^{-1}ut')
  =
  (t''vt''^{-1})(t''t)u.
\]
This gives \(t'^{-1}ut'=u\), hence
\[
  u\in\operatorname{Centr}_G(T)(S')=T(S'),
\]
and therefore \(u=1\). Similarly, \(v=1\).

\SGARefTarget{sga3:xxii:cor:4.2.8}\paragraph{\SGARefLink{sga3:xxii:cor:4.2.8}{Corollary 4.2.8}.}
\textit{One has}
\[
  \operatorname{Ker}(f)
  =
  \prod_{\alpha\in R_-}K_\alpha
  \mathbin{\cdot}\operatorname{Ker}(f_T)\mathbin{\cdot}
  \prod_{\alpha\in R_+}K_\alpha,
\]
\textit{where, for each \(\alpha\in R\), \(K_\alpha\) denotes the
finite \(S\)-group}
\[
  K_\alpha
  =
  \operatorname{Ker}\!\left(
    U_\alpha\longrightarrow U_\alpha^{\otimes q(\alpha)}
  \right)
  \simeq
  \boldsymbol{\alpha}_{q(\alpha),S}.
\]

To apply this corollary, we make the following definition.

\SGARefTarget{sga3:xxii:def:4.2.9}\paragraph{\SGARefLink{sga3:xxii:def:4.2.9}{Definition 4.2.9}.}
Let \(S\) be a scheme and let \(G\) and \(G'\) be two reductive
\(S\)-groups. A faithfully flat and finite morphism of \(S\)-groups
\(f:G\to G'\) \textup{(i.e.\ one that is surjective and has finite
kernel over \(S\))} is called an \emph{isogeny}. If, moreover,
\(\operatorname{Ker}(f)\) is a central subgroup of \(G\), then \(f\)
is called a \emph{central isogeny}.

\SGARefTarget{sga3:xxii:prop:4.2.10}\paragraph{\SGARefLink{sga3:xxii:prop:4.2.10}{Proposition 4.2.10}.}
\textit{Let \(f:G\to G'\) be a morphism of split groups. For \(f\)
to be an isogeny \textup{(respectively, a central isogeny)}, it is
necessary and sufficient that \(f_T\) be an isogeny, i.e.\ that
\(\mathscr R(f)\) be injective with finite cokernel
\textup{(respectively, and that \(q(\alpha)=1\) for every
\(\alpha\in R\))}.}

Indeed, by \SGARefLink{sga3:xxii:cor:4.2.8}{4.2.8}, \(\operatorname{Ker}(f)\) is finite over \(S\)
if and only if \(\operatorname{Ker}(f_T)\) is finite over \(S\),
and \(\operatorname{Ker}(f)\subset T\) if and only if every
\(q(\alpha)\) is equal to \(1\). In that case
\(\operatorname{Ker}(f)\) is central, since it is of multiplicative
type and invariant; cf.\ \SGARefLink{sga3:ix:corollary:5:5}{Exposé IX, 5.5}.

\SGARefTarget{sga3:xxii:rem:4.2.11}\paragraph{\SGARefLink{sga3:xxii:rem:4.2.11}{Remark 4.2.11}.}
\textup{(a)} Thus \(f:G\to G'\) is a central isogeny if and only if
\[
  \mathscr R(f):\mathscr R(G')\longrightarrow\mathscr R(G)
\]
is an isogeny in the sense of \SGARefLink{sga3:XXI:6.2}{Exposé XXI, 6.2}. Moreover, in this
case, with the notation of \emph{loc.\ cit.}, one has
\[
  \operatorname{Ker}(f)
  =
  D_S\!\left(K\bigl(\mathscr R(f)\bigr)\right),
  \qquad
  K\bigl(\mathscr R(f)\bigr)
  =
  \operatorname{Coker}\bigl(\mathscr R(f)\bigr).
\]

\textup{\SGARefLink{sga3:xxii:def:1.5:item:b}{(b)}} If \(G\) and \(G'\) are semisimple, every morphism of
split groups \(f:G\to G'\) is an isogeny.

\textup{(c)} If \(f:G\to G'\) is faithfully flat and finite and
\(G\) is reductive \textup{(respectively, semisimple)}, then \(G'\)
is so as well. Indeed, \(G'\) is of finite presentation over \(S\)
\textup{(Exposé V, 9.1)}, affine over \(S\)
\textup{(EGA II, 6.7.1)}, smooth over \(S\)
\textup{(Exposé VI, 9.2)}, and has connected reductive
\textup{(respectively, semisimple)} geometric fibers by \SGARefLink{sga3:xix:paragraph:1-7}{Exposé XIX,
1.7}.

\SGARefLink{sga3:xxii:def:4.2.1}{Definition 4.2.1} may seem arbitrary. It is justified by the
following proposition, which, for simplicity, we shall state for
semisimple groups.

A morphism \(f:G\to G'\) of reductive \(S\)-groups is said to be
\emph{splittable} if there exist splittings of \(G\) and \(G'\)
with which \(f\) is compatible. We then have the following.

\SGARefTarget{sga3:xxii:prop:4.2.12}\paragraph{\SGARefLink{sga3:xxii:prop:4.2.12}{Proposition 4.2.12}.}
\textit{Let \(S\) be a scheme, let \(G\) and \(G'\) be two
semisimple \(S\)-groups, and let \(f:G\to G'\) be a morphism of
groups. Let \(s\in S\). The following conditions are equivalent:}

\begin{enumerate}
\item[\textup{\SGAInlineRefTarget{sga3:xxii:prop:4.2.12:item:i}(i)}]
\(\operatorname{Ker}(f_{\bar s})\) is finite
\textup{(\(\Longleftrightarrow e\) is isolated in
\(\operatorname{Ker}(f)(\bar s)\))}, and \(f_{\bar s}\) is
surjective, i.e.\ \(f_{\bar s}\) is an isogeny.

\item[\textup{\SGAInlineRefTarget{sga3:xxii:prop:4.2.12:item:ii}(ii)}]
\(f_{\bar s}\) is splittable.

\item[\textup{\SGAInlineRefTarget{sga3:xxii:prop:4.2.12:item:iii}(iii)}]
There exists an étale morphism \(S'\to S\) covering \(s\) such that
\(f_{S'}:G_{S'}\to G'_{S'}\) is splittable.
\end{enumerate}

Clearly, \textup{\SGARefLink{sga3:xxii:prop:4.2.12:item:iii}{(iii)}} is equivalent to \textup{\SGARefLink{sga3:xxii:prop:4.2.12:item:ii}{(ii)}}, and
\textup{\SGARefLink{sga3:xxii:prop:4.2.12:item:ii}{(ii)}} implies \textup{\SGARefLink{sga3:xxii:prop:4.2.12:item:i}{(i)}} by \SGARefLink{sga3:xxii:rem:4.2.11}{4.2.11}\textup{\SGARefLink{sga3:xxii:def:1.5:item:b}{(b)}}. This is
where the hypothesis that \(G\) and \(G'\) are semisimple enters;
the other implications are valid for reductive groups.

We now prove that \textup{\SGARefLink{sga3:xxii:prop:4.2.12:item:i}{(i)}} implies \textup{\SGARefLink{sga3:xxii:prop:4.2.12:item:iii}{(iii)}}. We may
suppose that \(G\) and \(G'\) are split in such a way that \(f\)
induces a morphism
\[
  f_T:T\longrightarrow T'
\]
\textup{(\SGARefLink{sga3:xxii:cor:2.3}{2.3} and \SGARefLink{sga3:xix:paragraph:2:8}{Exposé XIX, 2.8})}. After restricting \(S\), we may
suppose that \(f_T=D_S(h)\), where \(h:M'\to M\) is a morphism of
groups. Let \(\alpha\in R\), and consider the composite morphism
\[
  p:
  W(\mathfrak g^\alpha)
  \xrightarrow{\ \exp_\alpha\ }
  G
  \xrightarrow{\ f\ }
  G'.
\]
Since \(\operatorname{Ker}(p_{\bar s})\) is finite,
\(p_{\bar s}\ne e\). On the other hand, \(f_{T_{\bar s}}\) is
surjective. We may therefore apply \SGARefLink{sga3:xxii:prop:4.1.9}{4.1.9}: there exist an open
subset \(V_\alpha\) of \(S\) containing \(s\), a root
\(\alpha'\in R'\), and an integer \(q(\alpha)\) such that
\(x\mapsto x^{q(\alpha)}\) is an endomorphism of
\(\mathbb G_{a,V_\alpha}\), and an isomorphism of
\(\mathcal O_{V_\alpha}\)-modules
\[
  f_\alpha:
  \left.(\mathfrak g^\alpha)^{\otimes q(\alpha)}
  \right|_{V_\alpha}
  \xrightarrow{\ \sim\ }
  \left.\mathfrak g'^{\alpha'}\right|_{V_\alpha},
\]
such that
\[
  f\bigl(\exp_\alpha(X_\alpha)\bigr)
  =
  \exp_{d(\alpha)}
  \bigl(f_\alpha(X_\alpha^{q(\alpha)})\bigr)
\]
and
\[
  \alpha'\circ f_T=h(\alpha')=q(\alpha)\alpha.
\]
We may replace \(S\) by the intersection of the \(V_\alpha\), for
\(\alpha\in R\). Put \(\alpha'=d(\alpha)\). It is clear that
\[
  d:R\longrightarrow R'
\]
is a bijection, since the kernel of \(h\) is finite
\textup{(\(f_{T_{\bar s}}\) being surjective)}. It remains only to
prove that
\[
  f_T\circ\alpha^*=q(\alpha)\alpha'^*,
\]
which follows by a trivial modification of the argument used in
\SGARefLink{sga3:xx:3.11}{Exposé XX, 3.11}.

In any event, as we have seen in the course of the proof,
\textup{\SGARefLink{sga3:xxii:prop:4.2.12:item:i}{(i)}} implies \textup{\SGARefLink{sga3:xxii:prop:4.2.12:item:iii}{(iii)}}. We have therefore proved the
following.

\SGARefTarget{sga3:xxii:cor:4.2.13}\paragraph{\SGARefLink{sga3:xxii:cor:4.2.13}{Corollary 4.2.13}.}
\textit{Let \(S\) be a scheme and let \(f:G\to G'\) be an isogeny
of reductive groups. Then \(f\) is locally splittable for the
étale topology.}

\SGARefTarget{sga3:xxii:section:4.3}\section*{4.3. Central
Quotients of Reductive Groups}
\addcontentsline{toc}{section}{\SGARefLink{sga3:xxii:section:4.3}{4.3}. Central Quotients of Reductive
Groups}

We first consider a special case.

\SGARefTarget{sga3:xxii:prop:4.3.1}\paragraph{\SGARefLink{sga3:xxii:prop:4.3.1}{Proposition 4.3.1}.}
\textit{Let \(S\) be a scheme, let \((G,T,M,R)\) be a split
\(S\)-group, let \(N\) be a subgroup of \(M\) containing \(R\),
and put}
\[
  Q=D_S(M/N)\subset\operatorname{Centr}(G).
\]
\textit{Then:}

\begin{enumerate}
\item[\textup{(i)}]
\textit{\(G'=G/Q\) is a reductive \(S\)-group, and \(T'=T/Q\) is a
maximal torus of \(G'\);}

\item[\textup{\SGAInlineRefTarget{sga3:xxii:prop:4.3.1:item:ii}(ii)}]
\textit{if \(T'\) is identified with \(D_S(N)\), then \(R\subset N\)
is a root system of \(G'\) relative to \(T'\),
\((G',T',N,R)\) is a splitting of \(G'\), and
\(\mathscr R(G')\) is canonically identified with the induced root
datum \textup{(\SGARefLink{sga3:XXI:6.5}{Exposé XXI, 6.5})} \(\mathscr R(G)_N\);}

\item[\textup{(iii)}]
\textit{the canonical morphism \(G\to G'\) is compatible with the
splittings, has root exponents equal to \(1\), and, by
functoriality, induces the canonical morphism \textup{(loc. cit.)}}
\[
  \mathscr R(G)_N\longrightarrow\mathscr R(G).
\]
\end{enumerate}

We know that \(G'=G/Q\) is representable by an affine group scheme
over \(S\) \textup{(\SGARefLink{sga3:viii:corollary:5:7}{Exposé VIII, 5.7})}, smooth over \(S\)
\textup{(Exposé VI\textsubscript{B}, \SGARefLink{prop:sga3-VIB-9.2}{9.2})}, and with connected
reductive geometric fibers \textup{(as quotients of reductive
groups; cf.\ \SGARefLink{sga3:xix:paragraph:1-7}{Exposé XIX, 1.7})}. Thus \(G'\) is a reductive
\(S\)-group.

It is clear that
\[
  T'=T/Q\simeq D_S(N)
\]
is a maximal torus of \(G'\). Next, choose a system \(R_+\) of
positive roots of \(R\). The open subset \(\Omega_{R_+}\) of \SGARefLink{sga3:xxii:subsection:4-2}{4.2} is
stable under \(Q\), and there is a canonical isomorphism
\[
  \Omega_{R_+}/Q
  \simeq
  \left(\prod_{\alpha\in R_-}U_\alpha\right)
  \mathop{\times}_S(T/Q)
  \mathop{\times}_S
  \left(\prod_{\alpha\in R_+}U_\alpha\right).
\]
Moreover, \(\Omega_{R_+}/Q\) is an open subset of \(G'\) containing
the identity section \textup{(cf.\ \SGARefLink{sga3:iv:proposition:4:7:2}{Exposé IV, 4.7.2} and \SGARefLink{sga3:iv:lemma:6:4:1}{6.4.1})}.

It follows that, if \(\mathfrak g'\) denotes the Lie algebra of
\(G'\) and \(\bar\alpha\) the character of \(T/Q\) induced by
\(\alpha\) \textup{(or, equivalently, defined by \(\alpha\in N\)
under the identification \(T/Q=D_S(N)\))}, then the canonical
morphism \(\mathfrak g\to\mathfrak g'\) induces, for every
\(\alpha\in R\), an isomorphism
\[
  \mathfrak g^\alpha
  \xrightarrow{\ \sim\ }
  \mathfrak g'^{\bar\alpha}.
\]

We have thus proved that \(R\) is a root system of \(G'\) relative
to \(T'\), and the proof is completed without difficulty.

\SGARefTarget{sga3:xxii:cor:4.3.2}\paragraph{\SGARefLink{sga3:xxii:cor:4.3.2}{Corollary 4.3.2}.}
\textit{Let \(S\) be a scheme, let \(G\) be a reductive
\(S\)-group, and let \(Q\) be a normal subgroup of \(G\) of
multiplicative type. Then \(Q\) is central in \(G\), the quotient
\(G/Q\) is representable by a reductive \(S\)-group, and the
canonical morphism \(G\to G/Q\) is locally splittable for the
étale topology \textup{(with root exponents equal to \(1\))}.}

The first assertion follows from \SGARefLink{sga3:ix:corollary:5:5}{Exposé IX, 5.5}; the others are
local for the étale topology and reduce to \SGARefLink{sga3:xxii:prop:4.3.1}{4.3.1}.

\SGARefTarget{sga3:xxii:def:4.3.3}\paragraph{\SGARefLink{sga3:xxii:def:4.3.3}{Definition 4.3.3}.}
Let \(G\) be a reductive \(S\)-group. We say that \(G\) is
\emph{adjoint} \textup{(respectively, \emph{simply connected})} if,
for every \(s\in S\), the type of \(G\) at \(s\) is given by an
adjoint \textup{(respectively, simply connected)} root datum, i.e.\
\textup{(\SGARefLink{sga3:XXI:6.2.6}{Exposé XXI, 6.2.6})} one for which \(M\) is generated by
\(R\) \textup{(respectively, \(M^*\) is generated by \(R^*\))}.

\SGARefTarget{sga3:xxii:prop:4.3.4}\paragraph{\SGARefLink{sga3:xxii:prop:4.3.4}{Proposition 4.3.4}.}
\textit{\textup{\SGAInlineRefTarget{sga3:xxii:prop:4.3.4:item:i}(i)} An adjoint \textup{(respectively, simply
connected)} reductive group is semisimple.}

\textit{\textup{\SGARefLink{sga3:xxii:prop:4.3.1:item:ii}{(ii)}} If \(T\) is a maximal torus of the adjoint
\textup{(respectively, simply connected)} reductive group \(G\),
and if \(\alpha\) is a root of \(G\) relative to \(T\), then the
infinitesimal root \(\bar\alpha\) is nonzero on every fiber
\textup{(respectively, \(\alpha^*\) is a monomorphism and the
infinitesimal coroot \(H_\alpha\) is nonzero on every fiber)}.}

Indeed, \textup{\SGARefLink{sga3:xxii:prop:4.3.4:item:i}{(i)}} is trivial; \textup{\SGARefLink{sga3:xxii:prop:4.3.1:item:ii}{(ii)}} may be checked on
the geometric fibers and follows at once from \SGARefLink{sga3:XXI:6.2.8}{Exposé XXI, 6.2.8}.

\SGARefTarget{sga3:xxii:prop:4.3.5}\paragraph{\SGARefLink{sga3:xxii:prop:4.3.5}{Proposition 4.3.5}.}
\textit{\textup{\SGARefLink{sga3:xxii:prop:4.3.4:item:i}{(i)}} The reductive group \(G\) is adjoint if and
only if}
\[
  \operatorname{Centr}(G)=\{e\}_S.
\]

\textit{\textup{\SGARefLink{sga3:xxii:prop:4.3.1:item:ii}{(ii)}} For every reductive group \(G\), the quotient
group \(G/\operatorname{Centr}(G)\) is an adjoint reductive group.}

Indeed, we may suppose that \(G\) is split. Then \textup{\SGARefLink{sga3:xxii:prop:4.3.4:item:i}{(i)}} is
immediate, since
\[
  \operatorname{Centr}(G)=D_S\bigl(M/\Gamma_0(R)\bigr),
\]
and \textup{\SGARefLink{sga3:xxii:prop:4.3.1:item:ii}{(ii)}} follows from \SGARefLink{sga3:xxii:prop:4.3.1}{4.3.1}.

\SGARefTarget{sga3:xxii:def:4.3.6}\paragraph{\SGARefLink{sga3:xxii:def:4.3.6}{Definition 4.3.6}.}
Let \(G\) be a reductive \(S\)-group. The \emph{adjoint group} of
\(G\), denoted by \(\operatorname{ad}(G)\), is the group
\(G/\operatorname{Centr}(G)\). The \emph{radical} of \(G\), denoted
by \(\operatorname{rad}(G)\), is the maximal torus
\textup{(unique by \SGARefLink{sga3:xii:1.12}{Exposé XII, 1.12})} of
\(\operatorname{Centr}(G)\). The \emph{semisimple group associated
with \(G\)} is the quotient \(G/\operatorname{rad}(G)\).

The preceding definitions are compatible with base change. If
\(s\in S\), then
\(\operatorname{rad}(G)_{\bar s}\) is indeed the radical of
\(G_{\bar s}\) in the usual sense \textup{(\SGARefLink{sga3:xii:1.6}{Exposé XIX, 1.6})}.

\SGARefTarget{sga3:xxii:par:4.3.7}\paragraph{\SGARefLink{sga3:xxii:par:4.3.7}{4.3.7}.}
If \((G,T,M,R)\) is a split group, then
\[
  \operatorname{rad}(G)=D_S(M/N),
  \qquad
  N=M\cap V(R).
\]
Thus the semisimple group associated with \(G\)
\textup{(as also the adjoint group of \(G\))} is equipped with a
canonical splitting \textup{(\SGARefLink{sga3:xxii:prop:4.3.1}{4.3.1})}, and we have a diagram of
split groups
\[
  G\longrightarrow \operatorname{ss}(G)
   \longrightarrow \operatorname{ad}(G)
\]
corresponding to the canonical diagram of root data
\textup{(\SGARefLink{sga3:XXI:6.5.5}{Exposé XXI, 6.5.5})}
\[
  \operatorname{ad}\bigl(\mathscr R(G)\bigr)
  \longrightarrow
  \operatorname{ss}\bigl(\mathscr R(G)\bigr)
  \longrightarrow
  \mathscr R(G).
\]

\SGARefTarget{sga3:xxii:rem:4.3.8}\paragraph{\SGARefLink{sga3:xxii:rem:4.3.8}{Remark 4.3.8}.}
\textit{Let \((G,T,M,R)\) be a split adjoint
\textup{(respectively, simply connected)} \(S\)-group, and let
\(\Delta\) be a system of simple roots of \(R\). Then the family
\(\{\alpha\}_{\alpha\in\Delta}\), respectively
\(\{\alpha^*\}_{\alpha\in\Delta}\), induces an isomorphism}
\[
  T\xrightarrow{\ \sim\ }(\mathbb G_{m,S})^\Delta,
  \qquad\textit{respectively}\qquad
  (\mathbb G_{m,S})^\Delta\xrightarrow{\ \sim\ }T.
\]

Indeed,
\[
  M=\Gamma_0(R)
  \qquad
  \textup{(respectively, \(M^*=\Gamma_0(R^*)\))},
\]
and \(\Delta\) \textup{(respectively,
\(\{\alpha^*\}_{\alpha\in\Delta}\))} is a basis of the free
abelian group \(\Gamma_0(R)\)
\textup{(respectively, \(\Gamma_0(R^*)\))}; see \SGARefLink{sga3:XXI:3.1.8}{Exposé XXI, 3.1.8}.

\SGARefTarget{sga3:xxii:rem:4.3.9}\paragraph{\SGARefLink{sga3:xxii:rem:4.3.9}{Remark 4.3.9}.}
The radical of a reductive group is a \emph{characteristic}
subgroup \textup{(i.e.\ stable under
\(\underline{\operatorname{Aut}}_{S\text{-gr.}}(G)\))}, by its
definition.

\SGARefTarget{sga3:xxii:section:5}\section*{5. Subgroups of Type
\textup{(R)}}
\addcontentsline{toc}{section}{5. Subgroups of Type (R)}

We are especially interested in reductive groups, but some of the
results that we shall establish hold more generally for a wider
class of groups: groups of type \textup{(RR)}.

\SGARefTarget{sga3:xxii:section:5.1}\subsection*{\SGAInlineRefTarget{sga3:xxii:subsection:5-1}5.1. Groups of
Type \textup{(RR)}}
\addcontentsline{toc}{subsection}{\SGARefLink{sga3:xxii:subsection:5-1}{5.1}. Groups of Type (RR)}

\SGARefTarget{sga3:xxii:def:5.1.1}\paragraph{\SGARefLink{sga3:xxii:def:5.1.1}{Definition 5.1.1}.}
Let \(S\) be a scheme and \(G\) an \(S\)-group scheme. We say that
\(G\) is \emph{of type \textup{(RR)}} if it satisfies the following
conditions:

\begin{enumerate}
\item[\textup{\SGAInlineRefTarget{sga3:xxii:def:5.1.1:item:i}(i)}]
\(G\) is smooth and of finite presentation over \(S\), with
connected fibers.

\item[\textup{\SGAInlineRefTarget{sga3:xxii:def:5.1.1:item:ii}(ii)}]
\(G\) locally possesses maximal tori for the fpqc topology.

\item[\textup{\SGAInlineRefTarget{sga3:xxii:def:5.1.1:item:iii}(iii)}]
For every \(s\in S\), every maximal torus \(T\) of
\(G_{\bar s}\), and every root \(\alpha\) of \(G_{\bar s}\)
relative to \(T_{\bar s}\) \textup{(\SGARefLink{sga3:xix:paragraph:1-10}{Exposé XIX, 1.10})},
\[
  \Lie(G_{\bar s})^\alpha
\]
has rank \(1\) \textup{(as a vector space over
\(\overline{\kappa(s)}\))}.

\item[\textup{\SGAInlineRefTarget{sga3:xxii:def:5.1.1:item:iv}(iv)}]
For every \(s\in S\) and every maximal torus \(T\) of
\(G_{\bar s}\), let \(R\) denote the set of roots of
\(G_{\bar s}\) relative to \(T\), and let \(\Gamma_0(R)\) be the
subgroup of
\[
  \operatorname{Hom}_{\bar s\text{-gr.}}
    (T,\mathbb G_{m,\bar s})
\]
generated by \(R\). Then the
content
of every root \(\alpha\in R\) in the free abelian group
\(\Gamma_0(R)\) \textup{(and hence a positive integer)} is
invertible on \(S\).
\end{enumerate}

\SGARefTarget{sga3:xxii:reminder:5.1.1.1}
\paragraph{Reminder \SGARefLink{sga3:xxii:reminder:5.1.1.1}{5.1.1.1}.}
\footnote{Editor's note: This reminder was added.}
Let us recall that, if \(G\) is a smooth connected algebraic group
over an algebraically closed field \(k\), a \emph{Cartan subgroup}
of \(G\) is the centralizer of a maximal torus of \(G\)
\textup{(\SGARefLink{sga3:xii:1.0}{Exposé XII, 1.0})}, and such a subgroup is smooth and
connected. For this, as well as for other characterizations of
Cartan subgroups, see \emph{Bible}, \S7.1, Theorem 1 in the case
where \(G\) is affine, and \SGARefLink{sga3:xii:6.6}{Exposé XII, Theorem 6.6} in the general
case.

If \(S\) is an arbitrary scheme and \(G\) is a smooth \(S\)-group
of finite type, a \emph{Cartan subgroup} of \(G\) is a smooth
\(S\)-subgroup \(C\) of \(G\) such that, for every \(s\in S\),
\(C_{\bar s}\) is a Cartan subgroup of \(G_{\bar s}\)
\textup{(\SGARefLink{sga3:xii:3.1}{Exposé XII, Definition 3.1})}.

\SGARefTarget{sga3:xxii:rem:5.1.2}\paragraph{\SGARefLink{sga3:xxii:rem:5.1.2}{Remark 5.1.2}.}
\textup{\SGAInlineRefTarget{sga3:xxii:rem:5.1.2:item:a}(a)} By \SGARefLink{sga3:xii:7.1}{Exposé XII, 7.1} \textup{(where the hypothesis that
\(G\) is separated is satisfied here because \(G\) has connected
fibers; cf.\ Exposé VI\textsubscript{B}, \SGARefLink{cor:sga3-VIB-5.5}{5.5})}, conditions
\textup{\SGARefLink{sga3:xxii:def:5.1.1:item:i}{(i)}} and \textup{\SGARefLink{sga3:xxii:def:5.1.1:item:ii}{(ii)}} imply that \(G\) locally possesses,
for the étale topology, maximal tori \textup{(respectively, Cartan
subgroups)}, which are locally conjugate for the étale topology.

\textup{\SGARefLink{sga3:xxii:def:1.5:item:b}{(b)}} The Cartan subgroups of \(G\) have connected fibers
\textup{(\SGARefLink{sga3:xii:6.6}{Exposé XII, 6.6})}.

\textup{\SGARefLink{sga3:xii:6.6:item:c}{(c)}} If \(G\) is affine over \(S\), conditions
\textup{(i)} and \textup{\SGARefLink{sga3:xxii:def:5.1.1:item:ii}{(ii)}} are respectively equivalent to:
\begin{enumerate}
\item[\textup{(i$'$)}]
\(G\) is smooth over \(S\) and has connected fibers;

\item[\textup{(ii$'$)}]
the reductive rank of the fibers of \(G\) is locally constant
\textup{(\SGARefLink{sga3:xii:1.7}{Exposé XII, 1.7})}.
\end{enumerate}

\textup{\SGAInlineRefTarget{sga3:xxii:rem:5.1.2:item:d}(d)} By \SGARefLink{sga3:xii:8.8}{Exposé XII, 8.8} \textup{\SGARefLink{sga3:xii:8.8:item:c:02}{(c)}} and \textup{\SGARefLink{sga3:xxii:rem:5.1.2:item:d}{(d)}},
\(G\) has a reductive center \(Z\), and for every \(s\in S\), in
the notation of \textup{\SGARefLink{sga3:xxii:def:5.1.1:item:iv}{(iv)}},
\[
  Z_{\bar s}
  =
  \bigcap_{\alpha\in R}\operatorname{Ker}(\alpha),
\]
whence
\[
  \operatorname{Hom}_{\bar s\text{-gr.}}
    \bigl((T/Z)_{\bar s},\mathbb G_{m,\bar s}\bigr)
  \simeq
  \Gamma_0(R).
\]

\textup{\SGAInlineRefTarget{sga3:xxii:rem:5.1.2:item:e}(e)} Condition \textup{\SGARefLink{sga3:xxii:def:5.1.1:item:iv}{(iv)}} holds in particular in the
following two cases:
\begin{enumerate}
\item[\textup{(1)}]
\(S\) has characteristic \(0\);

\item[\textup{(2)}]
every root is an indivisible element of the group generated by the
roots.
\end{enumerate}

\textup{(f)} The property of being of type \textup{(RR)} is stable
under base change and is local for the fpqc topology.

Remarks \textup{(c)} and \textup{\SGARefLink{sga3:xxii:rem:5.1.2:item:e}{(e)}} immediately give the
following proposition.

\SGARefTarget{sga3:xxii:prop:5.1.3}
\paragraph{\SGARefLink{sga3:xxii:prop:5.1.3}{Proposition 5.1.3}.}
\textit{A reductive group is of type \textup{(RR)}.}

\SGARefTarget{sga3:xxii:prop:5.1.4}
\paragraph{\SGARefLink{sga3:xxii:prop:5.1.4}{Proposition 5.1.4}.}
\textit{Let \(S\) be a scheme, let \(G\) be an \(S\)-group of type
\textup{(RR)}, and let \(Q\) be a central subgroup of \(G\), of
finite presentation over \(S\), such that the quotient \(G/Q\) is
representable \textup{(for example, \(G\) affine over \(S\) and
\(Q\) of multiplicative type \textup{(\SGARefLink{sga3:ix:proposition:2:3}{Exposé IX, 2.3})}, or \(S\)
artinian \textup{(Exposé VI\textsubscript{A}, \SGARefLink{thm:sga3-VIA-3.3.2}{3.3.2})})}. Then
\(G/Q\) is an \(S\)-group of type \textup{(RR)}.}

Indeed, \(G/Q\) is smooth over \(S\)
\textup{(Exposé VI\textsubscript{B}, \SGARefLink{prop:sga3-VIB-9.2}{9.2})}, of finite presentation
over \(S\) \textup{(\SGARefLink{sga3:v:proposition:9:1}{Exposé V, 9.1})}, and has connected fibers, so
condition \textup{\SGARefLink{concl:sga3-VIB-9.2-i}{(i)}} holds. On the other hand, condition
\textup{\SGARefLink{concl:sga3-VIB-9.2-ii}{(ii)}} follows from \SGARefLink{sga3:xii:7.6}{Exposé XII, 7.6}; it remains to verify
conditions \textup{\SGARefLink{sga3:xxii:def:5.1.1:item:iii}{(iii)}} and \textup{\SGARefLink{sga3:xxii:def:5.1.1:item:iv}{(iv)}}.

Put \(G'=G/Q\), let
\[
  u:G\longrightarrow G'
\]
be the canonical morphism, and let \(T'=u(T)\) be the maximal torus
of \(G'\) that is the image of \(T\)
\textup{(cf.\ \SGARefLink{sga3:xii:7.1}{Exposé XII, 7.1} \textup{\SGARefLink{sga3:xii:7.1:item:e}{(e)}})}. For each
\(\alpha\in R\), again denote by \(\alpha\) the character of
\(T'\) induced by \(\alpha\); indeed,
\[
  Q\cap T
  \subset
  \bigcap_{\alpha\in R}\operatorname{Ker}(\alpha)
\]
by \SGARefLink{sga3:xxii:rem:5.1.2}{5.1.2} \textup{\SGARefLink{sga3:xxii:rem:5.1.2:item:d}{(d)}}. We first prove:

\SGARefTarget{sga3:xxii:lemma:5.1.5}
\paragraph{\SGARefLink{sga3:xxii:lemma:5.1.5}{Lemma 5.1.5}.}
\textit{Under the hypotheses of \SGARefLink{sga3:xxii:prop:5.1.4}{5.1.4}, let \(T=D_S(M)\) be a
trivialized maximal torus of \(G\), and suppose that the
decomposition of \(\mathfrak g=\Lie(G)\) under
\(\operatorname{Ad}(T)\) has the form}
\[
  \mathfrak g
  =
  \mathfrak g^0
  \oplus
  \coprod_{\alpha\in R}\mathfrak g^\alpha,
  \qquad
  R\subset M-\{0\},
\]
\textit{where, for every \(s\in S\),
\(\mathfrak g^\alpha(s)\ne0\) for \(\alpha\in R\)
\textup{(thus \(\mathfrak g^\alpha\) is an invertible
\(\mathcal O_S\)-module for every \(\alpha\in R\), and \(R\) is
the set of roots of \(G_{\bar s}\) relative to \(T_{\bar s}\), for
every \(s\in S\))}. Then the Lie algebra \(\mathfrak g'\) of \(G'\)
decomposes under \(\operatorname{Ad}(T')\) as follows:}
\[
  \mathfrak g'
  =
  \mathfrak g'^{\,0}
  \oplus
  \coprod_{\alpha\in R}\mathfrak g'^{\,\alpha},
\]
\textit{and \(\Lie(u)\) induces an isomorphism from
\(\mathfrak g^\alpha\) onto \(\mathfrak g'^{\,\alpha}\).}

Indeed, let
\[
  p=\Lie(u):\mathfrak g\longrightarrow\mathfrak g'.
\]
One immediately has
\(p(\mathfrak g^\alpha)\subset\mathfrak g'^{\,\alpha}\) for every
\(\alpha\in R\), and
\(p(\mathfrak g^0)\subset\mathfrak g'^{\,0}\). Since
\[
  \operatorname{Ker}(p)
  =
  \Lie(Q)
  \subset
  \Lie\bigl(\operatorname{Centr}_G(T)\bigr)
  =
  \mathfrak g^0,
\]
\(p\) induces a monomorphism from \(\mathfrak g^\alpha\) into
\(\mathfrak g'^{\,\alpha}\), for every \(\alpha\in R\).

To prove the lemma, it is enough to do so when \(S\) is the
spectrum of an algebraically closed field; in view of the preceding
remarks, it is then enough to prove that
\[
  \operatorname{rk}(\mathfrak g')
  =
  \operatorname{rk}(\mathfrak g'^{\,0})
  +
  \operatorname{Card}(R).
\]
Set
\[
  C=\operatorname{Centr}_G(T),
  \qquad
  C'=\operatorname{Centr}_{G'}(T').
\]
By \SGARefLink{sga3:xii:7.1}{Exposé XII, 7.1} \textup{\SGARefLink{sga3:xii:7.1:item:e}{(e)}}, \(u\) induces a faithfully flat
morphism \(C\to C'\) with kernel \(Q\). Hence
\[
  \dim C'+\dim Q=\dim C.
\]
But \(G\), \(G'\), \(C\), and \(C'\) are smooth, and therefore
\[
\begin{aligned}
  \dim G
  &=
  \operatorname{rk}(\mathfrak g)
   =
  \operatorname{rk}(\mathfrak g^0)+\operatorname{Card}(R)
   =
  \dim C+\operatorname{Card}(R)
  \\
  &=
  \dim Q+\dim C'+\operatorname{Card}(R)
   =
  \dim Q+\operatorname{rk}(\mathfrak g'^{\,0})
  +\operatorname{Card}(R),
  \\
  \operatorname{rk}(\mathfrak g')
  &=
  \dim G'
   =
  \dim G-\dim Q.
\end{aligned}
\]
It follows that
\[
  \operatorname{rk}(\mathfrak g')
  =
  \operatorname{rk}(\mathfrak g'^{\,0})
  +
  \operatorname{Card}(R),
\]
which is the desired relation.

We return to the proof of \SGARefLink{sga3:xxii:prop:5.1.4}{5.1.4}. We may suppose that \(S\) is the
spectrum of an algebraically closed field. Let \(T\) be a maximal
torus of \(G\). Applying \SGARefLink{sga3:xxii:lemma:5.1.5}{5.1.5}, we immediately obtain
\textup{\SGARefLink{sga3:xxii:def:5.1.1:item:iii}{(iii)}} and \textup{\SGARefLink{sga3:xxii:def:5.1.1:item:iv}{(iv)}} for \(G/Q\).

To use the preceding proposition, we introduce a definition.

\SGARefTarget{sga3:xxii:def:5.1.6}\paragraph{\SGARefLink{sga3:xxii:def:5.1.6}{Definition 5.1.6}.}
We say that the \(S\)-group scheme \(G\) is \emph{of type
\textup{(RA)}} if it is of type \textup{(RR)} and also satisfies
condition \textup{(iv$'$)} below, which is stronger than
\textup{(iv)}:
\begin{enumerate}
\item[\textup{(iv$'$)}]
For every \(s\in S\) and every maximal torus \(T\) of
\(G_{\bar s}\), the content of every root of \(G_{\bar s}\)
relative to \(T\), in
\[
  \operatorname{Hom}_{\bar s\text{-gr.}}
    (T,\mathbb G_{m,\bar s}),
\]
is invertible on \(S\).
\end{enumerate}

\SGARefTarget{sga3:xxii:rem:5.1.7}\paragraph{\SGARefLink{sga3:xxii:rem:5.1.7}{Remarks 5.1.7}.}
\textup{\SGARefLink{sga3:xxii:rem:5.1.2:item:a}{(a)}} An adjoint reductive \(S\)-group is of type
\textup{(RA)}.

\textup{\SGARefLink{sga3:xxii:def:1.5:item:b}{(b)}} If \(S\) has characteristic \(0\), every group of
type \textup{(RR)} is of type \textup{(RA)}.

\textup{(c)} The property of being of type \textup{(RA)} is stable
under base change and is local for the fpqc topology.

Remark \textup{\SGARefLink{sga3:xxii:rem:5.1.2:item:a}{(a)}} above admits the following generalization.

\SGARefTarget{sga3:xxii:prop:5.1.8}\paragraph{\SGARefLink{sga3:xxii:prop:5.1.8}{Proposition 5.1.8}.}
\textit{Let \(S\) be a scheme, let \(G\) be an \(S\)-group of type
\textup{(RR)}, and let \(Z\) be its reductive center. Suppose that
the quotient \(G/Z\) is representable \textup{(for example, if
\(G\) is affine over \(S\), or if \(S\) is artinian)}. Then \(G/Z\)
is of type \textup{(RA)}.}

Indeed, this follows immediately from \SGARefLink{sga3:xxii:prop:5.1.4}{5.1.4}, \SGARefLink{sga3:xxii:lemma:5.1.5}{5.1.5}, and \SGARefLink{sga3:xxii:rem:5.1.2}{5.1.2}
\textup{\SGARefLink{sga3:xxii:rem:5.1.2:item:d}{(d)}}.

\SGARefTarget{sga3:xxii:rem:5.1.9}\paragraph{\SGARefLink{sga3:xxii:rem:5.1.9}{Remark 5.1.9}.}
If \(G\) is of type \textup{(RR)} and \(T\) is a maximal torus of
\(G\), \SGARefLink{sga3:xix:proposition:6-3}{Exposé XIX, 6.3} applies. In particular,
\(W_G(T)\) is étale, quasi-finite, and separated over \(S\).

\SGARefTarget{sga3:xxii:section:5.2}
\subsection*{\SGAInlineRefTarget{sga3:xxii:subsection:5-2}5.2. Subgroups of Type \textup{(R)}}
\addcontentsline{toc}{subsection}{\SGARefLink{sga3:xxii:subsection:5-2}{5.2}. Subgroups of Type (R)}

\SGARefTarget{sga3:xxii:def:5.2.1}\paragraph{\SGARefLink{sga3:xxii:def:5.2.1}{Definition 5.2.1}.}
Let \(S\) be a scheme, let \(G\) be a smooth \(S\)-group scheme of
finite presentation with connected fibers, and let \(H\) be a group subscheme of \(G\). We say that
\(H\) is \emph{of type \textup{(R)}} if:
\begin{enumerate}
\item[\textup{\SGAInlineRefTarget{sga3:xxii:def:5.2.1:item:i}(i)}]
\(H\) is smooth and of finite presentation over \(S\), and has
connected fibers;

\item[\textup{\SGAInlineRefTarget{sga3:xxii:def:5.2.1:item:ii}(ii)}]
for every \(s\in S\), \(H_{\bar s}\) contains a Cartan subgroup of
\(G_{\bar s}\).
\end{enumerate}
This notion is stable under base change and is local for the fpqc
topology.

\SGARefTarget{sga3:xxii:recall:5.2.2}\paragraph{Recall \SGARefLink{sga3:xxii:recall:5.2.2}{5.2.2}.}
\textup{(cf.\ \SGARefLink{sga3:xii:7.9}{Exposé XII, 7.9})} Under the preceding conditions:
\begin{enumerate}
\item[\textup{\SGAInlineRefTarget{sga3:xxii:def:5.2.1:item:a}(a)}]
\[
  H=\underline{\operatorname{Norm}}_G(H)^0.
\]

\item[\textup{\SGAInlineRefTarget{sga3:xxii:def:5.2.1:item:b}(b)}]
If \(G\) locally possesses Cartan subgroups (respectively maximal
tori) for the étale topology, then so does \(H\), and the Cartan
subgroups (respectively maximal tori) of \(H\) are Cartan subgroups
(respectively maximal tori) of \(G\).
\end{enumerate}

\SGARefTarget{sga3:xxii:examples:5.2.3}\paragraph{\SGARefLink{sga3:xxii:examples:5.2.3}{Examples 5.2.3}.}
\textup{\SGARefLink{sga3:xxii:def:5.2.1:item:a}{(a)}} \emph{Borel subgroups.} A Borel subgroup of \(G\) is a
subgroup \(H\) of type \textup{(R)} whose geometric fibers are
Borel subgroups of the corresponding fibers of
\(G\).

\textup{(b)} \emph{Parabolic subgroups.} A parabolic subgroup of
\(G\) is a subgroup of type \textup{(R)} whose geometric fibers
contain Borel subgroups.

Other examples are given by the following propositions.

\SGARefTarget{sga3:xxii:prop:5.2.4}\paragraph{\SGARefLink{sga3:xxii:prop:5.2.4}{Proposition 5.2.4}.}
\textit{Let \(G\) be as in \SGARefLink{sga3:xxii:def:5.2.1}{5.2.1}, let \(K\subset H\) be two group
subschemes of \(G\), and suppose that \(H\) is of type
\textup{(R)}. Then \(K\) is a subgroup of type \textup{(R)} of
\(H\) if and only if it is a subgroup of type \textup{(R)} of
\(G\).}

Indeed,\footnote{Editor's note: We have expanded the argument that
follows.} let \(s\in S\). Since \(H\) is of type \textup{(R)},
every maximal torus of \(H_{\bar s}\) is a maximal torus of
\(G_{\bar s}\), and the same therefore holds for Cartan subgroups.

\SGARefTarget{sga3:xxii:prop:5.2.5}\paragraph{\SGARefLink{sga3:xxii:prop:5.2.5}{Proposition 5.2.5}.}
\textit{Let \(G\) be as in \SGARefLink{sga3:xxii:def:5.2.1}{5.2.1}, let \(T\) be a maximal torus of
\(G\), let \(Q\) be a subtorus of \(T\), and put
\(Z=\underline{\operatorname{Centr}}_G(Q)\). If \(H\) is a
subgroup of type \textup{(R)} of \(G\) containing \(T\), then
\(H\cap Z\) is a subgroup of type \textup{(R)} of \(Z\).}

First recall that \(Z\) is a closed group subscheme of \(G\), of
finite presentation (\SGARefLink{sga3:xi:corollary:6:11}{Exposé XI, 6.11}), with connected fibers
(\SGARefLink{sga3:xii:6.6}{Exposé XII, 6.6}), and smooth over \(S\) (\SGARefLink{sga3:xi:corollary:2:4}{Exposé XI, 2.4}), so it
satisfies the conditions imposed in \SGARefLink{sga3:xxii:def:5.2.1}{Definition 5.2.1}. Likewise,
\(H\cap Z\) is of finite presentation, smooth, and has connected
fibers (since
\(H\cap Z=\underline{\operatorname{Centr}}_H(Q)\)); moreover,
\[
  H\cap Z\supset\underline{\operatorname{Centr}}_G(T).
\]

\SGARefTarget{sga3:xxii:prop:5.2.6}\paragraph{\SGARefLink{sga3:xxii:prop:5.2.6}{Proposition 5.2.6}.}
\textit{Let \(S\) be a scheme, let \(G\) be an \(S\)-group of type
\textup{(RR)} \textup{(respectively \textup{(RA)})}, and let \(H\)
be a subgroup of type \textup{(R)} of \(G\). Then \(H\) is an
\(S\)-group of type \textup{(RR)} \textup{(respectively
\textup{(RA)})}.}

Indeed, \textup{\SGARefLink{sga3:xxii:def:5.2.1:item:i}{(i)}} is clear, \textup{\SGARefLink{sga3:xxii:def:5.2.1:item:ii}{(ii)}} follows from \SGARefLink{sga3:xxii:recall:5.2.2}{5.2.2}
\textup{\SGARefLink{sga3:xxii:def:5.2.1:item:b}{(b)}}, and \textup{\SGARefLink{sga3:xxii:def:5.1.1:item:iii}{(iii)}} and \textup{\SGARefLink{sga3:xxii:def:5.1.1:item:iv}{(iv)}}
\textup{(respectively \textup{(iv$'$)})} are to be verified when
\(S\) is the spectrum of an algebraically closed field. Then \(H\)
contains a maximal torus \(T\) of \(G\), and therefore also
\[
  C=\underline{\operatorname{Centr}}_G(T),
\]
and the assertions to be proved follow immediately from:

\SGARefTarget{sga3:xxii:lemma:5.2.7}\paragraph{\SGARefLink{sga3:xxii:lemma:5.2.7}{Lemma 5.2.7}.}
\textit{Let \(S\) be a scheme, let \(G\) be an \(S\)-group of type
\textup{(RR)}, and let \(T\) be a maximal torus of \(G\) equipped
with a trivialization \(T\simeq D_S(M)\). Suppose that}
\[
  \mathfrak g
  =
  \mathfrak g^0
  \oplus
  \coprod_{\alpha\in R}\mathfrak g^\alpha
\]
\textit{\textup{(the \(\mathfrak g^\alpha\) are then invertible
\(\mathcal O_S\)-modules)}. Let \(H\) be a subgroup of type
\textup{(R)} containing
\(C=\underline{\operatorname{Centr}}_G(T)\)
\textup{(that is, containing \(T\))}. Then
\(\mathfrak h=\Lie(H/S)\) is locally on \(S\) of the form}
\[
  \mathfrak g^0
  +
  \coprod_{\alpha\in R'}\mathfrak g^\alpha
  =
  \mathfrak g_{R'}.
\]
\textit{More precisely, for each \(s\in S\), put}
\[
  R'(s)
  =
  \bigl\{
    \alpha\in R
    \mathrel{\big|}
    \mathfrak g^\alpha(s)\subset\mathfrak h(s)
  \bigr\}.
\]
\textit{Then \(R'(s)\) is a locally constant function of \(s\). If
\(U\) is an open subset of \(S\) on which \(R'(s)=R'\), then}
\[
  \mathfrak h_U
  =
  \mathfrak g_U^0
  \oplus
  \coprod_{\alpha\in R'}\mathfrak g_U^\alpha.
\]

Indeed, \(\mathfrak h\) is a submodule of \(\mathfrak g\), locally
a direct factor, containing \(\mathfrak g^0\), and stable under
\(T\).

\SGARefTarget{sga3:xxii:section:5.3}
\subsection*{\SGAInlineRefTarget{sga3:xxii:subsection:5-3}5.3. Strict Transporter of Two Subgroups of Type
\textup{(R)}. Applications}
\addcontentsline{toc}{subsection}{\SGARefLink{sga3:xxii:subsection:5-3}{5.3}. Strict Transporter of Two
Subgroups of Type (R). Applications}

\SGARefTarget{sga3:xxii:recall:5.3.0}\paragraph{Recall \SGARefLink{sga3:xxii:recall:5.3.0}{5.3.0}.}
\footnote{Editor's note: This reminder was added; it is used in the
proofs of \SGARefLink{sga3:xxii:prop:5.3.1}{5.3.1} and \SGARefLink{sga3:xxii:cor:5.3.4}{5.3.4}.} Let \(S\) be a scheme, let \(G\) be a
smooth \(S\)-group, put \(\mathfrak g=\Lie(G/S)\), and let
\(\mathfrak h\) be an \(\mathcal O_S\)-submodule of \(\mathfrak g\)
that is locally a direct factor. The \(\mathcal O_S\)-algebra
\[
  \mathscr A=\operatorname{Sym}(\omega^1_{G/S})
\]
is locally free; hence the \(S\)-scheme
\[
  \underline{\operatorname{Lie}}(G/S)
  =
  W(\mathfrak g)
  =
  \operatorname{Spec}(\mathscr A)
\]
is essentially free in the sense of \SGARefLink{sga3:viii:definition:6:1}{Exposé VIII, 6.1}. Since
\(W(\mathfrak h)\) is a closed subscheme of
\(\underline{\operatorname{Lie}}(G/S)\), of finite presentation
over \(\underline{\operatorname{Lie}}(G/S)\), it follows from
\SGARefLink{sga3:viii:examples:6:5}{Exposé VIII, 6.5} \textup{\SGARefLink{sga3:viii:examples:6:5:item:a:01}{(a)}} that
\[
  N=\underline{\operatorname{Norm}}_G(\mathfrak h)
\]
is represented by a closed group subscheme of \(G\), of finite
presentation over \(G\). (See also the additions \SGARefLink{sga3:xxii:par:6.2.3}{6.2.3} and \SGARefLink{sga3:xxii:prop:6.2.4}{6.2.4}
\textup{\SGARefLink{sga3:viii:examples:6:5:item:a:01}{(a)}} in Exposé VI\textsubscript{B}.) On the other hand, by
\SGARefLink{sga3:ii:theorem:5-3-1}{Exposé II, 5.3.1},
\[
  \underline{\operatorname{Lie}}(N/S)
  =
  \underline{\operatorname{Norm}}_
    {\underline{\operatorname{Lie}}(G/S)}(\mathfrak h).
\]

Finally, by Exposé VI\textsubscript{B}, \SGARefLink{thm:sga3-VIB-3.10}{3.10}, if \(N\) is smooth
over \(S\) at the points of the unit section, then the group
subfunctor \(N^0\) (defined in Exposé VI\textsubscript{B}, \SGARefLink{def:sga3-VIB-3.1}{3.1}) is
represented by an open group subscheme of \(N\), smooth over \(S\).

\SGARefTarget{sga3:xxii:prop:5.3.1}\paragraph{\SGARefLink{sga3:xxii:prop:5.3.1}{Proposition 5.3.1}.}
\textit{Let \(S\) be a scheme, let \(G\) be an \(S\)-group of type
\textup{(RA)} \textup{(\SGARefLink{sga3:xxii:def:5.1.6}{5.1.6})}, let \(H\) be a subgroup of type
\textup{(R)} of \(G\), and let
\(\mathfrak g\supset\mathfrak h\) be their Lie algebras. Then
\(\underline{\operatorname{Norm}}_G(\mathfrak h)\)
\textup{(which, by \SGARefLink{sga3:xxii:recall:5.3.0}{5.3.0}, is representable by a closed subscheme of
\(G\), of finite presentation over \(S\))} is smooth over \(S\) at
every point of the unit section, and}
\[
  \underline{\operatorname{Norm}}_G(\mathfrak h)^0=H.
\]

\footnote{Editor's note: We have expanded the original in the
argument that follows.} \emph{Proof.} Put
\[
  N=\underline{\operatorname{Norm}}_G(\mathfrak h),
  \qquad
  \mathfrak n=\Lie(N/S).
\]
We have \(H\subset N\), and by \SGARefLink{sga3:ii:theorem:5-3-1}{Exposé II, 5.3.1}, for every
\(s\in S\),
\[
  \mathfrak h(s)
  \subset
  \mathfrak n(s)
  =
  \operatorname{Norm}_{\mathfrak g(s)}(\mathfrak h(s)).
\]
Now, by \SGARefLink{sga3:xxii:lemma:5.3.2}{5.3.2} below,
\[
  \mathfrak h(s)
  =
  \operatorname{Norm}_{\mathfrak g(s)}(\mathfrak h(s)),
\]
and since \(H\) is smooth over \(S\),
\[
  \dim_{\kappa(s)}\mathfrak h(s)=\dim H_s
\]
\textup{(cf.\ \textup{[DG70]}, \S II.5, Theorem 2.1)}. Hence
\[
  \dim_{\kappa(s)}\mathfrak n(s)
  =
  \dim_{\kappa(s)}\mathfrak h(s)
  =
  \dim H_s
  \leq
  \dim N_s.
\]
It follows that
\[
  N_s^0=H_s^0=H_s,
\]
since \(H\) has connected fibers. Consequently the group subfunctor
\(N^0\) \textup{(defined in Exposé VI\textsubscript{B}, \SGARefLink{def:sga3-VIB-3.1}{3.1})} is
represented by the smooth \(S\)-group \(H\). This proves \SGARefLink{sga3:xxii:prop:5.3.1}{5.3.1},
modulo the following lemma.

\SGARefTarget{sga3:xxii:lemma:5.3.2}\paragraph{\SGARefLink{sga3:xxii:lemma:5.3.2}{Lemma 5.3.2}.}
\textit{Under the conditions of \SGARefLink{sga3:xxii:lemma:5.2.7}{5.2.7}, if \(G\) is of type
\textup{(RA)}, then, for every \(s\in S\),}
\[
  \operatorname{Norm}_{\mathfrak g(s)}(\mathfrak h(s))
  =
  \mathfrak h(s).
\]

Indeed, we are reduced to the case where \(S\) is the spectrum of a
field, and hence where
\[
  \mathfrak h=\mathfrak g_{R'}
\]
for some \(R'\subset R\). But we already have
\[
  \operatorname{Transp}_{\mathfrak g}
    (\mathfrak t,\mathfrak h)
  =
  \mathfrak h.
\]
Indeed, if \(H\in\mathfrak t\) and
\(X\in\mathfrak g^\alpha\), then
\[
  [H,X]=\bar\alpha(H)X,
\]
where \(\bar\alpha:\mathfrak t\to\mathcal O_S\) is the morphism
derived from \(\alpha\). Condition \textup{(iv$'$)} says precisely
that \(\bar\alpha\ne0\) for every \(\alpha\in R\).

\SGARefTarget{sga3:xxii:cor:5.3.3}\paragraph{\SGARefLink{sga3:xxii:cor:5.3.3}{Corollary 5.3.3}.}
\textit{Let \(S\) be a scheme, let \(G\) be an \(S\)-group of type
\textup{(RA)}, let \(H\) and \(H'\) be two subgroups of type
\textup{(R)} of \(G\), and let \(\mathfrak h\) and
\(\mathfrak h'\) be their Lie algebras. Then}
\[
  H=H'
  \quad\Longleftrightarrow\quad
  \mathfrak h=\mathfrak h'.
\]

\SGARefTarget{sga3:xxii:cor:5.3.4}\paragraph{\SGARefLink{sga3:xxii:cor:5.3.4}{Corollary 5.3.4}.}
\textit{Under the conditions of \SGARefLink{sga3:xxii:lemma:5.2.7}{5.2.7}, with \(G\) of type
\textup{(RA)}, the maps}
\[
  H\longmapsto\Lie(H/S),
  \qquad
  \mathfrak h\longmapsto
    \underline{\operatorname{Norm}}_G(\mathfrak h)^0
\]
\textit{give a bijective correspondence between the set of
subgroups of type \textup{(R)} of \(G\) containing \(T\), and the
set of Lie subalgebras of \(\mathfrak g\) containing
\(\mathfrak g^0\), stable under \(T\), and whose normalizer in
\(G\) is smooth over \(S\) at every point of the unit section.}

Indeed,\footnote{Editor's note: The proof that follows was added.}
let \(\mathfrak h\) be a Lie subalgebra of \(\mathfrak g\)
satisfying the properties above. By \SGARefLink{sga3:xxii:recall:5.3.0}{5.3.0},
\[
  H=\underline{\operatorname{Norm}}_G(\mathfrak h)^0
\]
is a smooth \(S\)-group scheme. Moreover, since
\[
  C=\underline{\operatorname{Centr}}_G(T)
\]
stabilizes every \(\mathfrak g^\alpha\) and has connected fibers
\textup{(\SGARefLink{sga3:xii:6.6}{Exposé XII, 6.6})}, we have \(C\subset H\). Therefore
\(H\) is a subgroup of \(G\) of type \textup{(R)}. By \SGARefLink{sga3:ii:theorem:5-3-1}{Exposé II,
5.3.1},
\[
  \underline{\operatorname{Lie}}(H)
  =
  \underline{\operatorname{Norm}}_{\mathfrak g}(\mathfrak h).
\]
Finally, by the proof of \SGARefLink{sga3:xxii:lemma:5.3.2}{5.3.2},
\[
  \underline{\operatorname{Norm}}_{\mathfrak g}(\mathfrak h)
  =
  \mathfrak h.
\]

\SGARefTarget{sga3:xxii:cor:5.3.5}\paragraph{\SGARefLink{sga3:xxii:cor:5.3.5}{Corollary 5.3.5}.}
\textit{Let \(S\) be a scheme, let \(G\) be an \(S\)-group of type
\textup{(RR)} \textup{(\SGARefLink{sga3:xxii:def:5.1.1}{5.1.1})}, let \(T\) be a maximal torus of
\(G\), and let \(H\) and \(H'\) be two subgroups of type
\textup{(R)} of \(G\), both containing \(T\). Then}
\[
  H=H'
  \quad\Longleftrightarrow\quad
  \mathfrak h=\mathfrak h'.
\]

By virtue of the finite-presentation hypotheses, we reduce as usual
\textup{(cf.\ EGA IV\textsubscript{3}, \S8, and Exposé
VI\textsubscript{B}, \S10)} to the case where \(S\) is noetherian.
It is then enough to verify that
\(\mathfrak h=\mathfrak h'\) implies
\[
  H_{S'}=H'_{S'}
\]
for every \(S'\) that is the spectrum of an artinian quotient of a
local ring of \(S\);\footnote{Editor's note: Indeed, let \(g\in G\),
let \(s\) be its image in \(S\), let \(\mathfrak m\) be the maximal
ideal of \(\mathcal O_{G,g}\), let \(\mathfrak n\) be that of
\(\mathcal O_{S,s}\), and let \(I\) (respectively \(I'\)) be the
kernel of the morphism from \(\mathcal O_{G,g}\) to
\(\mathcal O_{H,g}\) (respectively \(\mathcal O_{H',g}\)), the
latter being the zero ring if \(g\notin H\) (respectively
\(g\notin H'\)). Since \(\mathcal O_{G,g}\) is noetherian, \(I\)
and \(I'\) are closed for the \(\mathfrak m\)-adic topology, hence
\emph{a fortiori} for the \(\mathfrak n\)-adic topology. It is
therefore enough to prove
\[
  I+\mathfrak n^n\mathcal O_{G,g}
  =
  I'+\mathfrak n^n\mathcal O_{G,g}
\]
for every \(n\in\mathbb N\).} we are therefore reduced to the case
where \(S\) is artinian, in which case we may apply \SGARefLink{sga3:xxii:prop:5.1.8}{5.1.8}.

Let
\[
  u:G\longrightarrow G'=G/Z
\]
be the canonical morphism, and let \(T'=T/Z\) be the maximal torus
of \(G'\) corresponding to \(T\). By \SGARefLink{sga3:xii:7.12}{Exposé XII, 7.12}, there exist
subgroups \(H_1\) and \(H'_1\) of type \textup{(R)} of \(G'\),
containing \(T'\), such that
\[
  H=u^{-1}(H_1),
  \qquad
  H'=u^{-1}(H'_1).
\]
It is enough to prove that \(H_1=H'_1\). But by \SGARefLink{sga3:xxii:lemma:5.2.7}{5.2.7} and \SGARefLink{sga3:xxii:lemma:5.1.5}{5.1.5},
\[
  \Lie(H_1)=\Lie(H'_1),
\]
and we are reduced to \SGARefLink{sga3:xxii:cor:5.3.3}{5.3.3}.

\SGARefTarget{sga3:xxii:rem:5.3.6}\paragraph{\SGARefLink{sga3:xxii:rem:5.3.6}{Remark 5.3.6}.}
The fact that \(H\) and \(H'\) contain the same maximal torus is
essential for the validity of \SGARefLink{sga3:xxii:cor:5.3.5}{5.3.5} when \(G\) is not of type
\textup{(RA)}. For example, consider the maximal tori of
\(\operatorname{SL}_{2,k}\), where \(k\) has characteristic
\(2\).

\SGARefTarget{sga3:xxii:cor:5.3.7}\paragraph{\SGARefLink{sga3:xxii:cor:5.3.7}{Corollary 5.3.7}.}
\textit{Let \(S\) be a scheme, let \(G\) be an \(S\)-group of type
\textup{(RR)}, let \(T\) be a maximal torus of \(G\), and let \(H\)
and \(H'\) be two subgroups of type \textup{(R)} of \(G\), both
containing \(T\). The set \(U\) of \(s\in S\) such that
\(H_s=H'_s\) is open and closed in \(S\), and}
\[
  H_U=H'_U.
\]

Indeed, this follows immediately from \SGARefLink{sga3:xxii:cor:5.3.5}{5.3.5} and \SGARefLink{sga3:xxii:lemma:5.2.7}{5.2.7}.

\SGARefTarget{sga3:xxii:cor:5.3.8}\paragraph{\SGARefLink{sga3:xxii:cor:5.3.8}{Corollary 5.3.8}.}
\textit{The ``functor of subgroups of type \textup{(R)} containing
\(T\)'', where \(T\) is a given maximal torus in a group \(G\) of
type \textup{(RR)}, is formally unramified \textup{(\SGARefLink{sga3:xi:definition:1:1}{Exposé XI,
1.1})}.}

\SGARefTarget{sga3:xxii:thm:5.3.9}\paragraph{\SGARefLink{sga3:xxii:thm:5.3.9}{Theorem 5.3.9}.}
\textit{Let \(G\) be an \(S\)-group of type \textup{(RR)}
\textup{(\SGARefLink{sga3:xxii:def:5.1.1}{5.1.1})}, and let \(H\) and \(H'\) be two subgroups of type
\textup{(R)} \textup{(\SGARefLink{sga3:xxii:def:5.2.1}{5.2.1})}. Let
\(\underline{\operatorname{Transt}}_G(H,H')\) be the strict
transporter of \(H\) into \(H'\), defined by}
\[
  \underline{\operatorname{Transt}}_G(H,H')(S')
  =
  \bigl\{
    g\in G(S')
    \mathrel{\big|}
    \operatorname{int}(g)H_{S'}=H'_{S'}
  \bigr\}.
\]
\textit{Then \(\underline{\operatorname{Transt}}_G(H,H')\) is
representable by a closed subscheme of \(G\), smooth and of finite
presentation over \(S\).}

The fact that
\(\underline{\operatorname{Transt}}_G(H,H')\) is representable by
a closed subscheme of \(G\), of finite presentation over \(S\),
follows from \SGARefLink{sga3:xi:corollary:6:11}{Exposé XI, 6.11} \textup{\SGARefLink{sga3:xi:corollary:6:11:item:a}{(a)}}. To prove that it is
smooth over \(S\), it is necessary to prove that if \(S\) is affine,
if \(S_0\) is the closed subscheme defined by a nilpotent ideal
\(J\), and if
\[
  g_0\in G(S_0),
  \qquad
  \operatorname{int}(g_0)H_0=H'_0,
\]
then there exists \(g\in G(S)\), lifting \(g_0\), such that
\[
  \operatorname{int}(g)H=H'.
\]
Since smoothness is local for the étale topology, we may suppose
that \(H\) contains a maximal torus \(T\) of \(G\).

Then \(T_0\) is a maximal torus of \(H_0\), and hence
\(\operatorname{int}(g_0)T_0\) is a maximal torus of \(H'_0\).
By \SGARefLink{sga3:ix:theorem:3:6:bis}{Exposé IX, 3.6 bis}, there exists a torus \(T'\) of \(H'\) such
that
\[
  T'_0=\operatorname{int}(g_0)T_0;
\]
by \SGARefLink{sga3:ix:corollary:3:3:bis}{Exposé IX, 3.3 bis}, there therefore exists \(g\in G(S)\), lifting
\(g_0\), such that
\[
  \operatorname{int}(g)T=T'.
\]
After replacing \(H\) by \(\operatorname{int}(g)H\), we may
therefore suppose that \(H\) and \(H'\) contain the same maximal
torus \(T\) and that \(H_0=H'_0\). But then \(H=H'\) by \SGARefLink{sga3:xxii:cor:5.3.7}{5.3.7}.
\qed

\SGARefTarget{sga3:xxii:cor:5.3.10}\paragraph{\SGARefLink{sga3:xxii:cor:5.3.10}{Corollary 5.3.10}.}
\textit{Let \(G\) be an \(S\)-group of type \textup{(RR)}, and let
\(H\) be a subgroup of type \textup{(R)} of \(G\). Then
\(\underline{\operatorname{Norm}}_G(H)\) is representable by a
closed group subscheme of \(G\), of finite presentation and smooth
over \(S\).}

Using now the reasoning employed in Exposé XI to deduce \SGARefLink{sga3:xi:corollary:5:4:bis}{5.4 bis}
from \SGARefLink{sga3:xi:corollary:5:2:bis}{5.2 bis}, we obtain:

\SGARefTarget{sga3:xxii:cor:5.3.11}\paragraph{\SGARefLink{sga3:xxii:cor:5.3.11}{Corollary 5.3.11}.}
\textit{Under the hypotheses of \SGARefLink{sga3:xxii:thm:5.3.9}{5.3.9}, the following conditions are
equivalent:}
\begin{enumerate}
  \item[\textup{\SGAInlineRefTarget{sga3:xxii:cor:5.3.11:item:i}(i)}] \textit{\(H\) and \(H'\) are locally conjugate
  in \(G\) for the étale topology;}

  \item[\textup{(i bis)}] \textit{the same holds for the fpqc
  topology;}

  \item[\textup{(ii)}] \textit{for every \(s\in S\),
  \(H_{\bar s}\) and \(H'_{\bar s}\) are conjugate by an element of
  \(G(\bar s)\);}

  \item[\textup{(ii bis)}] \textit{the structural morphism}
  \[
    \underline{\operatorname{Transt}}_G(H,H')
      \longrightarrow S
  \]
  \textit{is surjective;}

  \item[\textup{\SGAInlineRefTarget{sga3:xxii:cor:5.3.11:item:iii}(iii)}]
  \textit{\(\underline{\operatorname{Transt}}_G(H,H')\) is a
  principal homogeneous bundle under the action of the smooth
  \(S\)-group scheme of finite presentation
  \(\underline{\operatorname{Norm}}_G(H)\).}
\end{enumerate}

We merely note that the nontrivial implication
\textup{\SGARefLink{sga3:xxii:cor:5.3.11:item:iii}{(iii)}}\(\Rightarrow\)\textup{\SGARefLink{sga3:xxii:cor:5.3.11:item:i}{(i)}} is Hensel's lemma.

Using now \textit{Bible}, \S6.4, Theorem 4
\textup{(= [Ch05], \S6.5, Theorem 5)} and \S9.3,
Theorem 1, we obtain from \SGARefLink{sga3:xxii:cor:5.3.10}{5.3.10} and \SGARefLink{sga3:xxii:cor:5.3.11}{5.3.11}:

\SGARefTarget{sga3:xxii:cor:5.3.12}\paragraph{\SGARefLink{sga3:xxii:cor:5.3.12}{Corollary 5.3.12}.}
\textit{Let \(G\) be an \(S\)-group of type \textup{(RR)}. The
Borel subgroups of \(G\) are closed in \(G\), equal to their
normalizers, and locally conjugate for the étale topology.}

\SGARefTarget{sga3:xxii:def:5.3.13}\paragraph{\SGARefLink{sga3:xxii:def:5.3.13}{Definition 5.3.13}.}
\textit{Let \(S\) be a scheme, and let \(G\) be a smooth
\(S\)-group scheme of finite presentation with connected
fibers. A \emph{Killing pair}
of \(G\) is a pair \(T\subset B\), where \(T\) is a maximal torus
of \(G\) and \(B\) is a Borel subgroup of \(G\) containing \(T\).}

Using now the conjugacy of maximal tori in \(B\)
\textup{(cf.\ \SGARefLink{sga3:xxii:rem:5.1.2}{5.1.2} \textup{\SGARefLink{sga3:xxii:rem:5.1.2:item:a}{(a)}} and \SGARefLink{sga3:xxii:prop:5.2.6}{5.2.6}, for example)}, we have:

\SGARefTarget{sga3:xxii:cor:5.3.14}\paragraph{\SGARefLink{sga3:xxii:cor:5.3.14}{Corollary 5.3.14}.}
\textit{Let \(G\) be an \(S\)-group of type \textup{(RR)}. The
Killing pairs of \(G\) are locally conjugate for the étale topology.}

\SGARefTarget{sga3:xxii:cor:5.3.15}\paragraph{\SGARefLink{sga3:xxii:cor:5.3.15}{Corollary 5.3.15}.}
\textit{Let \(G\) be an \(S\)-group of type \textup{(RR)}. Let
\(T\) be a maximal torus of \(G\), and let}
\[
  W_G(T)
  =
  \underline{\operatorname{Norm}}_G(T)/
  \underline{\operatorname{Centr}}_G(T)
\]
\textit{be the corresponding Weyl group \textup{(\SGARefLink{sga3:xix:proposition:6:3}{Exposé XIX,
6.3})}. The ``functor of Borel subgroups of \(G\) containing \(T\)''
is formally principal homogeneous under \(W_G(T)\).}

This follows immediately from \SGARefLink{sga3:xxii:cor:5.3.14}{5.3.14} and from the fact that, if
\(B\) is a Borel subgroup of \(G\) containing \(T\), then
\[
  \underline{\operatorname{Norm}}_G(T)\cap B
  =
  \underline{\operatorname{Centr}}_G(T),
\]
cf.\ \SGARefLink{sga3:xiv:4.4}{Exposé XIV, 4.4}.

\SGARefTarget{sga3:xxii:prop:5.3.16}\paragraph{\SGARefLink{sga3:xxii:prop:5.3.16}{Proposition 5.3.16}.}
\textit{Let \(G\) be an \(S\)-group of type \textup{(RR)}, let
\(H\) be a subgroup of type \textup{(R)}, and let
\(N=\underline{\operatorname{Norm}}_G(H)\) be its normalizer
\textup{(\SGARefLink{sga3:xxii:cor:5.3.10}{5.3.10})}. Let \(T\) be a maximal torus of \(H\), and let
\(W_H(T)\) and \(W_N(T)\) be the corresponding Weyl groups
\textup{(étale, quasi-finite, and separated by \SGARefLink{sga3:xix:proposition:6:3}{Exposé XIX, 6.3})}.
There is the following exact sequence of sheaves \textup{(for the
étale topology)}, whose morphisms are induced by}
\[
  \underline{\operatorname{Norm}}_H(T)
  \longrightarrow
  \underline{\operatorname{Norm}}_N(T)
  \longrightarrow
  N/H:
\]
\[
  1
  \longrightarrow W_H(T)
  \longrightarrow W_N(T)
  \longrightarrow N/H
  \longrightarrow 1.
\]

The only nontrivial point is that the last arrow is an epimorphism.
Let \(n\in N(S')\), where \(S'\to S\). The two maximal tori \(T\)
and \(\operatorname{int}(n)T\) of \(H\) are locally conjugate in
\(H\) for the étale topology. Thus there is a covering family
\(\{S'_i\to S'\}\) and, for every \(i\), an
\(h_i\in H(S'_i)\) such that
\[
  \operatorname{int}(h_i)T=\operatorname{int}(n)T.
\]
Consequently
\[
  nh_i^{-1}\in \underline{\operatorname{Norm}}_N(T),
\]
which gives the desired result.

\SGARefTarget{sga3:xxii:rem:5.3.17}\paragraph{\SGARefLink{sga3:xxii:rem:5.3.17}{Remark 5.3.17}.}
The group \(W_N(T)\) may be described as follows. Suppose that we
have reduced to the situation of \SGARefLink{sga3:xxii:lemma:5.2.7}{5.2.7}, with
\(\mathfrak h=\mathfrak g_{R'}\). Then \(W_N(T)\) is equal to
\(\underline{\operatorname{Norm}}_W(R')\), the sheaf of sections
of \(W=W_G(T)\) which, acting on \(R\), normalize \(R'\). Indeed,
by \SGARefLink{sga3:xxii:cor:5.3.5}{5.3.5},
\[
  \underline{\operatorname{Norm}}_N(T)
  =
  \underline{\operatorname{Norm}}_G(H)
    \cap \underline{\operatorname{Norm}}_G(T)
  =
  \underline{\operatorname{Norm}}_G(\mathfrak h)
    \cap \underline{\operatorname{Norm}}_G(T).
\]

\SGARefTarget{sga3:xxii:cor:5.3.18}\paragraph{\SGARefLink{sga3:xxii:cor:5.3.18}{Corollary 5.3.18}.}
\textit{Let \(G\) be an \(S\)-group of type \textup{(RR)}, and let
\(H\) be a subgroup of type \textup{(R)}. Suppose that ``the Weyl
groups of \(G\) are finite'', that is, that for every \(S'\to S\)
and every maximal torus \(T\) of \(G_{S'}\), the étale
\(S'\)-scheme}
\[
  \underline{\operatorname{Norm}}_{G_{S'}}(T)/
  \underline{\operatorname{Centr}}_{G_{S'}}(T)
\]
\textit{is finite \textup{(cf.\ \SGARefLink{sga3:xix:proposition:6:3}{Exposé XIX, 6.3} \textup{\SGARefLink{sga3:xxii:cor:5.3.18:item:iii}{(iii)}})}.
Then the following conditions are equivalent:}
\begin{enumerate}
  \item[\textup{\SGAInlineRefTarget{sga3:xxii:cor:5.3.18:item:i}(i)}] \textit{\(H\) is closed in \(G\);}

  \item[\textup{\SGAInlineRefTarget{sga3:xxii:cor:5.3.18:item:ii}(ii)}]
  \textit{\(\underline{\operatorname{Norm}}_G(H)/H\) is
  representable by a finite étale \(S\)-scheme;}

  \item[\textup{\SGAInlineRefTarget{sga3:xxii:cor:5.3.18:item:iii}(iii)}]
  \textit{``the Weyl groups of \(H\) are finite''.}
\end{enumerate}

Indeed, we may suppose that \(H\) has a maximal torus \(T\). By
\SGARefLink{sga3:xxii:cor:5.3.10}{5.3.10}, \(N=\underline{\operatorname{Norm}}_G(H)\) is closed in
\(G\); therefore \(W_N(T)\) is closed in \(W_G(T)\), and hence is
finite over \(S\). Clearly \textup{\SGARefLink{sga3:xxii:cor:5.3.18:item:i}{(i)}} implies \textup{\SGARefLink{sga3:xxii:cor:5.3.18:item:iii}{(iii)}}, and
\textup{\SGARefLink{sga3:xxii:cor:5.3.18:item:iii}{(iii)}} implies \textup{\SGARefLink{sga3:xxii:cor:5.3.18:item:ii}{(ii)}} by the exact sequence of
\SGARefLink{sga3:xxii:prop:5.3.16}{5.3.16}. Finally, \textup{\SGARefLink{sga3:xxii:cor:5.3.18:item:ii}{(ii)}} implies \textup{\SGARefLink{sga3:xxii:cor:5.3.18:item:i}{(i)}}, since if
\(N/H\) is finite, it is separated, so \(H\) is closed in \(N\),
and hence in \(G\).

\SGARefTarget{sga3:xxii:rem:5.3.19}\paragraph{\SGARefLink{sga3:xxii:rem:5.3.19}{Remark 5.3.19}.}
When \(G\) is reductive, the preceding conditions on \(H\) seem
always to be satisfied. We prove them below in most cases.

\SGARefTarget{sga3:xxii:section:5.4}
\subsection*{\SGAInlineRefTarget{sga3:xxii:subsection:5-4}5.4. Subgroups of Type \textup{(R)} of a Split
Reductive Group \textup{(Generalities)}}
\addcontentsline{toc}{subsection}{\SGARefLink{sga3:xxii:subsection:5-4}{5.4}. Subgroups of Type (R) of a
Split Reductive Group (Generalities)}

\SGARefTarget{sga3:xxii:5.4.1}\paragraph{\SGARefLink{sga3:xxii:5.4.1}{5.4.1}.}
If \(H\) is a subgroup of type \textup{(R)} of the reductive group
\(G\), then \(H\) locally for the étale topology contains a maximal
torus of \(G\) \textup{(\SGARefLink{sga3:xxii:recall:5.2.2}{5.2.2})}. By \SGARefLink{sga3:xxii:cor:2.3}{2.3}, locally for the étale
topology we may suppose that \(G\) is split relative to this torus.
Thus let \((G,T,M,R)\) be a split \(S\)-group, and let \(H\) be a
subgroup of type \textup{(R)} of \(G\) containing \(T\). By \SGARefLink{sga3:xxii:cor:5.3.5}{5.3.5},
such a subgroup is characterized by its Lie algebra, which
\textup{(\SGARefLink{sga3:xxii:lemma:5.2.7}{5.2.7})} is locally on \(S\) of the form
\(\mathfrak g_{R'}\):
\[
  \mathfrak g_{R'}
  =
  \mathfrak t
  \oplus
  \prod_{\alpha\in R'}\mathfrak g^\alpha.
\]

\SGARefTarget{sga3:xxii:def:5.4.2}\paragraph{\SGARefLink{sga3:xxii:def:5.4.2}{Definition 5.4.2}.}
\textit{Let \((G,T,M,R)\) be a split \(S\)-group. A subset \(R'\)
of \(R\) is said to be of type \textup{(R)} if
\(\mathfrak g_{R'}\) is the Lie algebra of a subgroup of type
\textup{(R)} of \(G\) containing \(T\). This subgroup, uniquely
determined by \(R'\), is denoted \(H_{R'}\).}

\SGARefTarget{sga3:xxii:lem:5.4.3}\paragraph{\SGARefLink{sga3:xxii:lem:5.4.3}{Lemma 5.4.3}.}
\textit{Under the preceding conditions, the following equivalences
hold:}
\[
\begin{aligned}
  H\cap Z_\alpha=T
    &\quad\Longleftrightarrow\quad
      \alpha\notin R',\ -\alpha\notin R',\\
  H\supset U_\alpha
    &\quad\Longleftrightarrow\quad
      \alpha\in R',\\
  H\cap U_\alpha=e
    &\quad\Longleftrightarrow\quad
      \alpha\notin R',\\
  H\supset Z_\alpha
    &\quad\Longleftrightarrow\quad
      \alpha\in R',\ -\alpha\in R'.
\end{aligned}
\]

Indeed, \(H\cap Z_\alpha\) is a subgroup of type \textup{(R)} of
\(Z_\alpha\), by \SGARefLink{sga3:xxii:prop:5.2.5}{5.2.5}; but a subgroup of type \textup{(R)} of
\(Z_\alpha\) containing \(T\) is locally equal to one of the
following subgroups:
\[
  T,\qquad T\cdot U_\alpha,\qquad
  T\cdot U_{-\alpha},\qquad Z_\alpha,
\]
by \SGARefLink{sga3:xxii:cor:5.3.5}{5.3.5}.

\SGARefTarget{sga3:xxii:lem:5.4.4}\paragraph{\SGARefLink{sga3:xxii:lem:5.4.4}{Lemma 5.4.4}.}
\textit{Under the conditions of \SGARefLink{sga3:xxii:def:5.4.2}{5.4.2}, let \(R_+\) be a system of
positive roots; choose orderings on \(R'\cap R_+\) and
\(R'\cap -R_+\). The morphism}
\[
  \Omega_{R_+,R'}
  =
  \prod_{\alpha\in R'\cap -R_+}U_\alpha
  \mathop{\times}\limits_S T
  \mathop{\times}\limits_S
  \prod_{\alpha\in R'\cap R_+}U_\alpha
  \longrightarrow G,
\]
\textit{induced by multiplication in \(G\), induces an open
immersion}
\[
  \Omega_{R_+,R'}\longrightarrow H_{R'}.
\]

Indeed, by \SGARefLink{sga3:xxii:lem:5.4.3}{5.4.3} this morphism factors through \(H_{R'}\), and
therefore induces an immersion
\(\Omega_{R_+,R'}\to H_{R'}\). We then argue as in \SGARefLink{sga3:xxii:par:4.1.1}{4.1.1}.

\SGARefTarget{sga3:xxii:prop:5.4.5}\paragraph{\SGARefLink{sga3:xxii:prop:5.4.5}{Proposition 5.4.5}.}
\textit{Let \((G,T,M,R)\) be a split \(S\)-group. Let \(R'\) and
\(R''\) be two subsets of \(R\) of type \textup{(R)}.}
\begin{enumerate}
  \item[\textup{\SGAInlineRefTarget{sga3:xxii:prop:5.4.5:item:i}(i)}]
  \textit{\(H_{R'}\cap H_{R''}\) is smooth at every point of the
  unit section, \(R'\cap R''\) is of type \textup{(R)}, and}
  \[
    \bigl(H_{R'}\cap H_{R''}\bigr)^0
    =
    H_{R'\cap R''}.
  \]

  \item[\textup{\SGAInlineRefTarget{sga3:xxii:prop:5.4.5:item:ii}(ii)}]
  \textit{The following equivalence holds:}
  \[
    H_{R'}\subset H_{R''}
    \quad\Longleftrightarrow\quad
    R'\subset R''.
  \]
\end{enumerate}

Indeed, \textup{\SGARefLink{sga3:xxii:prop:5.4.5:item:ii}{(ii)}} follows immediately from \textup{\SGARefLink{sga3:xxii:prop:5.4.5:item:i}{(i)}}. To
prove \textup{\SGARefLink{sga3:xxii:prop:5.4.5:item:i}{(i)}}, it suffices to prove that
\(H_{R'}\cap H_{R''}\) is smooth at every point of the unit section:
its neutral component \textup{(cf.\ Exposé
VI\textsubscript{B}, \SGARefLink{thm:sga3-VIB-3.10}{3.10})} will then be a subgroup of type
\textup{(R)} containing \(T\), and hence equal to
\(H_{R'\cap R''}\). But
\[
  \Omega_{R_+,R'}\cap\Omega_{R_+,R''}
  =
  \Omega_{R_+,R'\cap R''}
\]
is an open subset of \(H_{R'}\cap H_{R''}\), contains the unit
section, and is smooth over \(S\).

\SGARefTarget{sga3:xxii:cor:5.4.6}\paragraph{\SGARefLink{sga3:xxii:cor:5.4.6}{Corollary 5.4.6}.}
\textit{Let \(S\) be a scheme, \(G\) a reductive \(S\)-group, \(T\)
a maximal torus of \(G\), and \(s\) a point of \(S\). If \(H\) and
\(H'\) are two subgroups of type \textup{(R)} containing \(T\), and
if \(H_s\subset H'_s\), then there is an open subset \(U\) of \(S\)
containing \(s\) such that \(H_U\subset H'_U\).}

Indeed, we may suppose that \(G\) is split relative to \(T\), and
the assertion then follows immediately from \SGARefLink{sga3:xxii:prop:5.4.5}{5.4.5} \textup{\SGARefLink{sga3:xxii:prop:5.4.5:item:ii}{(ii)}}.

We are thus led to ask which subsets \(R'\) of \(R\) are of type
\textup{(R)}. We may suppose that the group is adjoint; one must
then verify that \(\mathfrak g_{R'}\) is a Lie algebra and that its
normalizer is smooth at every point of the unit section. The most
important case is given by the following theorem.

\SGARefTarget{sga3:xxii:thm:5.4.7}\paragraph{\SGARefLink{sga3:xxii:thm:5.4.7}{Theorem 5.4.7}.}
\textit{Every closed subset \(R'\) of \(R\) is of type
\textup{(R)}.}

Recall \textup{(\SGARefLink{sga3:XXI:3.1.4}{Exposé XXI, 3.1.4})} that a subset \(R'\) of \(R\)
is said to be \emph{closed} if
\[
  \alpha,\beta\in R',\quad \alpha+\beta\in R
  \quad\Longrightarrow\quad
  \alpha+\beta\in R'.
\]

\SGARefTarget{sga3:xxii:rem:5.4.8}\paragraph{\SGARefLink{sga3:xxii:rem:5.4.8}{Remark 5.4.8}.}
We shall see later \textup{(\SGARefLink{sga3:XXIII:6.6}{Exposé XXIII, 6.6})} that, if
\(6\cdot1_S\ne0\)---for example, if \(S\) has residue characteristics
different from \(2\) and \(3\)---the fact that
\(\mathfrak g_{R'}\) is a Lie algebra already implies that \(R'\)
is closed. In that case a subset \(R'\) is of type \textup{(R)} if
and only if it is closed. Thus \SGARefLink{sga3:xxii:thm:5.4.7}{Theorem 5.4.7} gives all the subsets
of type \textup{(R)} that are ``independent of the
characteristic.''

Let us first prove the following lemma.

\SGARefTarget{sga3:xxii:lem:5.4.9}\paragraph{\SGARefLink{sga3:xxii:lem:5.4.9}{Lemma 5.4.9}.}
\textit{For each \(\alpha\in R\), choose an element}
\[
  X_\alpha\in\Gamma(S,\mathfrak g^\alpha)^\times.
\]
\textit{Let \(\alpha,\beta\in R\), with
\(\alpha+\beta\ne0\), and let \(q\) be the largest integer \(i\)
such that \(\beta+i\alpha\in R\).\begingroup
\renewcommand{\thefootnote}{*}\footnote{Translator's note: the French text prints \(\alpha+i\beta\) at this point; the displayed
expansion and the proof require \(\beta+i\alpha\).}
\addtocounter{footnote}{-1}\endgroup There exist
uniquely determined sections}
\[
  M_{\alpha,\beta,i}\in\Gamma(S,\mathcal O_S)
\]
\textit{such that}
\[
  \operatorname{Ad}\bigl(\exp_\alpha(xX_\alpha)\bigr)(X_\beta)
  =
  X_\beta+
  \sum_{i=1}^{q}
    M_{\alpha,\beta,i}x^iX_{\beta+i\alpha}
\]
\textit{for every \(x\in\mathbb G_a(S')\), where \(S'\to S\).}

The map
\[
  x\longmapsto
  \operatorname{Ad}\bigl(\exp_\alpha(xX_\alpha)\bigr)(X_\beta)
\]
defines a morphism
\[
  \mathbb G_{a,S}\longrightarrow
  \mathbf W(\mathfrak g)\simeq\mathbb G_{a,S}^{m}.
\]
There are therefore uniquely determined sections
\(Y_n\in\Gamma(S,\mathfrak g)\) such that
\[
  \operatorname{Ad}\bigl(\exp_\alpha(xX_\alpha)\bigr)(X_\beta)
  =
  \sum_{n\geq0}x^nY_n.
\]
Applying the inner automorphism defined by a section \(t\) of
\(T\), we immediately find
\[
  \operatorname{Ad}(t)(Y_n)
  =
  \beta(t)\alpha(t)^nY_n,
\]
and hence \(Y_n\in\Gamma(S,\mathfrak g^{\beta+n\alpha})\). Since
\(\alpha\) and \(\beta\) are not proportional, none of the
\(\beta+n\alpha\) is zero. Thus \(Y_n=0\) for \(n>q\), and
\[
  Y_n=M_{\alpha,\beta,n}X_{\beta+n\alpha}
\]
for \(0\leq n\leq q\), where
\(M_{\alpha,\beta,n}\in\mathbb G_a(S)\) is uniquely determined.
Setting \(x=0\) in the formula obtained gives \(Y_0=X_\beta\),
which completes the proof.

\SGARefTarget{sga3:xxii:rem:5.4.10}\paragraph{\SGARefLink{sga3:xxii:rem:5.4.10}{Remark 5.4.10}.}
Differentiating the preceding formula at \(x=0\), we find
\[
  [X_\alpha,X_\beta]
  =
  \begin{cases}
    N_{\alpha,\beta}X_{\alpha+\beta},
      &\text{where \(N_{\alpha,\beta}=M_{\alpha,\beta,1}\), if
        \(\alpha+\beta\in R\),}\\
    0,
      &\text{if \(\alpha+\beta\notin R\) and
        \(\alpha+\beta\ne0\).}
  \end{cases}
\]

Let us now prove \SGARefLink{sga3:xxii:thm:5.4.7}{5.4.7}. If \(R'\) is a closed subset of \(R\), then
\(\mathfrak g_{R'}\) is a Lie subalgebra of \(\mathfrak g\), by
\SGARefLink{sga3:xxii:rem:5.4.10}{5.4.10} and \SGARefLink{sga3:xx:2.10}{Exposé XX, 2.10}, formula \textup{(3)}. By \SGARefLink{sga3:xxii:lem:5.4.9}{5.4.9} and
\SGARefLink{sga3:xx:2.10}{Exposé XX, 2.10}, formula \textup{\SGARefLink{sga3:xxii:section:2}{(2)}}, \(U_\alpha\) normalizes
\(\mathfrak g_{R'}\) for every \(\alpha\in R'\). Choose a system of
positive roots \(R_+\), and consider the open subset
\(\Omega_{R_+}\) of \SGARefLink{sga3:xxii:prop:4.1.2}{4.1.2}. Let \(\Omega_{R_+,R'}\) be the closed
subscheme of \(\Omega_{R_+}\) defined by
\[
  \Omega_{R_+,R'}
  =
  \left(
    \prod_{\alpha\in R'\cap -R_+}U_\alpha
  \right)
  \cdot T\cdot
  \left(
    \prod_{\alpha\in R'\cap R_+}U_\alpha
  \right).
\]
The canonical immersion \(\Omega_{R_+,R'}\to G\) factors through
\[
  i:\Omega_{R_+,R'}
  \longrightarrow
  \underline{\operatorname{Norm}}_G(\mathfrak g_{R'}).
\]
Suppose that \(G\) is adjoint. The tangent map of \(i\) at the
points of the unit section is bijective by \SGARefLink{sga3:xxii:lemma:5.3.2}{5.3.2}. In particular,
\(i\) is étale at every point of the unit section, and hence is a
local immersion at every point of the unit section.
Therefore
\(\underline{\operatorname{Norm}}_G(\mathfrak g_{R'})\) is smooth
at every point of the unit section, as was to be proved.

\SGARefTarget{sga3:xxii:section:5.5}
\subsection*{\SGAInlineRefTarget{sga3:xxii:subsection:5-5}5.5. Borel Subgroups of a Split Reductive Group}
\addcontentsline{toc}{subsection}{\SGARefLink{sga3:xxii:subsection:5-5}{5.5}. Borel Subgroups of a Split
Reductive Group}

\SGARefTarget{sga3:xxii:prop:5.5.1}\paragraph{\SGARefLink{sga3:xxii:prop:5.5.1}{Proposition 5.5.1}.}
\textit{Let \((G,T,M,R)\) be a split \(S\)-group. For every system
of positive roots \(R_+\) of \(R\), the subgroup \(H_{R_+}\)
\textup{(which exists by \SGARefLink{sga3:xxii:thm:5.4.7}{5.4.7})} is a Borel subgroup of \(G\).
Moreover, for every ordering on \(R_+\), the morphism induced by
multiplication in \(G\),}
\[
  T\mathop{\times}\limits_S
  \prod_{\alpha\in R_+}U_\alpha
  \longrightarrow G,
\]
\textit{is a closed immersion with image \(H_{R_+}\). We write
\(B_{R_+}=H_{R_+}\).}

By the definition of Borel subgroups, the first assertion may be
verified after replacing \(S\) by the spectrum of an algebraically
closed field. Let \(B\) be the Borel subgroup of \(G\) containing
\(T\) and corresponding to the system of positive roots \(R_+\)
\textup{(\emph{Bible}, \S10.4, Proposition 9)}. The Lie algebra of
\(B\) is \(\mathfrak g_{R_+}\); therefore \(B=H_{R_+}\) by \SGARefLink{sga3:xxii:cor:5.3.5}{5.3.5}.

For the second assertion, the morphism in the statement induces an
open immersion
\[
  i:
  T\mathop{\times}\limits_S
  \prod_{\alpha\in R_+}U_\alpha
  \longrightarrow H_{R_+}
\]
by \SGARefLink{sga3:xxii:lem:5.4.4}{5.4.4}. But \(i\) is surjective
\textup{(\emph{Bible}, \S15.1, Corollary 1 to Proposition 1)}.

\SGARefTarget{sga3:xxii:cor:5.5.2}\paragraph{\SGARefLink{sga3:xxii:cor:5.5.2}{Corollary 5.5.2}.}
\textit{Choose an arbitrary ordering on \(R_+\), and for each
\(\alpha\in R_+\) choose an}
\[
  X_\alpha\in\Gamma(S,\mathfrak g^\alpha)^\times.
\]
\textit{Let \(\alpha,\beta\in R_+\). For every pair
\((i,j)\in\mathbb N^*\times\mathbb N^*\) such that
\(i\alpha+j\beta\in R\), there is a unique section}
\[
  C_{i,j,\alpha,\beta}\in\Gamma(S,\mathcal O_S)
\]
\textit{such that, for all \(x,y\in\mathbb G_a(S')\), where
\(S'\to S\), one has}
\[
  \exp_\alpha(xX_\alpha)
  \exp_\beta(yX_\beta)
  \exp_\alpha(xX_\alpha)^{-1}
  =
  \exp_\beta(yX_\beta)
  \prod_{\substack{i,j\in\mathbb N^*\\ i\alpha+j\beta\in R}}
    \exp_{i\alpha+j\beta}
      \bigl(C_{i,j,\alpha,\beta}x^iy^jX_{i\alpha+j\beta}\bigr).
\]

If \(\alpha=\beta\), the assertion is trivial. Suppose therefore
that \(\alpha\ne\beta\). By the proposition there are unique
morphisms
\[
  F_0:\mathbb G_{a,S}^{\,2}\longrightarrow T,
  \qquad
  F_\gamma:\mathbb G_{a,S}^{\,2}\longrightarrow\mathbb G_{a,S}
  \quad(\gamma\in R_+)
\]
such that
\[
  \exp(xX_\alpha)\exp(yX_\beta)\exp(xX_\alpha)^{-1}
  =
  F_0(x,y)
  \prod_{\gamma\in R_+}
    \exp\bigl(F_\gamma(x,y)X_\gamma\bigr).
\]
Let \(t\in T(S')\), where \(S'\to S\). Applying
\(\operatorname{int}(t)\) to this formula, we immediately obtain
\[
  \tag{1}\label{sga3:xxii:subsection:5-5:equation:1}
  F_0\bigl(\alpha(t)x,\beta(t)y\bigr)=F_0(x,y),
\]
\[
  \tag{2}\label{sga3:xxii:subsection:5-5:equation:2}
  F_\gamma\bigl(\alpha(t)x,\beta(t)y\bigr)
  =
  \gamma(t)F_\gamma(x,y).
\]
Since \(\alpha\) and \(\beta\) are two linearly independent
characters of \(T\) over \(\mathbb Q\), the first relation implies,
as usual, that \(F_0\) is constant, and hence
\(F_0(x,y)=e\). Write next
\[
  F_\gamma(x,y)=\sum a_{ij}x^iy^j,
  \qquad a_{ij}\in\Gamma(S,\mathcal O_S).
\]
Substituting in relation \textup{\SGARefLink{sga3:xxii:subsection:5-5:equation:2}{(2)}} and comparing the
polynomials on both sides gives
\[
  a_{ij}
  \bigl(\alpha(t)^i\beta(t)^j-\gamma(t)\bigr)=0.
\]
If \(\gamma\ne i\alpha+j\beta\), we know
\textup{(\SGARefLink{sga3:xix:lemma:4:13}{Exposé XIX, 4.13})} that there exist a faithfully flat
quasi-compact morphism \(S'\to S\) and an element \(t\in T(S')\)
such that
\[
  \alpha(t)^i\beta(t)^j-\gamma(t)=1.
\]
Thus \(a_{ij}=0\) over \(S'\), and hence over \(S\). If
\(\gamma=i\alpha+j\beta\), set
\(a_{ij}=C_{i,j,\alpha,\beta}\). Setting \(x=0\)
\textup{(respectively, \(y=0\))}, one obtains
\[
  C_{0,1,\alpha,\beta}=1
  \qquad
  \textup{(respectively, \(C_{1,0,\alpha,\beta}=0\)).}
\]

\SGARefTarget{sga3:xxii:rem:5.5.3}\paragraph{\SGARefLink{sga3:xxii:rem:5.5.3}{Remark 5.5.3}.}
Differentiating with respect to \(y\) at \(y=0\) and comparing
with \SGARefLink{sga3:xxii:lem:5.4.9}{5.4.9}, one finds
\[
  C_{i,1,\alpha,\beta}=M_{\alpha,\beta,i}.
\]

\SGARefTarget{sga3:xxii:cor:5.5.4}\paragraph{\SGARefLink{sga3:xxii:cor:5.5.4}{Corollary 5.5.4}.}
\textit{Let \(S\) be a scheme, let \(G\) be a reductive
\(S\)-group, let \(T\) be a maximal torus of \(G\), and let
\(\alpha\ne\beta\) be two roots of \(G\) relative to \(T\) such
that \(\alpha+\beta\) is nontrivial on every fiber. Order
arbitrarily the set of roots \(i\alpha+j\beta\), where
\((i,j)\in\mathbb N^*\times\mathbb N^*\). For all
\(i,j\in\mathbb N^*\) such that \(i\alpha+j\beta\in R\), there is
a unique morphism of \(\mathcal O_S\)-modules}
\[
  f_{\alpha,\beta,i,j}:
  (\mathfrak g^\alpha)^{\otimes i}
  \otimes
  (\mathfrak g^\beta)^{\otimes j}
  \longrightarrow
  \mathfrak g^{i\alpha+j\beta}.
\]
\textit{such that, for every \(S'\to S\) and all
\(X\in\mathbf W(\mathfrak g^\alpha)(S')\),
\(Y\in\mathbf W(\mathfrak g^\beta)(S')\), one has
\textup{(the exponentials in the right-hand side being understood
in the sense of \SGARefLink{sga3:xxii:rem:1.2}{1.2})}}
\[
  \exp_\alpha(X)\exp_\beta(Y)\exp_\alpha(-X)
  =
  \exp_\beta(Y)
  \prod_{(i,j)}
    \exp_{i\alpha+j\beta}
      \bigl(f_{\alpha,\beta,i,j}(X^{\otimes i}\otimes
      Y^{\otimes j})\bigr).
\]

The assertion is local for the fpqc topology. By \S2 we may
therefore suppose that \(G\) is split relative to \(T\), with
\(\alpha\) and \(\beta\) constant in the splitting. Since
\(\alpha+\beta\ne0\), there is a system of positive roots \(R_+\)
containing \(\alpha\) and \(\beta\) \textup{(\SGARefLink{sga3:XXI:3.5.4}{Exposé XXI, 3.5.4})},
and we are reduced to \SGARefLink{sga3:xxii:cor:5.5.2}{5.5.2}.

\SGARefTarget{sga3:xxii:cor:5.5.5}\paragraph{\SGARefLink{sga3:xxii:cor:5.5.5}{Corollary 5.5.5}.}
\textit{Let \(S\) be a scheme and \(G\) a reductive
\(S\)-group.}
\begin{enumerate}
\item[\SGAInlineRefTarget{sga3:xxii:cor:5.5.5:item:i}(i)] \textit{\(G\) has Borel subgroups locally for the
étale topology. If \(T\) is a maximal torus of \(G\), then,
locally for the étale topology, \(G\) also has Borel subgroups
containing \(T\).}
\item[\SGAInlineRefTarget{sga3:xxii:corollary:5-5-5:item:ii}\SGAInlineRefTarget{sga3:xxii:cor:5.5.5:item:ii}(ii)] \textit{If \(T\) is a maximal torus of \(G\), the
``functor of Borel subgroups of \(G\) containing \(T\)'' is
represented by a principal homogeneous bundle under \(W_G(T)\).}
\item[\SGAInlineRefTarget{sga3:xxii:cor:5.5.5:item:iii}(iii)] \textit{If \((G,T,M,R)\) is split, every Borel
subgroup \(B\) of \(G\) containing \(T\) is locally on \(S\) of
the form \(B_{R_+}\), where \(R_+\) is a system of positive roots
of \(R\).}
\item[\SGAInlineRefTarget{sga3:xxii:cor:5.5.5:item:iv}(iv)] \textit{If \(T\subset B\) is a Killing pair of \(G\),
there is an étale covering family \(\{S_i\to S\}\), and for every
\(i\) a splitting \((G_{S_i},T_{S_i},M_i,R_i)\) and a system of
positive roots \(R_{i+}\) of \(R_i\), such that
\(B_{S_i}=B_{R_{i+}}\).}
\end{enumerate}

Indeed, \textup{\SGARefLink{sga3:xxii:cor:5.5.5:item:i}{(i)}} follows from \SGARefLink{sga3:xxii:cor:2.3}{2.3} and \SGARefLink{sga3:xxii:prop:5.5.1}{5.5.1},
\textup{\SGARefLink{sga3:xxii:cor:5.5.5:item:ii}{(ii)}} from \textup{\SGARefLink{sga3:xxii:cor:5.5.5:item:i}{(i)}} and \SGARefLink{sga3:xxii:cor:5.3.15}{5.3.15},
\textup{\SGARefLink{sga3:xxii:cor:5.5.5:item:iii}{(iii)}} from \textup{\SGARefLink{sga3:xxii:cor:5.5.5:item:ii}{(ii)}} and \SGARefLink{sga3:xxii:prop:5.5.1}{5.5.1}, and
\textup{\SGARefLink{sga3:xxii:cor:5.5.5:item:iv}{(iv)}} from \textup{\SGARefLink{sga3:xxii:cor:5.5.5:item:iii}{(iii)}} and \SGARefLink{sga3:xxii:cor:2.3}{2.3}.

\SGARefTarget{sga3:xxii:lem:5.5.6}\paragraph{\SGARefLink{sga3:xxii:lem:5.5.6}{Lemma 5.5.6}.}
\textit{Choose on the group \(\Gamma_0(R)\) generated by the
roots a totally ordered group structure such that the roots
\(>0\) are the elements of \(R_+\)
\textup{(cf. \SGARefLink{sga3:XXI:3.5.6}{Exposé XXI, 3.5.6})}
Let \(\alpha_1<\cdots<\alpha_N\) be the elements of \(R_+\).
Consider the isomorphism}
\[
  f:
  T\mathop{\times}\limits_S U_{\alpha_1}
  \mathop{\times}\limits_S\cdots
  \mathop{\times}\limits_S U_{\alpha_N}
  \longrightarrow B_{R_+}
\]
\textit{induced by multiplication in \(G\). For
\(i=1,\ldots,N\), put}
\[
  U_{\geq i}
  =
  f\left(
    U_{\alpha_i}\mathop{\times}\limits_S\cdots
    \mathop{\times}\limits_S U_{\alpha_N}
  \right).
\]
\begin{enumerate}
\item[\SGAInlineRefTarget{sga3:xxii:lem:5.5.6:item:i}(i)] \textit{Each \(U_{\geq i}\) is an invariant subgroup
of \(B_{R_+}\).}
\item[\SGAInlineRefTarget{sga3:xxii:lem:5.5.6:item:ii}(ii)] \textit{For \(1\leq i\leq N-1\),
\(U_{\geq i}\) identifies with the semidirect product}
\[
  U_{\geq i}=U_{\alpha_i}\cdot U_{\geq i+1}.
\]
\item[\SGAInlineRefTarget{sga3:xxii:lem:5.5.6:item:iii}(iii)] \textit{\(B_{R_+}\) identifies with the semidirect
product}
\[
  B_{R_+}=T\cdot U_{\geq1}.
\]
\item[\SGAInlineRefTarget{sga3:xxii:lem:5.5.6:item:iv}(iv)] \textit{For \(1\leq i\leq N-1\), the inner
automorphisms of \(U_{\geq1}\) act trivially on
\(U_{\geq i}/U_{\geq i+1}\), identified with \(U_{\alpha_i}\)
by \textup{\SGARefLink{sga3:xxii:lem:5.5.6:item:ii}{(ii)}}.}
\end{enumerate}

We first prove, by descending induction on \(i\), the following
assertion:
\[
  U_{\geq i}\ \text{is an invariant subgroup of }B_{R_+},
  \quad\text{the semidirect product of }
  U_{\alpha_i}\text{ and }U_{\geq i+1}.
\]
This is true for \(i=N\); suppose it has been proved for \(i+1\),
and let us prove it for \(i\). As schemes, we have
\[
  U_{\geq i}=U_{\alpha_i}\cdot U_{\geq i+1}.
\]
We must show that \(U_{\geq i}\) is invariant under the inner
automorphisms of \(B_{R_+}\). This is clear for
\(\operatorname{int}(t)\), where \(t\in T(S)\); it is therefore
enough to consider \(\operatorname{int}(x)\), where
\(x\in U_\alpha(S)\) and \(\alpha\in R_+\). Since
\(U_{\geq i+1}\) is invariant, it is enough to prove that
\[
  \operatorname{int}(x)(U_{\alpha_i})\subset U_{\geq i}.
\]
By \SGARefLink{sga3:xxii:cor:5.5.2}{5.5.2}, if \(y\in U_{\alpha_i}(S')\), then
\[
  y^{-1}xyx^{-1}\in U_{\geq i+1}(S'),
\]
which implies that
\[
  \operatorname{int}(x)(y)\in U_{\geq i}(S').
\]

It remains to prove that \(U_{\geq i}\) is a subgroup. Let
\(x,y\in U_{\geq i}(S)\), and write
\[
  x=px',\qquad y=qy',
\]
where \(p,q\in U_{\alpha_i}(S)\) and
\(x',y'\in U_{\geq i+1}(S)\). We then have
\[
  xy=px'qy'=pq(q^{-1}x'q)y'
  \in U_{\alpha_i}(S')U_{\geq i+1}(S'),
\]
and likewise
\[
  x^{-1}
  =
  p^{-1}\bigl(p{x'}^{-1}p^{-1}\bigr)
  \in U_{\alpha_i}(S')U_{\geq i+1}(S').
\]
We have thus proved \textup{\SGARefLink{sga3:xxii:lem:5.5.6:item:i}{(i)}} and \textup{\SGARefLink{sga3:xxii:lem:5.5.6:item:ii}{(ii)}}, as well as
\textup{\SGARefLink{sga3:xxii:lem:5.5.6:item:iv}{(iv)}} along the way. Assertion \textup{\SGARefLink{sga3:xxii:lem:5.5.6:item:iii}{(iii)}} is a trivial
consequence of \SGARefLink{sga3:xxii:prop:5.5.1}{5.5.1}.

\SGARefTarget{sga3:xxii:lem:5.5.7}\paragraph{\SGARefLink{sga3:xxii:lem:5.5.7}{Lemma 5.5.7}.}
\textit{With the preceding notation, choose, for each
\(1\leq i\leq N\), an element}
\[
  X_i\in\Gamma(S,\mathfrak g^{\alpha_i})^\times,
\]
\textit{and consider the isomorphism}
\[
  a:\mathbb G_{a,S}^{\,N}\longrightarrow U_{\geq1}
\]
\textit{defined on underlying sets by}
\[
  a(x_1,\ldots,x_N)
  =
  \exp_{\alpha_1}(x_1X_1)\cdots
  \exp_{\alpha_N}(x_NX_N).
\]
\textit{There is a unique family \((Q_i)_{1\leq i\leq N}\) of
polynomials}
\[
  Q_i
  =
  Q_i(x_1,\ldots,x_N,y_1,\ldots,y_N)
\]
\textit{with coefficients in \(\Gamma(S,\mathcal O_S)\), such
that}
\[
\begin{split}
  &a(x_1,\ldots,x_N)a(y_1,\ldots,y_N)\\
  &\qquad =
  a\bigl(
    Q_1(x_1,\ldots,x_N,y_1,\ldots,y_N),
    \ldots,
    Q_N(x_1,\ldots,x_N,y_1,\ldots,y_N)
  \bigr).
\end{split}
\]
\textit{Moreover, the coefficients of these polynomials belong
to the subring of \(\Gamma(S,\mathcal O_S)\) generated by the
\(C_{i,j,\alpha,\beta}\) of \SGARefLink{sga3:xxii:cor:5.5.2}{5.5.2}
\(\bigl(\alpha,\beta\in R_+,\ i,j\in\mathbb N^*\bigr)\), and
each \(Q_i\) is of the form}
\[
\begin{split}
  Q_i(x_1,\ldots,x_N,y_1,\ldots,y_N)
  ={}&x_i+y_i\\
  &+
  Q_i'(x_1,\ldots,x_{i-1},y_1,\ldots,y_{i-1}).
\end{split}
\]

The existence and uniqueness of the \(Q_i\) are immediate from
the fact that \(a\) is an isomorphism of schemes. Let
\(z,z',z''\) denote sections of \(U_{\geq i+1}\). We have
\[
\begin{split}
  &a(x_1,\ldots,x_N)a(y_1,\ldots,y_N)\\
  &\quad =
  a(x_1,\ldots,x_{i-1},0,\ldots,0)
  \exp_{\alpha_i}(x_iX_i)z\\
  &\qquad\quad{}\cdot
  a(y_1,\ldots,y_{i-1},0,\ldots,0)
  \exp_{\alpha_i}(y_iX_i)z'.
\end{split}
\]
Using \SGARefLink{sga3:xxii:lem:5.5.6}{5.5.6} \textup{\SGARefLink{sga3:xxii:lem:5.5.6:item:i}{(i)}} and \textup{\SGARefLink{sga3:xxii:lem:5.5.6:item:iv}{(iv)}}, the right-hand side
can be written
\[
\begin{split}
  &a(x_1,\ldots,x_{i-1},0,\ldots,0)
  a(y_1,\ldots,y_{i-1},0,\ldots,0)\\
  &\qquad{}\cdot
  \exp_{\alpha_i}\bigl((x_i+y_i)X_i\bigr)z'',
\end{split}
\]
and, applying \SGARefLink{sga3:xxii:lem:5.5.6}{5.5.6} once again, this gives
\[
  Q_i(x_1,\ldots,x_N,y_1,\ldots,y_N)
  =
  x_i+y_i+
  Q_i'(x_1,\ldots,x_{i-1},y_1,\ldots,y_{i-1}),
\]
where
\[
  Q_i'(x_1,\ldots,x_{i-1},y_1,\ldots,y_{i-1})
  =
  Q_i(x_1,\ldots,x_{i-1},0,\ldots,0,
      y_1,\ldots,y_{i-1},0,\ldots,0);
\]
this is precisely the required form.\footnote{Translator's note: The
French source prints \(x_i\) as the first argument on the right-hand
side of this formula. We have written \(x_1\), as required by the
defining specialization and by the preceding calculation.}

It remains to prove the assertion concerning the coefficients of the
polynomials \(Q_i\). Let \(A\) be the subring of
\(\Gamma(S,\mathcal O_S)\) generated by the
\(C_{i,j,\alpha,\beta}\), where
\(\alpha,\beta\in R_+\) and \(i,j\in\mathbb N^*\).
We prove by descending induction on \(i\) that, if
\[
  x_1=\cdots=x_{i-1}=0
  \qquad\text{and}\qquad
  y_1=\cdots=y_{i-1}=0,
\]
that is, if \(a(x_1,\ldots,x_N)\) and \(a(y_1,\ldots,y_N)\)
are sections of \(U_{\geq i}\), then the polynomials
\[
  R_j(x_i,\ldots,x_N,y_i,\ldots,y_N)
  =
  Q_j(x_1,\ldots,x_N,y_1,\ldots,y_N)
\]
have coefficients in \(A\). This is trivial for \(i=N\), and
also for \(j<i\), since \(R_j=0\) when \(j<i\). Let \(i<N\);
suppose the assertion proved for \(i+1\), and let us prove it
for \(i\) and \(j\geq i\). We have
\[
\begin{split}
  a(0,\ldots,0,x_i,\ldots,x_N)
  &=
  \exp(x_iX_i)
  a(0,\ldots,0,x_{i+1},\ldots,x_N)\\
  &=
  \exp(x_iX_i)Z_i.
\end{split}
\]
Similarly, write
\[
  a(0,\ldots,y_i,\ldots,y_N)=\exp(y_iX_i)T_i.
\]
Then
\[
\begin{split}
  &a(0,\ldots,x_i,\ldots,x_N)
   a(0,\ldots,y_i,\ldots,y_N)\\
  &\qquad =
  \exp\bigl((x_i+y_i)X_i\bigr)
  \operatorname{int}\bigl(\exp(-y_iX_i)\bigr)(Z_i)\cdot T_i.
\end{split}
\]
Now\footnote{Translator's note: In the second conjugating exponential
the French source prints the undefined symbol \(Y_i\). We have
written \(X_i\), as required by the preceding identity and by the
definition of \(Z_i\).}
\[
\begin{split}
  \operatorname{int}\bigl(\exp(-y_iX_i)\bigr)(Z_i)
  &=
  \operatorname{int}\bigl(\exp(-y_iX_i)\bigr)
  \bigl(
    \exp(x_{i+1}X_{i+1})\cdots\exp(x_NX_N)
  \bigr)
\end{split}
\]
is a product of \(N-i-1\) sections of \(U_{\geq i+1}\), whose
coefficients in the decomposition
\[
  U_{\geq i+1}
  =
  U_{\alpha_{i+1}}\cdots U_{\alpha_N}
\]
are polynomials in \(y_i,x_{i+1},\ldots,x_N\) with coefficients
in \(A\), by \SGARefLink{sga3:xxii:cor:5.5.2}{5.5.2}. Applying the induction hypothesis, we deduce
that the coefficients of
\[
  \operatorname{int}\bigl(\exp(-y_iX_i)\bigr)(Z_i)\cdot T_i
\]
are likewise polynomials with coefficients in \(A\), which
completes the proof.

The preceding induction immediately gives a proof of the following.

\SGARefTarget{sga3:xxii:lem:5.5.8}\paragraph{\SGARefLink{sga3:xxii:lem:5.5.8}{Lemma 5.5.8}.}
\textit{With the notation of \SGARefLink{sga3:xxii:lem:5.5.6}{5.5.6}, let, for each
\(i=1,\ldots,N\),}
\[
  f_i:U_{\alpha_i}\longrightarrow H
\]
\textit{be a group morphism, where \(H\) is an \(S\)-group
functor. For the morphism}
\[
  f:U_{\geq1}\longrightarrow H
\]
\textit{defined by}
\[
  f\bigl(
    \exp(x_1X_1)\cdots\exp(x_NX_N)
  \bigr)
  =
  f_1\bigl(\exp(x_1X_1)\bigr)\cdots
  f_N\bigl(\exp(x_NX_N)\bigr)
\]
\textit{to be a group morphism, it is necessary and sufficient
that, for every pair \(i<j\),}
\[
\begin{split}
  &f_j\bigl(\exp(x_jX_j)\bigr)
   f_i\bigl(\exp(x_iX_i)\bigr)
   f_j\bigl(\exp(-x_jX_j)\bigr)\\
  &\qquad =
  f\bigl(
    \exp(x_jX_j)\exp(x_iX_i)\exp(-x_jX_j)
  \bigr).
\end{split}
\]

\SGARefTarget{sga3:xxii:sec:5.6}
\subsection*{\SGAInlineRefTarget{sga3:xxii:subsection:5-6}5.6. Subgroups of Type \textup{(R)} with Solvable
Fibers}

\SGARefTarget{sga3:xxii:prop:5.6.1}
\paragraph{\SGARefLink{sga3:xxii:prop:5.6.1}{Proposition 5.6.1}.}
\textit{Let \((G,T,M,R)\) be a split \(S\)-group, let \(R'\) be
a subset of \(R\) of type \textup{(R)} \textup{(\SGARefLink{sga3:xxii:def:5.4.2}{5.4.2})}, and let
\(H_{R'}\) be the corresponding subgroup of \(G\). The following
conditions are equivalent:}
\begin{enumerate}
\item[\textup{\SGAInlineRefTarget{sga3:xxii:prop:5.6.1:item:i}(i)}]
\textit{\(H_{R'}\) has solvable geometric fibers.}

\item[\textup{\SGAInlineRefTarget{sga3:xxii:prop:5.6.1:item:ii}(ii)}]
\textit{There exists a system of positive roots \(R_+\) such that
\(R'\subset R_+\), and hence \(H_{R'}\subset B_{R_+}\)
\textup{(cf.\ \SGARefLink{sga3:xxii:prop:5.4.5}{5.4.5})}.}

\item[\textup{\SGAInlineRefTarget{sga3:xxii:prop:5.6.1:item:iii}(iii)}]
\textit{\(R'\cap(-R')=\varnothing\).}

\item[\textup{\SGAInlineRefTarget{sga3:xxii:prop:5.6.1:item:iv}(iv)}]
\textit{For every ordering of \(R'\), the morphism induced by
multiplication in \(G\)}
\[
  T\times_S\prod_{\alpha\in R'}U_\alpha
  \longrightarrow H_{R'}
\]
\textit{is an isomorphism.}

\item[\textup{\SGAInlineRefTarget{sga3:xxii:prop:5.6.1:item:v}(v)}]
\textit{\(
  H_{R'}\cap\underline{\operatorname{Norm}}_G(T)=T.
\)}

\item[\textup{\SGAInlineRefTarget{sga3:xxii:prop:5.6.1:item:vi}(vi)}]
\textit{For every subset \(R''\) of \(R\) of type \textup{(R)},
one has \textup{(cf.\ \SGARefLink{sga3:xxii:prop:5.4.5}{5.4.5})}}
\[
  H_{R'}\cap
  \underline{\operatorname{Norm}}_G(H_{R''})
  =
  H_{R'\cap R''}.
\]
\end{enumerate}

We shall prove these equivalences according to the following logical
scheme.

\SGARefTarget{sga3:xxii:diagram:5.6.1:implication-scheme}
\begin{center}
  \begin{tikzcd}[column sep=5.2em,row sep=2.9em]
  \textup{(iii)}
    \arrow[dd,Rightarrow]
  &
  \textup{(ii)}
    \arrow[l,Rightarrow]
    \arrow[r,Rightarrow]
    \arrow[d,Rightarrow]
  &
  \textup{(vi)}
    \arrow[dd,Rightarrow]
  \\
  &
  \textup{(iv)}
    \arrow[dr,Rightarrow]
  &
  \\
  \textup{(i)}
    \arrow[ur,Rightarrow]
  &&
  \textup{(v)}.
    \arrow[ll,Rightarrow]
\end{tikzcd}

\end{center}

Clearly \textup{\SGARefLink{sga3:xxii:prop:5.6.1:item:ii}{(ii)}} implies \textup{\SGARefLink{sga3:xxii:prop:5.6.1:item:iii}{(iii)}}, and \textup{\SGARefLink{sga3:xxii:prop:5.6.1:item:vi}{(vi)}}
implies \textup{\SGARefLink{sga3:xxii:prop:5.6.1:item:v}{(v)}} \textup{(take \(R''=\varnothing\))}. By \SGARefLink{sga3:xxii:cor:5.4.6}{5.4.6},
it is enough to verify that \textup{\SGARefLink{sga3:xxii:prop:5.6.1:item:i}{(i)}} implies \textup{\SGARefLink{sga3:xxii:prop:5.6.1:item:ii}{(ii)}} on
the geometric fibers. If \(S\) is the spectrum of an algebraically
closed field, then \(H_{R'}\) is contained in a Borel subgroup
containing \(T\), and this subgroup is therefore of the form
\(H_{R_+}\) \textup{(\SGARefLink{sga3:xxii:cor:5.5.5}{5.5.5}\textup{\SGARefLink{sga3:xxii:cor:5.5.5:item:iii}{(iii)}})}.

Likewise, the implication from \textup{\SGARefLink{sga3:xxii:prop:5.6.1:item:iii}{(iii)}} to \textup{\SGARefLink{sga3:xxii:prop:5.6.1:item:i}{(i)}} may be
checked on the geometric fibers. Suppose that \textup{\SGARefLink{sga3:xxii:prop:5.6.1:item:iii}{(iii)}} holds.
If \(H_{R'}\) were not solvable, there would exist a subtorus \(Q\)
of \(T\), of codimension \(1\) in \(T\), such that
\[
  \underline{\operatorname{Centr}}_{H_{R'}}(Q)
\]
is not solvable \textup{(\emph{Bible}, \S10.4, Proposition 8)}.
But \(\underline{\operatorname{Centr}}_{H_{R'}}(Q)\) has Lie algebra
\(\mathfrak g_{R''}\), where \(R''\) is the set of roots in \(R'\)
which vanish on \(Q\). Thus, by \textup{\SGARefLink{sga3:xxii:prop:5.6.1:item:iii}{(iii)}},
\[
  R''=\varnothing\qquad\text{or}\qquad R''=\{\alpha\}.
\]
Consequently,
\(\underline{\operatorname{Centr}}_{H_{R'}}(Q)\), which is a subgroup
of \(G\) of type \textup{(R)}, is \(T\) or \(T\cdot U_\alpha\), and
is therefore solvable, contrary to the hypothesis.

Likewise, the implication from \textup{\SGARefLink{sga3:xxii:prop:5.6.1:item:ii}{(ii)}} to \textup{\SGARefLink{sga3:xxii:prop:5.6.1:item:iv}{(iv)}} may be
checked on the geometric fibers \textup{(since the schemes in
question are flat \(S\)-schemes of finite presentation)}. By
\emph{Bible}, \S13.2, Theorem 1\textup{\SGARefLink{sga3:xxii:rem:5.1.2:item:d}{(d)}}, the morphism in question
is bijective. It induces an isomorphism on the tangent spaces at the
origin, and the conclusion follows as usual \textup{(cf.\ \SGARefLink{sga3:xxii:par:4.1.1}{4.1.1})}.

The implication from \textup{\SGARefLink{sga3:xxii:prop:5.6.1:item:iv}{(iv)}} to \textup{\SGARefLink{sga3:xxii:prop:5.6.1:item:v}{(v)}} follows from
\SGARefLink{sga3:xxii:lem:4.2.7}{4.2.7}. To prove that \textup{\SGARefLink{sga3:xxii:prop:5.6.1:item:v}{(v)}} implies \textup{\SGARefLink{sga3:xxii:prop:5.6.1:item:i}{(i)}}, we may again
reduce to the case in which \(S\) is the spectrum of an algebraically
closed field; the conclusion then follows from \emph{Bible}, \S10.3,
the corollary to \SGARefLink{sga3:xxii:section:6}{Proposition 6}, and \S9.3, Corollary 3 to Theorem 1.

It remains only to prove the implication from \textup{\SGARefLink{sga3:xxii:prop:5.6.1:item:ii}{(ii)}} to
\textup{\SGARefLink{sga3:xxii:prop:5.6.1:item:vi}{(vi)}}. We may reduce to the case in which \(G\) is adjoint.
Then, by \SGARefLink{sga3:xxii:cor:5.3.3}{5.3.3},
\[
  \underline{\operatorname{Norm}}_G(H_{R''})
  =
  \underline{\operatorname{Norm}}_G(\mathfrak g_{R''})
  \subset
  \underline{\operatorname{Transp}}_G
    (\mathfrak t,\mathfrak g_{R''}).
  \tag{39}\label{sga3:xxii:subsection:5-6:equation:39}
\]
By \SGARefLink{sga3:xxii:prop:5.4.5}{5.4.5}, it is enough to prove
\[
  H_{R'}(S)\cap
  \underline{\operatorname{Transp}}_{G(S)}
    (\mathfrak t,\mathfrak g_{R''})
  \subset
  H_{R'\cap R''}(S).
  \tag{x}\label{sga3:xxii:subsection:5-6:equation:x}
\]
We first prove a lemma.

\SGARefTarget{sga3:xxii:lem:5.6.2}
\paragraph{\SGARefLink{sga3:xxii:lem:5.6.2}{Lemma 5.6.2}.}
\textit{With the notation of \SGARefLink{sga3:xxii:lem:5.5.7}{5.5.7}, let}
\[
  u=\exp(x_1X_1)\cdots\exp(x_NX_N),
\]
\textit{where \(x_i\in\mathbb G_a(S)\). Let \(m\), with
\(1\leq m\leq N\), be such that \(x_i=0\) for \(i<m\).}
\begin{enumerate}
\item[\textup{\SGAInlineRefTarget{sga3:xxii:lem:5.6.2:item:a}(a)}]
\textit{If \(H\in\Gamma(S,\mathfrak t)\), the component of
\(\operatorname{Ad}(u)H\) in \(\mathfrak g^{\alpha_m}\) is}
\[
  -\overline{\alpha_m}(H)x_mX_m.
\]

\item[\textup{(b)}]
\textit{If \(Y\in\Gamma(S,\mathfrak g^{-\alpha_m})\), the component
of \(\operatorname{Ad}(u)Y\) in \(\mathfrak t\) is, with the notation
of \SGARefLink{sga3:xx:2.6}{Exposé XX, 2.6},}
\[
  x_m\langle X_m,Y\rangle H_{\alpha_m}.
\]
\end{enumerate}

\textit{Proof.}
Set
\[
  \mathfrak u
  =
  \mathfrak g^{\alpha_{m+1}}+\cdots+\mathfrak g^{\alpha_N}.
\]
By \SGARefLink{sga3:xxii:lem:5.4.9}{5.4.9}, one has
\[
  \operatorname{Ad}\bigl(\exp(x_iX_i)\bigr)\mathfrak u
  \subset\mathfrak u
  \qquad\text{for }i>m.
\]
By \SGARefLink{sga3:xx:2.10}{Exposé XX, 2.10},
\[
  \operatorname{Ad}\bigl(\exp(x_iX_i)\bigr)H
  =
  H-\overline{\alpha_i}(H)x_iX_i.
\]
It follows immediately that, modulo \(\mathfrak u\),
\[
  \operatorname{Ad}(u)H
  =
  \operatorname{Ad}\bigl(\exp(x_mX_m)\bigr)H
  =
  H-\overline{\alpha_m}(H)x_mX_m,
\]
which proves the first assertion.

Likewise, set
\[
  \mathfrak n
  =
  \mathfrak g^{\alpha_1}+\cdots+\mathfrak g^{\alpha_N},
  \qquad
  u_1
  =
  \exp(x_{m+1}X_{m+1})\cdots\exp(x_NX_N).
\]
For \(i>m\), one has \(\alpha_i>\alpha_m\); hence, by \SGARefLink{sga3:xxii:lem:5.4.9}{5.4.9}, modulo
\(\mathfrak n\),
\[
  \operatorname{Ad}(u_1)Y\equiv Y,
  \qquad\text{and therefore}\qquad
  \operatorname{Ad}(u)Y
  \equiv
  \operatorname{Ad}\bigl(\exp(x_mX_m)\bigr)Y.
\]
Applying \SGARefLink{sga3:xx:2.10}{Exposé XX, 2.10}, one therefore obtains, modulo
\(\mathfrak n\),
\[
  \operatorname{Ad}(u)Y-Y
  \equiv
  x_m\langle X_m,Y\rangle H_{\alpha_m},
\]
which proves the second assertion.

Let us return to the proof of inclusion \textup{\SGARefLink{sga3:xxii:subsection:5-6:equation:x}{(x)}}. Suppose that
there exists
\[
  h\in H_{R'}(S),
  \qquad
  h\notin H_{R'\cap R''}(S),
\]
such that
\[
  \operatorname{Ad}(h)\mathfrak t\subset\mathfrak g_{R''}.
\]
We may write
\[
  h=t\exp(x_1X_1)\cdots\exp(x_NX_N).
\]
Since \(h\notin H_{R'\cap R''}(S)\), there is a least \(m\) such
that
\[
  t\exp(x_1X_1)\cdots\exp(x_{m-1}X_{m-1})
  \in H_{R'\cap R''}(S),
  \qquad
  \alpha_m\notin R'',
  \qquad
  x_m\neq0.
\]
Then
\[
  h'=\exp(x_mX_m)\cdots\exp(x_NX_N)
\]
still satisfies the conditions imposed above on \(h\). By \SGARefLink{sga3:xxii:lem:5.6.2}{Lemma
5.6.2}, for every \(H\in\Gamma(S,\mathfrak t)\), the component of
\(\operatorname{Ad}(h')H\) in \(\mathfrak g^{\alpha_m}\) is
\[
  -\overline{\alpha_m}(H)x_mX_m.
\]
\footnote{Translator's note: In this sentence the French source prints
\(\overline{\alpha_m}(X)\); \SGARefLink{sga3:xxii:lem:5.6.2}{Lemma 5.6.2}\textup{\SGARefLink{sga3:xxii:lem:5.6.2:item:a}{(a)}} and the following
sentence require \(H\).}
It therefore follows from the hypothesis on \(h\) and from the
definition of \(m\) that
\[
  \overline{\alpha_m}(H)=0
  \qquad\text{for every }H\in\Gamma(S,\mathfrak t).
\]
This is impossible, since \(G\) is adjoint and
\(\overline{\alpha_m}\) is nonzero on every fiber.

\SGARefTarget{sga3:xxii:rem:5.6.2bis}
\SGAInlineRefTarget{sga3:xxii:remark:5-6-2-bis}\paragraph{\SGARefLink{sga3:xxii:remark:5-6-2-bis}{Remark 5.6.2 bis}.}
With the notation of \SGARefLink{sga3:xxii:lem:5.6.2}{Lemma 5.6.2}, suppose that
\(\operatorname{Ad}(u)\) induces the identity on \(\mathfrak t\) and
on \(\mathfrak g^{-\alpha_m}\). Then \(x_m=0\). Indeed,
\[
  x_m\overline{\alpha_m}=0,
  \qquad
  x_mH_{\alpha_m}=0.
\]
If \(\alpha_m\notin2M\), then \(\overline{\alpha_m}\) is nonzero on
every fiber; if \(\alpha_m\in2M\), then
\(\alpha_m^\ast\notin2M^\ast\), and \(H_{\alpha_m}\) is nonzero on
every fiber. Thus \(x_m=0\). It follows that \(u=e\) if
\(\operatorname{Ad}(u)\) acts trivially on \(\mathfrak g\).

\SGARefTarget{sga3:xxii:rem:5.6.3}
\paragraph{\SGARefLink{sga3:xxii:rem:5.6.3}{Remark 5.6.3}.}
Let \(H\) be a subgroup of type \textup{(R)} of a reductive
\(S\)-group \(G\), and suppose that its geometric fibers are
solvable. Then \(H\) is closed in \(G\), and
\[
  \underline{\operatorname{Norm}}_G(H)/H
\]
is represented by a separated finite étale \(S\)-scheme. This
follows from \SGARefLink{sga3:xxii:cor:5.3.18}{5.3.18} and either \SGARefLink{sga3:xxii:rem:3.5}{3.5} or \SGARefLink{sga3:xix:theorem:2:5}{Exposé XIX, 2.5}\textup{\SGARefLink{sga3:xix:theorem:2:5:item:ii}{(ii)}}.

\SGARefTarget{sga3:xxii:cor:5.6.4}
\paragraph{\SGARefLink{sga3:xxii:cor:5.6.4}{Corollary 5.6.4}.}
\textit{Let \((G,T,M,R)\) be a split reductive \(S\)-group. If
\(R'\subset R\) is closed and \(R'\cap(-R')=\varnothing\), then
\(R'\) is contained in a system of positive roots.}

\textit{Proof.}
By \SGARefLink{sga3:xxii:thm:5.4.7}{5.4.7}, \(R'\) is of type \textup{(R)}, so the result follows from
\SGARefLink{sga3:xxii:prop:5.6.1}{Proposition 5.6.1}.

\SGARefTarget{sga3:xxii:cor:5.6.5}
\paragraph{\SGARefLink{sga3:xxii:cor:5.6.5}{Corollary 5.6.5}.}
\textit{Under the hypotheses of \SGARefLink{sga3:xxii:prop:5.6.1}{Proposition 5.6.1}, multiplication
in \(G\) induces an isomorphism}
\[
  \prod_{\alpha\in R'}U_\alpha
  \xrightarrow{\ \sim\ } U_{R'},
\]
\textit{where \(U_{R'}\) is a closed subgroup scheme of \(G\),
smooth over \(S\), with connected unipotent geometric fibers, and is
independent of the chosen ordering of \(R'\). Moreover, \(H_{R'}\)
is the semidirect product \(T\cdot U_{R'}\), with \(U_{R'}\)
invariant in \(H_{R'}\).}

\textit{Proof.}
If \(R'\subset R_+\), then
\[
  U_{R'}=H_{R'}\cap U_{\geq1},
\]
with the notation of \SGARefLink{sga3:xxii:lem:5.5.6}{5.5.6}, is a closed subgroup of \(G\), of finite
presentation, and invariant in \(H_{R'}\). By \SGARefLink{sga3:xxii:prop:5.6.1}{Proposition
5.6.1}\textup{\SGARefLink{sga3:xxii:prop:5.6.1:item:iv}{(iv)}}, one has \(H_{R'}=T\cdot U_{R'}\), from which all
the assertions follow.

\SGARefTarget{sga3:xxii:rem:5.6.6}
\paragraph{\SGARefLink{sga3:xxii:rem:5.6.6}{Remark 5.6.6}.}
In particular, \(U_{R_+}\) is the group \(U_{\geq1}\) of \SGARefLink{sga3:xxii:lem:5.5.6}{5.5.6}.

Let us draw some corollaries from the preceding results concerning
groups of type \(U_{R'}\).

\SGARefTarget{sga3:xxii:cor:5.6.7}
\paragraph{\SGARefLink{sga3:xxii:cor:5.6.7}{Corollary 5.6.7}.}
\textit{Let \((G,T,M,R)\) be a split reductive \(S\)-group, and let
\(R'\) and \(R''\) be two subsets of \(R\) of type \textup{(R)},
with \(R'\cap(-R')=\varnothing\).}
\begin{enumerate}
\item[\textup{\SGAInlineRefTarget{sga3:xxii:cor:5.6.7:item:i}(i)}]
\textit{One has}
\[
  U_{R'}\cap
  \underline{\operatorname{Norm}}_G(H_{R''})
  =
  U_{R'\cap R''}.
\]

\item[\textup{\SGAInlineRefTarget{sga3:xxii:cor:5.6.7:item:ii}(ii)}]
\textit{Suppose that \(R'\) is closed. If the conditions
\(\alpha\in R'\), \(\beta\in R''\), and
\(\alpha+\beta\in R\) imply that \(\alpha+\beta\in R'\), then
\(H_{R''}\) normalizes \(U_{R'}\).}
\end{enumerate}

\textit{Proof.}
Assertion \textup{\SGARefLink{sga3:xxii:cor:5.6.7:item:i}{(i)}} follows at once from \SGARefLink{sga3:xxii:cor:5.6.5}{5.6.5} and
\SGARefLink{sga3:xxii:prop:5.6.1}{5.6.1}\textup{\SGARefLink{sga3:xxii:prop:5.6.1:item:vi}{(vi)}}. To prove \textup{\SGARefLink{sga3:xxii:cor:5.6.7:item:ii}{(ii)}}, it is enough, by \SGARefLink{sga3:xxii:lem:5.4.4}{5.4.4},
to show that \(T\) and each \(U_\beta\), for \(\beta\in R''\),
normalize \(U_{R'}\).\footnote{Editor's note: Indeed, since \(G\)
has connected fibers, it is separated over \(S\) by Exposé
VI\textsubscript{B}, \SGARefLink{cor:sga3-VIB-5.5}{5.5}; hence, by \SGARefLink{sga3:xi:corollary:6:11}{Exposé XI, 6.11} \textup{(see
also Addendum \SGARefLink{cor:sga3-VIB-6.5.5}{6.5.5} in Exposé VI\textsubscript{B})},
\(\underline{\operatorname{Norm}}_G(U_{R'})\) is represented by a
closed subgroup scheme \(N\) of \(G\). If \(N\) contains \(T\) and
the \(U_\beta\), for \(\beta\in R''\), it contains the big cell
\(\Omega_{R_+,R''}\). This cell is schematically dense in
\(H_{R''}\) by 5.4.4 and EGA IV\textsubscript{3}, 11.10.10
\textup{(the fibers of \(H_{R''}\) are integral, and
\(\Omega_{R_+,R''}\) contains a nonempty open subset of every
fiber)}. It follows that \(H_{R''}\subset N\), and therefore that
\(H_{R''}\) normalizes \(U_{R'}\).}
For \(T\) this is trivial; for \(U_\beta\) it follows from \SGARefLink{sga3:xxii:cor:5.5.2}{5.5.2} and
\SGARefLink{sga3:XXI:3.1.2}{Exposé XXI, 3.1.2}.

\SGARefTarget{sga3:xxii:cor:5.6.8}
\paragraph{\SGARefLink{sga3:xxii:cor:5.6.8}{Corollary 5.6.8}.}
\textit{Let \((G,T,M,R)\) be a split \(S\)-group, let \(R_+\) be a
system of positive roots, and let \(\alpha\) be a simple root of
\(R_+\) \textup{(that is, an element of \(R_+\) such that
\(R_+\setminus\{\alpha\}\) is closed)}. Set}
\[
  U_{\widehat\alpha}=U_{R_+\setminus\{\alpha\}}.
\]
\textit{Then:}
\begin{enumerate}
\item[\textup{\SGAInlineRefTarget{sga3:xxii:cor:5.6.8:item:i}(i)}]
\textit{\(U_{\widehat\alpha}\) is a normal subgroup of \(B_{R_+}\);}

\item[\textup{\SGAInlineRefTarget{sga3:xxii:cor:5.6.8:item:ii}(ii)}]
\textit{\(U_{R_+}\) is the semidirect product of
\(U_{\widehat\alpha}\) by \(U_\alpha\);}

\item[\textup{\SGAInlineRefTarget{sga3:xxii:cor:5.6.8:item:iii}(iii)}]
\textit{\(U_{-\alpha}\) normalizes \(U_{\widehat\alpha}\);}

\item[\textup{\SGAInlineRefTarget{sga3:xxii:cor:5.6.8:item:iv}(iv)}]
\textit{\(Z_\alpha\) normalizes \(U_{\widehat\alpha}\).}
\end{enumerate}
\textit{If one similarly defines}
\[
  U_{-\widehat\alpha}
  =
  U_{R_-\setminus\{-\alpha\}},
  \qquad
  R_-=-R_+,
\]
\textit{then}
\[
  \Omega_{R_+}
  =
  U_{-\widehat\alpha}\cdot U_{-\alpha}\cdot T\cdot
  U_\alpha\cdot U_{\widehat\alpha}.
\]

\textit{Proof.}
Assertion \textup{\SGARefLink{sga3:xxii:cor:5.6.8:item:ii}{(ii)}} follows from \SGARefLink{sga3:xxii:cor:5.6.5}{5.6.5}, and \textup{\SGARefLink{sga3:xxii:cor:5.6.8:item:i}{(i)}} from
\SGARefLink{sga3:xxii:cor:5.6.7}{5.6.7}\textup{\SGARefLink{sga3:xxii:cor:5.6.8:item:ii}{(ii)}}. Likewise, \textup{\SGARefLink{sga3:xxii:cor:5.6.8:item:iii}{(iii)}} follows from \SGARefLink{sga3:xxii:cor:5.5.2}{5.5.2}:
indeed, if \(\beta\in R_+\), with \(\beta\neq\alpha\), no
combination
\[
  i(-\alpha)+j\beta,
  \qquad i,j>0,
\]
can be negative, since \(\beta\) contains at least one simple root
different from \(\alpha\). Assertion \textup{\SGARefLink{sga3:xxii:cor:5.6.8:item:iv}{(iv)}} then follows from
\textup{\SGARefLink{sga3:xxii:cor:5.6.8:item:i}{(i)}} and \textup{\SGARefLink{sga3:xxii:cor:5.6.8:item:iii}{(iii)}}, because
\[
  U_{-\alpha}\cdot T\cdot U_\alpha
\]
is schematically dense in \(Z_\alpha\). Finally, the last assertion
follows from \textup{\SGARefLink{sga3:xxii:cor:5.6.8:item:ii}{(ii)}} and its analogue for \(U_{R_-}\).

Let us return to the general situation.

\SGARefTarget{sga3:xxii:prop:5.6.9}
\paragraph{\SGARefLink{sga3:xxii:prop:5.6.9}{Proposition 5.6.9}.}
\textit{Let \(S\) be a scheme, let \(G\) be a reductive \(S\)-group,
and let \(H\) be a subgroup of type \textup{(R)} with solvable
geometric fibers.}
\begin{enumerate}
\item[\textup{\SGAInlineRefTarget{sga3:xxii:prop:5.6.9:item:i}(i)}]
\textit{The group}
\[
  \mathrm D_S(H)
  =
  \underline{\operatorname{Hom}}_{S\text{-gr.}}
    (H,\mathbb G_{m,S})
\]
\textit{is represented by a twisted constant \(S\)-group whose type
at \(s\in S\) is}
\[
  \mathbb Z^{\operatorname{rgred}(G_s)}.
\]
\textit{The biduality morphism \textup{(Exposé VIII, \S1)}}
\[
  f:H\longrightarrow\mathrm D_S\bigl(\mathrm D_S(H)\bigr)
\]
\textit{is smooth and surjective.}

\item[\textup{\SGAInlineRefTarget{sga3:xxii:prop:5.6.9:item:ii}(ii)}]
\textit{The kernel \(H^u\) of \(f\) is the largest closed normal
subgroup scheme of \(H\), smooth over \(S\), with connected
unipotent geometric fibers. It is called the unipotent part of
\(H\), and is also denoted}
\[
  H^u=\operatorname{rad}^u(H).
\]
\textit{Moreover, \(H^u\) is the sheaf of commutators of \(H\):
every morphism of groups from \(H\) to an \(S\)-presheaf of
commutative groups that is separated for the fppf topology vanishes
on \(H^u\), and therefore factors through}
\[
  H/H^u=\mathrm D_S\bigl(\mathrm D_S(H)\bigr).
\]

\item[\textup{\SGAInlineRefTarget{sga3:xxii:prop:5.6.9:item:iii}(iii)}]
\textit{If \(T\) is a maximal torus of \(H\), the morphism
\(T\longrightarrow H\) induces isomorphisms}
\[
  \mathrm D_S(H)\xrightarrow{\ \sim\ }\mathrm D_S(T),
  \qquad
  T\xrightarrow{\ \sim\ }\mathrm D_S\bigl(\mathrm D_S(H)\bigr).
\]
\textit{Moreover, \(H\) is identified with the semidirect product
of \(H^u\) by \(T\).}

\item[\textup{\SGAInlineRefTarget{sga3:xxii:prop:5.6.9:item:iv}(iv)}]
\textit{In the situation of \SGARefLink{sga3:xxii:prop:5.6.1}{5.6.1}, if \(H=H_{R'}\), then}
\[
  H^u=U_{R'}.
\]
\end{enumerate}

\textit{Proof.}
The assertions of the proposition are local for the étale topology
\textup{(\SGARefLink{sga3:x:corollary:5-5}{Exposé X, 5.5})}. We may therefore reduce to the situation
of \SGARefLink{sga3:xxii:prop:5.6.1}{5.6.1}. We then know from \SGARefLink{sga3:xxii:cor:5.6.5}{5.6.5} that \(H_{R'}\) is the semidirect
product of \(U_{R'}\) by \(T\). Let us show that \(U_{R'}\) is the
sheaf of commutators of \(H_{R'}\). Since
\(H_{R'}/U_{R'}=T\) is commutative, it suffices to prove that every
morphism of groups
\[
  \phi:H_{R'}\longrightarrow V
\]
as in \textup{\SGARefLink{sga3:xxii:prop:5.6.9:item:ii}{(ii)}} vanishes on \(U_{R'}\). It is enough to prove
that \(\phi\) vanishes on each \(U_\alpha\), for \(\alpha\in R'\).
If \(t\in T(S')\) and \(X\in W(\mathfrak g^\alpha)(S')\), then
\[
  1
  =
  \phi\bigl(t\exp_\alpha(X)t^{-1}\exp_\alpha(-X)\bigr)
  =
  \phi\bigl(\exp_\alpha((\alpha(t)-1)X)\bigr).
\]
Since \(\alpha:T\longrightarrow\mathbb G_{m,S}\) is faithfully
flat, it follows at once that \(\phi\) vanishes on
\(U_\alpha^\times\). But every section of \(U_\alpha\) is locally
the sum of two sections of \(U_\alpha^\times\). Thus
\[
  \underline{\operatorname{Hom}}_{S\text{-gr.}}(H,V)
  =
  \underline{\operatorname{Hom}}_{S\text{-gr.}}
    (H/U_{R'},V)
\]
for every \(V\) as above. Applying this result to
\(V=\mathbb G_{m,S}\), we immediately obtain \textup{\SGARefLink{sga3:xxii:prop:5.6.9:item:i}{(i)}} and
\textup{\SGARefLink{sga3:xxii:prop:5.6.9:item:iii}{(iii)}}, then \textup{\SGARefLink{sga3:xxii:prop:5.6.9:item:iv}{(iv)}} and the second assertion of
\textup{\SGARefLink{sga3:xxii:prop:5.6.9:item:ii}{(ii)}}.

It remains to prove the first assertion of \textup{\SGARefLink{sga3:xxii:prop:5.6.9:item:ii}{(ii)}}. The only
nontrivial point is that every closed normal subgroup \(U\) of
\(H\), smooth over \(S\), with connected unipotent geometric
fibers, is a subgroup of \(H^u\). We first have:

\SGARefTarget{sga3:xxii:lem:5.6.10}
\paragraph{\SGARefLink{sga3:xxii:lem:5.6.10}{Lemma 5.6.10}.}
\textit{Let \(G\) be a reductive \(S\)-group, \(T\) a maximal
torus, and \(U\) a subgroup scheme of \(G\), smooth over \(S\),
with unipotent geometric fibers, and normalized by \(T\). Then}
\[
  U\cap T=e.
\]

\textit{Proof.}
Indeed, since
\[
  T=\underline{\operatorname{Centr}}_G(T),
\]
one has \(U\cap T=U^T\), where the right-hand side denotes the
invariants under \(\operatorname{int}(T)\). By \SGARefLink{sga3:xix:lemma:1-4}{Exposé XIX, 1.4},
\(U\cap T\) is therefore smooth over \(S\). It is also radicial
over \(S\): for every \(s\in S\), the group
\[
  U(\overline s)\cap T(\overline s)
\]
consists of elements that are both unipotent and semisimple. This
proves the lemma.

We return to the proof of \SGARefLink{sga3:xxii:prop:5.6.9}{5.6.9}\textup{\SGARefLink{sga3:xxii:prop:5.6.9:item:ii}{(ii)}}. If \(U\) is a normal
subgroup of \(H\) as above, then the semidirect product \(T\cdot U\)
is a subgroup of \(G\) of type \textup{(R)}, with solvable
geometric fibers. We may therefore suppose it to be of the form
\(H_{R''}\), with \(R''\subset R'\). It is enough to prove that
\(U=U_{R''}\), and we are thus reduced to the case
\(H=T\cdot U\). Since \(H/U\) is commutative, the French text says
that \(U\) is a subsheaf of the sheaf of commutators of \(H\),
which is \(H^u\).\footnote{Translator's note: The implication in the
French is stated in the wrong direction: commutativity of \(H/U\)
gives \(H^u\subset U\). The desired inclusion follows because
\(H/H^u\) is a torus whereas \(U\) has connected unipotent
geometric fibers.}
\hfill\(\square\)

Let us note that we have in fact just proved:

\SGARefTarget{sga3:xxii:prop:5.6.11}
\paragraph{\SGARefLink{sga3:xxii:prop:5.6.11}{Proposition 5.6.11}.}
\textit{Let \(S\) be a scheme, \(G\) a reductive \(S\)-group, and
\(T\) a maximal torus of \(G\). The maps}
\[
  H\longmapsto H^u,
  \qquad
  U\longmapsto T\cdot U
\]
\textit{are mutually inverse bijections between the set of subgroups
\(H\) of \(G\) of type \textup{(R)}, containing \(T\), and having
solvable geometric fibers, and the set of subgroups \(U\) of \(G\),
smooth over \(S\), normalized by \(T\), and having connected
unipotent geometric fibers.}

In particular, when \((G,T,M,R)\) is split, the groups \(H_{R'}\)
and \(U_{R'}\) correspond.

\SGARefTarget{sga3:xxii:cor:5.6.12}
\paragraph{\SGARefLink{sga3:xxii:cor:5.6.12}{Corollary 5.6.12}.}
\textit{Let \(S\) be a scheme, let \((G,T,M,R)\) be a split
\(S\)-group \textup{(respectively, and let \(R_+\) be a system of
positive roots of \(R\) defining the Borel subgroup \(B\))}.
Every smooth group subscheme of \(G\) with connected unipotent
geometric fibers \textup{(respectively, every smooth group subscheme
of \(B^u\))}, normalized by \(T\), is locally on \(S\) of the form
\(U_{R'}\), where \(R'\) is a subset of \(R\) contained in a system
of positive roots \textup{(respectively, a subset of \(R_+\))} and
is of type \textup{(R)}.}

\textit{Proof.}
For the ``respectively'' case, it is enough to observe that the
geometric fibers of the given group are unipotent and connected by
the \emph{Bible}, \S13.2, Theorem 1\textup{\SGARefLink{sga3:xxii:rem:5.1.2:item:d}{(d)}}.
\hfill\(\square\)

\SGARefLink{sga3:xxii:prop:5.6.9}{Proposition 5.6.9} also has the following corollary:

\SGARefTarget{sga3:xxii:cor:5.6.13}
\paragraph{\SGARefLink{sga3:xxii:cor:5.6.13}{Corollary 5.6.13}.}
\textit{Let \(S\) be a scheme, \(G\) a reductive \(S\)-group, and
\(H\) a subgroup of type \textup{(R)} with solvable geometric
fibers. Let \(\underline{\operatorname{Tor}}(H)\) be the functor of
maximal tori of \(H\):}
\[
  \underline{\operatorname{Tor}}(H)(S')
  =
  \{\,\text{maximal tori of }H_{S'}\,\}.
\]
\textit{Then \(\underline{\operatorname{Tor}}(H)\) is represented by
an affine smooth \(S\)-scheme, which is a principal homogeneous
bundle under \(H^u\) for the action}
\[
  (h,T)\longmapsto\operatorname{int}(h)T.
\]

\textit{Proof.}
Indeed, if \(T\) and \(T'\) are two maximal tori of \(H_{S'}\), there
exists a unique section \(h\in H^u(S')\) such that
\(\operatorname{int}(h)T=T'\). The uniqueness of \(h\) follows at
once from the equality
\[
  \underline{\operatorname{Norm}}_G(T)\cap H^u=e
\]
\textup{(cf., for example, \SGARefLink{sga3:xxii:prop:5.6.1}{5.6.1})}; it is therefore enough to prove
the existence of \(h\) locally for the étale topology. By \SGARefLink{sga3:xxii:prop:5.2.6}{5.2.6} and
\SGARefLink{sga3:xxii:rem:5.1.2}{5.1.2}\textup{\SGARefLink{sga3:xxii:rem:5.1.2:item:a}{(a)}}, we may suppose that \(T\) and \(T'\) are conjugate
by a section of \(H\), whence the desired conclusion because
\[
  H=H^u\cdot T
\]
by \SGARefLink{sga3:xxii:prop:5.6.9}{5.6.9}\textup{\SGARefLink{sga3:xxii:prop:5.6.9:item:iii}{(iii)}}. It follows that
\(\underline{\operatorname{Tor}}(H)\) is a principal homogeneous
sheaf under \(H^u\), which is affine and smooth over \(S\); the
assertion follows at once.
\hfill\(\square\)

\SGARefTarget{sga3:xxii:sec:5.7}
\subsection*{\SGAInlineRefTarget{sga3:xxii:subsection:5-7}5.7. Bruhat's Theorem}
\addcontentsline{toc}{subsection}{\SGARefLink{sga3:xxii:subsection:5-7}{5.7}. Bruhat's Theorem}

\SGARefTarget{sga3:xxii:recall:5.7.1}
\paragraph{Recall \SGARefLink{sga3:xxii:recall:5.7.1}{5.7.1}.}
\textit{Let \(k\) be an algebraically closed field, \(G\) a
reductive \(k\)-group, \(B\) a Borel subgroup of \(G\), \(T\) a
maximal torus of \(B\), and
\(N=\underline{\operatorname{Norm}}_G(T)\). Then}
\[
  G(k)=B(k)N(k)B(k);
\]
\textit{this is Bruhat's theorem \textup{(\emph{Bible}, \S13.4,
Corollary 1 to Theorem 3)}. More precisely, with the notation of
\SGARefLink{sga3:xxii:par:3.6}{3.6}, the sets}
\[
  B(k)N_w(k)B(k)=B^u(k)N_w(k)B^u(k)
\]
\textit{form a partition of \(G(k)\) as \(w\) ranges over
\((N/T)(k)\). If \(B'\) is another Borel subgroup of \(G\)
containing \(T\), the sets \(B'(k)N_w(k)B(k)\) likewise form a
partition of \(G(k)\). Indeed, if \(y\in N(k)\) is such that
\(\operatorname{int}(y)B=B'\), then}
\[
  yB(k)N_w(k)B(k)=B'(k)N_{yw}(k)B(k).
\]

\SGARefTarget{sga3:xxii:def:5.7.2}
\paragraph{\SGARefLink{sga3:xxii:def:5.7.2}{Definition 5.7.2}.}
\textit{Let \((G,T,M,R)\) be a split \(S\)-group, let \(R_-\) be a
system of positive roots of \(R\), and let \(B'=B_{R_-}\) be the
Borel subgroup it defines. For \(w\in W\), put
\textup{(cf. \SGARefLink{sga3:xxii:cor:5.6.5}{5.6.5})}}
\[
  R_-^w=R_-\cap w(R_-),
  \qquad
  B'{}^u_w=U_{R_-^w}
  =\prod_{\alpha\in R_-^w}U_\alpha.
\]
\textit{If
\(n_w\in\underline{\operatorname{Norm}}_G(T)(S)\) is a
representative of \(w\) \textup{(\SGARefLink{sga3:xxii:cor:3.8}{3.8})}, one may also write}
\[
  B'{}^u_w=B'{}^u\cap\operatorname{int}(n_w)B'{}^u.
\]

\SGARefTarget{sga3:xxii:lem:5.7.3}
\paragraph{\SGARefLink{sga3:xxii:lem:5.7.3}{Lemma 5.7.3}.}
\textit{Let \((G,T,M,R)\) be a split \(S\)-group, let \(R_+\) be a
system of positive roots of \(R\), let \(R_-=-R_+\), and let \(B\)
\textup{(respectively \(B'\))} be the Borel subgroup of \(G\)
defined by \(R_+\) \textup{(respectively \(R_-\))}. Let \(w\in W\),
and let \(N_w\) and \(B'{}^u_w\) be the corresponding subschemes of
\(G\) \textup{(\SGARefLink{sga3:xxii:cor:3.8}{3.8} and \SGARefLink{sga3:xxii:def:5.7.2}{5.7.2})}.}
\begin{enumerate}
\item[\textup{\SGAInlineRefTarget{sga3:xxii:lem:5.7.3:item:i}(i)}]
\textit{The sheaf \(B'\cdot N_w\cdot B\), the image of the
morphism}
\[
  B'\mathop{\times}\limits_S N_w
  \mathop{\times}\limits_S B
  \longrightarrow G
\]
\textit{induced by multiplication in \(G\), is represented by a
subscheme of \(G\) \textup{(in fact, by a closed subscheme of the
open subscheme \(n_w\Omega_{R_+}\))}.}

\item[\textup{\SGAInlineRefTarget{sga3:xxii:lem:5.7.3:item:ii}(ii)}]
\textit{The morphism}
\[
  B'{}^u_w\mathop{\times}\limits_S N_w
  \mathop{\times}\limits_S B^u
  \longrightarrow G,
\]
\textit{induced by multiplication in \(G\), is an immersion whose
image is the preceding subscheme.}
\end{enumerate}

\textit{Proof.}
Let us first show that the morphism in \textup{\SGARefLink{sga3:xxii:lem:5.7.3:item:ii}{(ii)}} is an
immersion. By definition, \(\operatorname{int}(n_w)^{-1}\) induces
a closed immersion of \(B'{}^u_w\) into \(B'{}^u\); hence the
morphism
\[
  (u,b)\longmapsto n_w^{-1}u\,n_wb
\]
induces a closed immersion
\[
  B'{}^u_w\mathop{\times}\limits_S B
  \longrightarrow \Omega_{R_+}.
\]
It follows immediately that the morphism in \textup{\SGARefLink{sga3:xxii:lem:5.7.3:item:ii}{(ii)}} induces a
closed immersion of its source into the open subscheme
\(n_w\Omega_{R_+}\). To prove \textup{\SGARefLink{sga3:xxii:lem:5.7.3:item:i}{(i)}}, it is enough to see that
\[
  B'(S)N_w(S)B(S)
  =
  B'{}^u_w(S)N_w(S)B^u(S).
\]
Now, if \(\alpha\in R\), then
\(\operatorname{int}(n_w)U_\alpha(S)=U_{w(\alpha)}(S)\). Thus, if
\(w^{-1}(\alpha)\in R_+\),
\[
\begin{aligned}
  U_\alpha(S)N_w(S)B(S)
  &=U_\alpha(S)n_wT(S)B^u(S)\\
  &=n_wU_{w^{-1}(\alpha)}(S)T(S)B^u(S)\\
  &=n_wB(S)=N_w(S)B^u(S).
\end{aligned}
\]
In view of the definition of \(R_-^w\), this immediately gives the
desired assertion.
\hfill\(\square\)

\SGARefTarget{sga3:xxii:thm:5.7.4}
\paragraph{\SGARefLink{sga3:xxii:thm:5.7.4}{Theorem 5.7.4}.}
\textit{Let \(S\) be a scheme, let \((G,T,M,R)\) be a split
\(S\)-group, let \(B\) be the Borel subgroup defined by the system
of positive roots \(R_+\), and let \(B'\) be the Borel subgroup
defined by \(R_-=-R_+\).}
\begin{enumerate}
\item[\textup{(i)}]
\textit{\textup{(Bruhat's theorem)} The schemes
\(B'{}^u_w\cdot N_w\cdot B\), as \(w\) ranges over \(W\), form a
partition of the underlying set of \(G\).}

\item[\textup{\SGAInlineRefTarget{sga3:xxii:thm:5.7.4:item:ii}(ii)}]
\textit{For each \(w\in W\), let \(n_w\) be a representative of
\(w\) in \(\underline{\operatorname{Norm}}_G(T)(S)\)
\textup{(\SGARefLink{sga3:xxii:cor:3.8}{3.8})}. Then the open subschemes}
\[
  n_w\Omega=n_wB'{}^u\cdot T\cdot B^u
\]
\textit{as \(w\) ranges over \(W\), form an open covering of \(G\).}
\end{enumerate}

\textit{Proof.}
Both assertions may be checked on the geometric fibers, where they
follow from \SGARefLink{sga3:xxii:recall:5.7.1}{5.7.1} and \SGARefLink{sga3:xxii:lem:5.7.3}{5.7.3}.
\hfill\(\square\)

\SGARefTarget{sga3:xxii:rem:5.7.5}
\paragraph{\SGARefLink{sga3:xxii:rem:5.7.5}{Remark 5.7.5}.}
Assertion \textup{\SGARefLink{sga3:xxii:lem:5.7.3:item:i}{(i)}} implies that if \(S\) is the spectrum of a
field, then \(G(S)\) is the disjoint union of the sets
\[
  B'{}^u_w(S)\cdot T(S)\cdot B^u(S).
\]
The corresponding assertion for an arbitrary \(S\), even a local
or Artinian one, is plainly false. Note, however, that
\textup{\SGARefLink{sga3:xxii:thm:5.7.4:item:ii}{(ii)}} implies that if \(S\) is local, then \(G(S)\) is
indeed the union of the \(n_w\Omega(S)\). In fact:

\SGARefTarget{sga3:xxii:cor:5.7.6}
\paragraph{\SGARefLink{sga3:xxii:cor:5.7.6}{Corollary 5.7.6}.}
\textit{Let \(\Delta\) be a system of simple roots of the split
group \(G\) over the local scheme \(S\).}
\begin{enumerate}
\item[\textup{\SGAInlineRefTarget{sga3:xxii:cor:5.7.6:item:i}(i)}]
\textit{Then \(G(S)\) is generated by \(T(S)\) and the groups
\(U_\alpha(S)\), for \(\alpha\in\Delta\cup-\Delta\).}

\item[\textup{\SGAInlineRefTarget{sga3:xxii:cor:5.7.6:item:ii}(ii)}]
\textit{If \(G\) is simply connected \textup{(\SGARefLink{sga3:xxii:def:4.3.3}{4.3.3})}, then
\(G(S)\) is already generated by the groups \(U_\alpha(S)\), for
\(\alpha\in\Delta\cup-\Delta\).}
\end{enumerate}

\textit{Proof.}
Let \(H\) be the subgroup of \(G(S)\) generated by the
\(U_\alpha(S)\), for \(\alpha\in\Delta\cup-\Delta\). First observe
that \(H\) contains a representative of every \(s_\alpha\), for
\(\alpha\in\Delta\), in
\(\underline{\operatorname{Norm}}_G(T)(S)\) \textup{(\SGARefLink{sga3:xx:3.1}{Exposé XX,
3.1})}, and hence a representative \(n_w\) of every \(w\in W\).
Since every \(\alpha\in R\) can be written as \(w(\alpha_0)\), with
\(w\in W\) and \(\alpha_0\in\Delta\), we have
\[
  U_\alpha(S)
  =
  \operatorname{int}(n_w)U_{\alpha_0}(S)
  \subset H.
\]
The subgroup generated by \(H\) and \(T(S)\) therefore contains
\(\Omega(S)\), and hence, by the preceding observation, is all of
\(G(S)\).

Now suppose that \(G\) is simply connected. Let us prove that
\(H\supset T(S)\). By \SGARefLink{sga3:xx:2.7}{Exposé XX, 2.7}, \(H\) contains
\(\alpha^*(\mathbb G_m(S))\) for every \(\alpha\in\Delta\), and it
remains only to apply \SGARefLink{sga3:xxii:rem:4.3.8}{4.3.8}.
\hfill\(\square\)

\SGARefTarget{sga3:xxii:rem:5.7.6.1}
\paragraph{\SGARefLink{sga3:xxii:rem:5.7.6.1}{Remark 5.7.6.1}.}
Instead of taking, for each \(\alpha\in\Delta\), both \(U_\alpha(S)\)
and \(U_{-\alpha}(S)\), it is enough to take \(U_\alpha(S)\) together
with a representative \(w_\alpha\) of the reflection \(s_\alpha\), or
else \(U_\alpha(S)\) together with a section of
\(U_{-\alpha}^{\times}\), \ldots.

\SGARefTarget{sga3:xxii:cor:5.7.7}
\paragraph{\SGARefLink{sga3:xxii:cor:5.7.7}{Corollary 5.7.7}.}
\textit{Suppose that \(G\) has semisimple rank \(1\), and choose
\(u_\alpha\in U_\alpha^\times(S)\). Then \(\Omega\) and
\(u_\alpha\Omega\) form an open covering of \(G\).}

\textit{Proof.}
If \(u_{-\alpha}\) is the section of \(U_{-\alpha}\) paired with
\(u_\alpha\) \textup{(cf.\ \SGARefLink{sga3:xxii:not:1.3}{1.3})}, then \SGARefLink{sga3:xxii:thm:5.7.4}{5.7.4} \textup{\SGARefLink{sga3:xxii:thm:5.7.4:item:ii}{(ii)}} gives
\[
  G
  =
  \Omega\cup
  u_{-\alpha}^{-1}u_\alpha u_{-\alpha}^{-1}\Omega,
\]
and hence
\[
  G
  =
  u_{-\alpha}G
  =
  u_{-\alpha}\Omega\cup
  u_\alpha u_{-\alpha}^{-1}\Omega
  =
  \Omega\cup u_\alpha\Omega.
\]
\hfill\(\square\)

\SGARefTarget{sga3:xxii:cor:5.7.8}
\paragraph{\SGARefLink{sga3:xxii:cor:5.7.8}{Corollary 5.7.8}.}
\textit{Let \(S\) be a scheme and let \(G\) be a reductive
\(S\)-group. Then \(G\) is essentially free over \(S\) in the sense
of \SGARefLink{sga3:viii:definition:6:1}{Exposé VIII, 6.1}.}

\textit{Proof.}
The assertion is local for the fpqc topology, so we may assume that
\(G\) is split. It then admits an open covering by open subschemes
isomorphic to
\[
  \mathbb G_{a,S}^{N}\mathop{\times}\limits_S\mathbb G_{m,S}^{n},
\]
and these are essentially free.
\hfill\(\square\)

\SGARefTarget{sga3:xxii:lem:5.7.9}
\paragraph{\SGARefLink{sga3:xxii:lem:5.7.9}{Lemma 5.7.9}.}
\textit{Under the hypotheses of \SGARefLink{sga3:xxii:thm:5.7.4}{5.7.4}, let \(\alpha\) be a simple
root of \(R_+\), and let \(u_\alpha\in U_\alpha^\times(S)\). For
every \(v\in U_{-\alpha}(S)\), one has}
\[
  \Omega\cdot v\subset\Omega\cup u_\alpha\cdot\Omega.
\]

\textit{Proof.}
We are comparing two open subschemes of \(G\), so it is enough to do
so when \(S\) is the spectrum of a field \(k\). We must therefore
prove
\[
  \Omega(k)v\subset\Omega(k)\cup u_\alpha\Omega(k).
\]
Now
\[
\begin{aligned}
  \Omega(k)v
  &=B'{}^u(k)T(k)B^u(k)v\\
  &=U_{-\widehat\alpha}(k)U_{-\alpha}(k)T(k)
    U_\alpha(k)U_{\widehat\alpha}(k)v\\
  &\subset
    U_{-\widehat\alpha}(k)Z_\alpha(k)
    U_{\widehat\alpha}(k)v.
\end{aligned}
\]
Here we use the decomposition of \SGARefLink{sga3:xxii:cor:5.6.8}{5.6.8}. Applying \SGARefLink{sga3:xxii:cor:5.6.8}{5.6.8}
\textup{\SGARefLink{sga3:xxii:cor:5.6.8:item:iii}{(iii)}} and then \SGARefLink{sga3:xxii:cor:5.7.7}{5.7.7} to the group \(Z_\alpha\), we obtain
\[
\begin{aligned}
  \Omega(k)v
  &\subset
    U_{-\widehat\alpha}(k)Z_\alpha(k)v
    U_{\widehat\alpha}(k)\\
  &\subset
    U_{-\widehat\alpha}(k)Z_\alpha(k)
    U_{\widehat\alpha}(k)\\
  &\subset
    U_{-\widehat\alpha}(k)U_{-\alpha}(k)T(k)
    U_\alpha(k)U_{\widehat\alpha}(k)\\
  &\qquad\cup
    U_{-\widehat\alpha}(k)u_\alpha U_{-\alpha}(k)T(k)
    U_\alpha(k)U_{\widehat\alpha}(k).
\end{aligned}
\]
Applying \SGARefLink{sga3:xxii:cor:5.6.8}{5.6.8} \textup{\SGARefLink{sga3:xxii:cor:5.6.8:item:iii}{(iii)}} once more, now with \(R_-\) in place
of \(R_+\), gives the result.
\hfill\(\square\)

\SGARefTarget{sga3:xxii:prop:5.7.10}
\paragraph{\SGARefLink{sga3:xxii:prop:5.7.10}{Proposition 5.7.10}.}
\textit{Under the hypotheses of \SGARefLink{sga3:xxii:thm:5.7.4}{5.7.4}, choose, for each simple
root \(\alpha\), an element
\(u_\alpha\in U_\alpha^\times(S)\). Let \(U_1\) be the submonoid of
\(B^u(S)\) generated by the \(u_\alpha\). Then the open subschemes
\(u\Omega\), for \(u\in U_1\), form an open covering of \(G\).}

\textit{Proof.}
Once again, we may assume that \(S\) is the spectrum of a field
\(k\). In view of \SGARefLink{sga3:xxii:cor:5.7.6}{5.7.6}, it is enough to prove that
\[
  \bigcup_{u\in U_1}u\Omega(k)
\]
is stable under right multiplication by \(T(k)\), \(U_\alpha(k)\),
and \(U_{-\alpha}(k)\), for every simple root \(\alpha\). This is
trivial in the first two cases, and in the last it follows from the
lemma.
\hfill\(\square\)

\SGARefTarget{sga3:xxii:rem:5.7.11}
\paragraph{\SGARefLink{sga3:xxii:rem:5.7.11}{Remark 5.7.11}.}
Let us point out a special case of \SGARefLink{sga3:xxii:def:5.7.2}{5.7.2}. If \(w=s_\alpha\) is the
reflection with respect to the simple root \(\alpha\), then
\[
  R_-\cap s_\alpha(R_-)=R_-\setminus\{-\alpha\}
\]
\textup{(\SGARefLink{sga3:XXI:3.3.1}{Exposé XXI, 3.3.1})}, and hence, in the notation of \SGARefLink{sga3:xxii:cor:5.6.8}{5.6.8},
\[
  B'{}^u_{s_\alpha}=U_{-\widehat\alpha}.
\]

\SGARefTarget{sga3:xxii:rem:5.7.12}
\paragraph{\SGARefLink{sga3:xxii:rem:5.7.12}{Remark 5.7.12}.}
In fact, the proof of \SGARefLink{sga3:xxii:prop:5.7.10}{5.7.10} immediately gives the following
statement. Under the hypotheses of \SGARefLink{sga3:xxii:prop:5.7.10}{5.7.10}, let \(\Gamma\) be a
submonoid of \(G(S)\). The open subschemes \(g\Omega\), for
\(g\in\Gamma\), form an open covering of \(G\) if and only if, for
every \(s\in S\) and every simple root \(\alpha\), one has
\[
  (u_\alpha)_{\bar{s}}B'{}^u(\bar{s})
  \subset
  \Gamma\cdot B'{}^u(\bar{s})\cdot T(\bar{s})\cdot B^u(\bar{s}).
\]

\SGARefTarget{sga3:xxii:rem:5.7.13}
\paragraph{\SGARefLink{sga3:xxii:rem:5.7.13}{Remark 5.7.13}.}
By \SGARefLink{sga3:xxii:cor:5.5.5}{5.5.5} \textup{\SGARefLink{sga3:xxii:cor:5.5.5:item:iii}{(iii)}}, arguing as in \SGARefLink{sga3:xxii:recall:5.7.1}{5.7.1}, one immediately
obtains the following variant of \SGARefLink{sga3:xxii:thm:5.7.4}{5.7.4}. Let \((G,T,M,R)\) be a
split \(S\)-group, and let \(B\) and \(B'\) be two Borel subgroups
of \(G\) containing \(T\). For every \(w\in W\), the sheaf
\(B'\cdot N_w\cdot B\) is representable by a subscheme of \(G\);
these subschemes, as \(w\) ranges over \(W\), form a partition of
the underlying set of \(G\). One can also give the analogue of
\SGARefLink{sga3:xxii:lem:5.7.3}{5.7.3} \textup{\SGARefLink{sga3:xxii:lem:5.7.3:item:ii}{(ii)}}: one must put
\[
  B'{}^u_w
  =
  B'{}^u\cap\operatorname{int}(n_w)\widetilde B{}^u,
\]
where \(\widetilde B\) is the Borel subgroup opposite to \(B\) with
respect to \(T\) \textup{(cf.\ \SGARefLink{sga3:xxii:prop:5.9.2}{5.9.2})}.

\SGARefTarget{sga3:xxii:prop:5.7.14}
\paragraph{\SGARefLink{sga3:xxii:prop:5.7.14}{Proposition 5.7.14}.}
\textit{Let \(S\) be a scheme, let \(G\) be a reductive
\(S\)-group, and let}
\[
  \operatorname{Ad}\colon
  G\longrightarrow GL_{\mathcal O_S}(\mathfrak g)
\]
\textit{be its adjoint representation. Then}
\[
  \operatorname{Ker}(\operatorname{Ad})
  =
  \underline{\operatorname{Centr}}(G).
\]
\textit{Equivalently, the canonical homomorphism induced by
\(\operatorname{Ad}\) upon passage to the quotient,}
\[
  \overline{\operatorname{Ad}}\colon
  G/\underline{\operatorname{Centr}}(G)
  =
  \operatorname{ad}(G)
  \longrightarrow
  GL_{\mathcal O_S}(\mathfrak g),
\]
\textit{is a monomorphism.}

\textit{Proof.}
We may assume that \(G\) is split. Choose a total order on
\(\Gamma_0(R)\) compatible with its group structure, and let
\(R_+\) be the set of positive roots. By \SGARefLink{sga3:xxii:thm:5.7.4}{5.7.4} \textup{\SGARefLink{sga3:xxii:thm:5.7.4:item:ii}{(ii)}} and
\SGARefLink{sga3:xxii:cor:4.1.6}{4.1.6}, it is enough to prove that, if \(n_w\) represents
\(w\in W\), if
\[
  u\in U(S),\qquad t\in T(S),\qquad v\in U_-(S),
\]
and if \(\operatorname{Ad}(n_wvtu)=\operatorname{id}\), then
\(w=e\), \(v=e\), and \(u=e\). For each \(m\in R\cup\{0\}\), put
\[
  \mathfrak g^{>m}
  =
  \bigoplus_{n>m}\mathfrak g^n,
  \qquad
  \mathfrak g^{<m}
  =
  \bigoplus_{n<m}\mathfrak g^n.
\]
Let \(X\in\Gamma(S,\mathfrak g^m)\); write
\[
  \operatorname{Ad}(tu)X
  =
  \operatorname{Ad}(v^{-1}n_w^{-1})X.
\]
Now
\[
\begin{aligned}
  \operatorname{Ad}(t)\operatorname{Ad}(u)X-m(t)X
  &\in\Gamma(S,\mathfrak g^{>m}),\\
  \operatorname{Ad}(v^{-1}n_w^{-1})X
  -\operatorname{Ad}(n_w^{-1})X
  &\in\Gamma\bigl(S,\mathfrak g^{<w^{-1}(m)}\bigr).
\end{aligned}
\]
If \(w\ne e\), there is an \(\alpha\in R\) such that
\(w^{-1}(\alpha)<\alpha\). On taking \(m=\alpha\), this gives a
contradiction, since
\[
\begin{aligned}
  \operatorname{Ad}(tu)X
  &\in
  \Gamma\bigl(S,\mathfrak g^\alpha+\mathfrak g^{>\alpha}\bigr)
  \cap
  \Gamma\bigl(
    S,
    \mathfrak g^{w^{-1}(\alpha)}
    +\mathfrak g^{<w^{-1}(\alpha)}
  \bigr)\\
  &=0.
\end{aligned}
\]
Thus \(w=e\), and we may choose \(n_w=e\). We then have
\[
  \operatorname{Ad}(v^{-1})X-X
  \in
  \Gamma\bigl(
    S,
    \mathfrak g^{<m}\cap
    (\mathfrak g^m+\mathfrak g^{>m})
  \bigr)
  =0,
\]
whence \(\operatorname{Ad}(v)X=X\) for every
\(X\in\Gamma(S,\mathfrak g^m)\), and therefore
\(\operatorname{Ad}(v)=\operatorname{id}\). Similarly,
\(\operatorname{Ad}(u)=\operatorname{id}\). The result now follows
from \SGARefLink{sga3:xxii:remark:5-6-2-bis}{5.6.2 bis}.
\hfill\(\square\)

\SGARefTarget{sga3:xxii:section:5.8}
\section*{5.8. Schemes Associated with a Reductive Group}
\addcontentsline{toc}{section}{\SGARefLink{sga3:xxii:section:5.8}{5.8}. Schemes Associated with a
Reductive Group}

\SGARefTarget{sga3:xxii:thm:5.8.1}
\paragraph{\SGARefLink{sga3:xxii:thm:5.8.1}{Theorem 5.8.1}.}
\textit{Let \(S\) be a scheme and let \(G\) be a reductive
\(S\)-group. Let \(\mathcal H\) be the functor of subgroups of type
\textup{(R)} of \(G\): for every \(S'\to S\),
\(\mathcal H(S')\) is the set of subgroups of type \textup{(R)} of
\(G_{S'}\) \textup{(cf.\ \SGARefLink{sga3:xxii:def:5.2.1}{5.2.1})}. Then \(\mathcal H\) is
representable by a quasi-projective \(S\)-scheme of finite
presentation over \(S\).}

\textit{Proof.}\footnote{Editor's note: The original is expanded in
what follows.}
Let
\[
  G'=G/\underline{\operatorname{Centr}}(G)
\]
be the adjoint group of \(G\) \textup{(\SGARefLink{sga3:xxii:def:4.3.6}{4.3.6})}, and let
\(u\colon G\to G'\) be the quotient morphism. By \SGARefLink{sga3:xii:7.12}{Exposé XII, 7.12},
the map
\[
  H'\longmapsto u^{-1}(H')
\]
is a bijection from the set of subgroups of type \textup{(R)} of
\(G'\) onto the set of subgroups of type \textup{(R)} of \(G\);
this remains true after every base change. Thus, replacing \(G\)
by \(G'\), we may assume that \(G\) is adjoint. Consider the
morphism
\[
  u\colon
  \mathcal H
  \longrightarrow
  \underline{\operatorname{Grass}}(\mathfrak g)
\]
that assigns to every subgroup of type \textup{(R)} its Lie algebra
\textup{(a submodule of \(\mathfrak g\) that is locally a direct
summand)}.
Then \(u\) is a monomorphism by \SGARefLink{sga3:xxii:cor:5.3.3}{5.3.3}. It is enough to prove that
it is representable by an immersion of finite presentation; in
other words, it is enough to prove the following assertion. For
every \(S_1\to S\), given a locally direct-summand submodule
\(\mathfrak h\) of \(\mathfrak g_{S_1}\), the morphisms
\(S'\to S_1\) for which \(\mathfrak h_{S'}\) is the Lie algebra of
a subgroup of type \textup{(R)} of \(G_{S'}\) are exactly those
that factor through a certain subscheme \(\Sigma\) of finite
presentation over \(S_1\). Replacing \(S_1\) by \(S\), we reduce
to \(S_1=S\), and we may moreover assume that \(S\) is affine;
then there exists an affine Noetherian scheme \(S_0\) such that
\(G\) \textup{(respectively, \(\mathfrak h\))} is obtained by base
change from an adjoint reductive \(S_0\)-group \(G_0\)
\textup{(respectively, from a locally direct-summand submodule
\(\mathfrak h_0\) of
\(\mathfrak g_0=\operatorname{Lie}(G_0/S_0)\))}. It is enough to
show that there exists a subscheme \(\Sigma_0\) of \(S_0\) having
the required properties, for then
\(\Sigma=\Sigma_0\times_{S_0}S\). Replacing \(S\) by \(S_0\), we
may therefore suppose that \(S\) is affine and Noetherian
\textup{(note that then every subscheme of \(S\) is of finite
presentation over \(S\))}. Finally, after replacing \(S\) by a
sufficiently small open subscheme, we may suppose that
\(\mathfrak g\) is free of rank \(n\) and that \(\mathfrak h\) is
a direct summand, free of rank \(r\).

We must first express the condition that \(\mathfrak h_{S'}\) be
a Lie subalgebra of \(\mathfrak g_{S'}\), that is, that the
morphism induced by the Lie bracket
\[
  \varphi\colon
  \mathfrak h\otimes\mathfrak h
  \xrightarrow{[\, ,\,]}
  \mathfrak g
\]
factor through \(\mathfrak h\). If
\((e_1,\ldots,e_n)\) is a basis of \(\mathfrak g\) such that
\((e_1,\ldots,e_r)\) is a basis of \(\mathfrak h\), then
\(\varphi\) is given by sections \(a_k^{ij}\) of
\(\mathcal O_S\), where \(i,j=1,\ldots,r\) and
\(k=1,\ldots,n\). The preceding condition is equivalent to
requiring that \(S'\to S\) factor through the closed subscheme of
\(S\) defined by the equations
\[
  a_k^{ij}=0
  \qquad
  \text{for \(k=r+1,\ldots,n\) and \(i,j=1,\ldots,r\).}
\]
We replace \(S\) by this closed subscheme.

Then, by \SGARefLink{sga3:xxii:recall:5.3.0}{5.3.0},
\[
  N=\underline{\operatorname{Norm}}_G(\mathfrak h)
\]
is a closed subgroup scheme of \(G\), of finite presentation over
\(G\). We must now express \textup{(cf.\ \SGARefLink{sga3:xxii:prop:5.3.1}{5.3.1})} the condition
that \(N_{S'}\) be smooth at every point of the unit section, of
relative dimension \(r=\operatorname{rank}(\mathfrak h)\), and
that the inclusion of \(\mathfrak h_{S'}\) in
\(\operatorname{Lie}(N_{S'}/S')\) be an equality.

Since \(N\) is affine over \(S\)\footnote{Editor's note: The
original then indicated that, denoting by
\(\mathfrak n_{N/G}\) the inverse image of the conormal sheaf
\(\mathcal N_{N/G}\) under the unit section \(S\to N\), the
condition that
\(\mathfrak n_{N_{S'}/G_{S'}}\to
\omega^1_{G_{S'}/S'}\)
be universally injective is equivalent to the condition that
\(S'\to S\) factor through a certain open subscheme of \(S\). We
were unable to justify this point, since formation of
\(\mathcal N_{N/G}\) does not commute with base change, and we
have replaced that argument by the one that follows, suggested by
O.\ Gabber.}
\textup{(being closed in \(G\))}, the unit section
\(\varepsilon\colon S\to N\) is a closed immersion. Thus
\(\varepsilon(S)\) is defined by a quasi-coherent ideal
\(\mathcal J\) of \(\mathcal A(N)\). Note that
\[
  \mathfrak n_{S/N}=\mathcal J/\mathcal J^2
  \simeq \varepsilon^*(\Omega^1_{N/S}),
\]
so its formation commutes with every base change \(S'\to S\).
By the equivalence \textup{\((c')\Longleftrightarrow(a)\)} in
\(\mathrm{EGA}\ \mathrm{IV}_4\), 17.12.1 \textup{(applied to
\(f\colon N\to S\) and \(j=\varepsilon\))}, the morphism
\(N_{S'}\to S'\) is smooth, of relative dimension \(r\), at every
point of \(\varepsilon(S')\) if and only if
\[
  \mathfrak n_{S'/N'}
  =
  \mathfrak n_{S/N}\otimes_{\mathcal O_S}\mathcal O_{S'}
\]
is locally free of rank \(r\), and the morphism
\[
  \varphi_n(S')\colon
  \operatorname{Sym}^n(\mathfrak n_{S'/N'})
  \longrightarrow
  \mathcal J_{S'}^n/\mathcal J_{S'}^{n+1}
\]
is an isomorphism for every \(n\geq1\). Put
\[
  K_n(S')=\ker\varphi_n(S').
\]
By \(\mathrm{TDTE}\) I, Lemma 3.6,
\(\mathfrak n_{S/N}\otimes\mathcal O_{S'}\) is locally free of
rank \(r\) if and only if \(S'\to S\) factors through a certain
subscheme \(Z\) of \(S\). Replacing \(S\) by \(Z\), we may
therefore suppose that
\(\mathfrak n_{S/N}=\mathcal J/\mathcal J^2\) is locally free of
rank \(r\). Then, for every \(S'\to S\),
\[
  (\mathcal J^2/\mathcal J^3)\otimes\mathcal O_{S'}
  =
  \mathcal J_{S'}^2/\mathcal J_{S'}^3,
\]
and consequently
\[
  K_2(S')=K_2(S)\otimes\mathcal O_{S'}.
\]
It follows that \(\varphi_2(S')\) is an isomorphism if and only if
\(S'\to S\) factors through the closed subscheme \(S_2\) of \(S\)
defined by the ideal generated by the image of
\[
  \operatorname{Sym}^2(\mathfrak n_{S/N})^*
  \otimes K_2(S)
  \longrightarrow
  \mathcal O_S.
\]
Over \(S_2\), the module \(\mathcal J^2/\mathcal J^3\) is then
isomorphic to \(\operatorname{Sym}^2(\mathfrak n_{S/N})\), hence
is locally free; the same argument shows that
\(\varphi_3(S')\) is an isomorphism if and only if \(S'\to S\)
factors through a certain closed subscheme \(S_3\) of \(S\), and
so on. We thus obtain that \(N_{S'}\to S'\) is smooth, of relative
dimension \(r\), at every point of \(\varepsilon(S')\) if and only
if \(S'\to S\) factors through the closed subscheme
\[
  Z=\bigcap_n S_n.
\]
But then, for every \(S'\to Z\),
\(\operatorname{Lie}(N_{S'}/S')\) is locally a direct summand of
rank \(r\) of \(\mathfrak g_{S'}\), and hence the inclusion
\[
  \mathfrak h_{S'}
  \subset
  \operatorname{Lie}(N_{S'}/S')
\]
is an equality. We then put \(H=N^0\).

Replacing \(S\) by \(Z\), it remains only to express the condition
that \(H_{s'}\) have the same reductive rank as \(G_{s'}\) at every
point \(s'\in S'\), or, equivalently, that \(H_s\) have the same
reductive rank as \(G_s\) at every point \(s\) in the
set-theoretic image of \(S'\) in \(S\). This condition defines an
open subset of \(S\) \textup{(\SGARefLink{sga3:xix:remark:6-2}{Exposé XIX, 6.2})}.
\hfill\(\square\)

\SGARefTarget{sga3:xxii:rem:after-5.8.1}
\paragraph{Remark.}
In general, the scheme \(\mathcal H\) is not smooth over \(S\). It
is nevertheless smooth if \(6\) is invertible on \(S\), or if
there exists a prime number \(p\) such that \(p\cdot1_S=0\)
\textup{(that is, if \(S\) has characteristic \(p>0\))}.

\SGARefTarget{sga3:xxii:cor:5.8.2}
\paragraph{\SGARefLink{sga3:xxii:cor:5.8.2}{Corollary 5.8.2}.}
\textit{Let \(S\) be a scheme, let \(G\) be a reductive
\(S\)-group, and let \(H\) be a subgroup of type \textup{(R)} of
\(G\). Recall from \SGARefLink{sga3:xxii:cor:5.3.10}{5.3.10} that
\(\underline{\operatorname{Norm}}_G(H)\) is representable by a
closed subgroup scheme of \(G\), smooth over \(S\). Then the
quotient sheaf
\[
  G/\underline{\operatorname{Norm}}_G(H)
\]
is representable by a quasi-projective \(S\)-scheme, smooth and of
finite presentation over \(S\) \textup{(in fact, by an open
subscheme of \(\mathcal H\))}.}

Indeed, consider the morphism
\[
  f\colon G\longrightarrow\mathcal H
\]
defined set-theoretically by
\[
  f(g)=\operatorname{int}(g)H.
\]
By \SGARefLink{sga3:xxii:thm:5.3.9}{5.3.9}, this morphism is smooth and of finite presentation, hence
open. Let \(V=f(G)\), equipped with its open-subscheme structure in
\(\mathcal H\). The morphism \(f\colon G\to V\) is a covering
morphism with kernel \(\underline{\operatorname{Norm}}_G(H)\).
Consequently
\(G/\underline{\operatorname{Norm}}_G(H)\) is representable by
\(V\) \textup{(cf.\ \SGARefLink{sga3:iv:lemma:4:6:5}{Exposé IV, 4.6.5})}.
\hfill\(\square\)

\SGARefTarget{sga3:xxii:cor:5.8.3}
\paragraph{\SGARefLink{sga3:xxii:cor:5.8.3}{Corollary 5.8.3}.}
{\itshape Let \(S\) be a scheme and let \(G\) be a reductive
\(S\)-group. Consider the functors
\(\underline{\operatorname{Tor}}(G)\),
\(\underline{\operatorname{Bor}}(G)\), and
\(\underline{\operatorname{Kil}}(G)\), defined by
\[
\begin{aligned}
  \underline{\operatorname{Tor}}(G)(S')
  &=
  \{\text{maximal tori of \(G_{S'}\)}\},\\
  \underline{\operatorname{Bor}}(G)(S')
  &=
  \{\text{Borel subgroups of \(G_{S'}\)}\},\\
  \underline{\operatorname{Kil}}(G)(S')
  &=
  \{\text{Killing pairs of \(G_{S'}\)
  \textup{ (cf.\ 5.3.13)}}\}.
\end{aligned}
\]
\textup{\SGARefLink{sga3:xxii:cor:5.7.6:item:i}{(i)}}
\(\underline{\operatorname{Tor}}(G)\),
\(\underline{\operatorname{Bor}}(G)\), and
\(\underline{\operatorname{Kil}}(G)\) are representable by smooth
\(S\)-schemes of finite presentation, with geometrically integral
fibres, and are respectively affine, projective, and affine over
\(S\).

\textup{\SGARefLink{sga3:xxii:cor:5.7.6:item:ii}{(ii)}}
The canonical morphism
\[
  \underline{\operatorname{Kil}}(G)
  \longrightarrow
  \underline{\operatorname{Tor}}(G)
\]
\textup{(respectively,
\(\underline{\operatorname{Kil}}(G)\to
\underline{\operatorname{Bor}}(G)\))}
is finite étale and surjective
\textup{(respectively, affine, smooth, and surjective)}.

\SGAInlineRefTarget{sga3:xxii:cor:5.8.3:item:iii}\textup{\SGARefLink{sga3:xxii:prop:5.6.9:item:iii}{(iii)}}
Let \(T\) be a maximal torus of \(G\)
\textup{(respectively, let \(B\) be a Borel subgroup of \(G\), or
let \(B\supset T\) be a Killing pair of \(G\))}. The morphism
\[
\begin{aligned}
  G&\longrightarrow\underline{\operatorname{Tor}}(G),
  &\text{resp.}\quad
  G&\longrightarrow\underline{\operatorname{Bor}}(G),
  &\text{resp.}\quad
  G&\longrightarrow\underline{\operatorname{Kil}}(G),
\end{aligned}
\]
defined by
\[
\begin{aligned}
  g&\longmapsto\operatorname{int}(g)T,
  &\text{respectively}\quad
  g&\longmapsto\operatorname{int}(g)B,
  &\text{respectively}\quad
  g&\longmapsto
  \bigl(\operatorname{int}(g)B,\operatorname{int}(g)T\bigr)
\end{aligned}
\]}

{\itshape induces an isomorphism}
\[
\begin{aligned}
  G/\underline{\operatorname{Norm}}_G(T)
  &\xrightarrow{\;\sim\;}
  \underline{\operatorname{Tor}}(G),
  &\text{resp.}\quad
  G/B
  &\xrightarrow{\;\sim\;}
  \underline{\operatorname{Bor}}(G),
  &\text{resp.}\quad
  G/T
  &\xrightarrow{\;\sim\;}
  \underline{\operatorname{Kil}}(G).
\end{aligned}
\]

One first sees that \textup{\SGARefLink{sga3:xxii:prop:5.6.9:item:iii}{(iii)}} follows from the conjugacy theorem
for maximal tori \textup{(respectively, Borel subgroups, respectively,
Killing pairs)} and from the facts that
\[
  \underline{\operatorname{Norm}}_G(B)=B,
  \qquad
  \underline{\operatorname{Norm}}_G(B)
  \cap
  \underline{\operatorname{Norm}}_G(T)=T,
\]
all of which were established above \textup{(\SGARefLink{sga3:xxii:rem:5.1.2}{5.1.2}, \SGARefLink{sga3:xxii:cor:5.3.12}{5.3.12}, \SGARefLink{sga3:xxii:cor:5.3.14}{5.3.14},
\SGARefLink{sga3:xxii:prop:5.6.1}{5.6.1})}.

It follows first that the canonical morphisms
\[
  \underline{\operatorname{Tor}}(G)\longrightarrow\mathcal H,
  \qquad
  \underline{\operatorname{Bor}}(G)\longrightarrow\mathcal H
\]
are representable, locally for the étale topology, by open immersions
\textup{(\SGARefLink{sga3:xxii:cor:5.8.2}{5.8.2} and \SGARefLink{sga3:xxii:rem:5.1.2}{5.1.2}, respectively \SGARefLink{sga3:xxii:cor:5.5.5}{5.5.5})}. Hence, by descent,
\(\underline{\operatorname{Tor}}(G)\) and
\(\underline{\operatorname{Bor}}(G)\) are representable by open
subschemes of \(\mathcal H\). Likewise,
\(\underline{\operatorname{Kil}}(G)\) is locally, for the étale
topology, representable by a scheme affine over the base
\textup{(\SGARefLink{sga3:ix:proposition:2:3}{Exposé IX, 2.3})}, and hence is representable by an affine
\(S\)-scheme, by descent for affine schemes.

The assertions of \textup{\SGARefLink{sga3:xxii:cor:5.7.6:item:ii}{(ii)}} follow at once from \SGARefLink{sga3:xxii:cor:5.5.5}{5.5.5}\textup{\SGARefLink{sga3:xxii:cor:5.5.5:item:ii}{(ii)}}
and \SGARefLink{sga3:xxii:cor:5.6.13}{5.6.13}. It already follows that
\(\underline{\operatorname{Tor}}(G)\) is affine over \(S\)
\textup{(EGA II, 6.7.1)}. It remains only to prove that
\(\underline{\operatorname{Bor}}(G)\) is projective over \(S\). We
already know that it is quasi-projective, so it remains to prove that
it is proper; but it has connected fibres, and therefore, by EGA
IV\(_3\), 15.7.10, one is reduced to proving this on geometric fibres.
If \(S\) is the spectrum of an algebraically closed field, then
\(\underline{\operatorname{Bor}}(G)=G/B\) by \textup{\SGARefLink{sga3:xxii:prop:5.6.9:item:iii}{(iii)}}, and the
conclusion follows from \emph{Bible}, \S6.4, Theorem 4
\textup{(or [Ch05], \S6.5, Theorem 5)}.
\hfill\(\square\)

\SGARefTarget{sga3:xxii:rem:5.8.4}
\paragraph{\SGARefLink{sga3:xxii:rem:5.8.4}{Remark 5.8.4}.}
Under the hypotheses of \SGARefLink{sga3:xxii:cor:5.8.3}{5.8.3}, let \(Q\) be a central subgroup of
multiplicative type of \(G\). The evident morphisms define
isomorphisms
\[
  \underline{\operatorname{Tor}}(G)
  \simeq
  \underline{\operatorname{Tor}}(G/Q),
  \qquad
  \underline{\operatorname{Bor}}(G)
  \simeq
  \underline{\operatorname{Bor}}(G/Q),
  \qquad
  \underline{\operatorname{Kil}}(G)
  \simeq
  \underline{\operatorname{Kil}}(G/Q).
\]

\SGARefTarget{sga3:xxii:cor:5.8.5}
\paragraph{\SGARefLink{sga3:xxii:cor:5.8.5}{Corollary 5.8.5}.}
{\itshape Let \(S\) be a scheme, \(G\) a reductive \(S\)-group, and
\(P\) a subgroup scheme of \(G\), smooth and of finite presentation
over \(S\). The following conditions are equivalent:

\textup{\SGARefLink{sga3:xxii:cor:5.7.6:item:i}{(i)}} For every \(s\in S\), \(P_{\bar s}\) is a parabolic
subgroup of \(G_{\bar s}\) \textup{(that is, the quotient scheme
\(G_{\bar s}/P_{\bar s}\) is proper over \(\bar s\), or equivalently
\(P_{\bar s}\) contains a Borel subgroup of \(G_{\bar s}\); cf.\
\emph{Bible}, \S6.4, Theorem 4, or [Ch05], \S6.5, Theorem 5)}.

\textup{\SGARefLink{sga3:xxii:cor:5.7.6:item:ii}{(ii)}} The quotient sheaf \(G/P\) is representable by an
\(S\)-scheme smooth and projective over \(S\).

Moreover, under these conditions, \(P\) is closed in \(G\), has
connected fibres, and
\[
  P=\underline{\operatorname{Norm}}_G(P).
\]}

Clearly \textup{\SGARefLink{sga3:xxii:cor:5.7.6:item:ii}{(ii)}} implies \textup{\SGARefLink{sga3:xxii:cor:5.7.6:item:i}{(i)}}. If \textup{\SGARefLink{sga3:xxii:cor:5.7.6:item:i}{(i)}} holds,
then
\[
  P(\bar s)
  =
  \underline{\operatorname{Norm}}_{G(\bar s)}(P_{\bar s})
\]
and \(P_{\bar s}\) is connected \textup{(for the first assertion,
cf.\ \emph{Bible}, \S12.3, \SGARefLink{sga3:xxii:section:4}{Lemma 4};\footnote{Editor's note: We have
supplied the details that follow.} the second assertion}
\textup{follows from this, since
\(P'=P_{\bar s}^{0}\) is a parabolic subgroup of \(G_{\bar s}\)
normalized by \(P_{\bar s}\), whence
\(P'(\bar s)=P(\bar s)\) and therefore \(P'=P_{\bar s}\))}. It
follows that \(P\) is of type \textup{(R)} and that
\[
  P=\underline{\operatorname{Norm}}_G(P),
\]
so that \(P\) is closed in \(G\). By \SGARefLink{sga3:xxii:cor:5.8.2}{5.8.2},
\[
  G/P=G/\underline{\operatorname{Norm}}_G(P)
\]
is representable by a quasi-projective \(S\)-scheme. Its fibres are
connected and proper, and it is therefore projective by the argument
of \SGARefLink{sga3:xxii:cor:5.8.3}{5.8.3}.
\hfill\(\square\)

\SGARefTarget{sga3:xxii:rem:5.8.6}
\paragraph{\SGARefLink{sga3:xxii:rem:5.8.6}{Remark 5.8.6}.}
The statements \SGARefLink{sga3:xxii:thm:5.8.1}{5.8.1}, \SGARefLink{sga3:xxii:cor:5.8.2}{5.8.2}, and \SGARefLink{sga3:xxii:cor:5.8.5}{5.8.5} are valid for an \(S\)-group
of type \textup{(RA)}, or for an \(S\)-group of type \textup{(RR)}
satisfying \SGARefLink{sga3:xxii:prop:5.1.8}{5.1.8}.

\SGARefTarget{sga3:xxii:rem:5.8.7}
\paragraph{\SGARefLink{sga3:xxii:rem:5.8.7}{Remark 5.8.7}.}
Through the inner automorphisms of \(G\), there are canonical actions
\[
\begin{aligned}
  G&\longrightarrow
  \underline{\operatorname{Aut}}_S
  \bigl(\underline{\operatorname{Tor}}(G)\bigr),
  &
  G&\longrightarrow
  \underline{\operatorname{Aut}}_S
  \bigl(\underline{\operatorname{Bor}}(G)\bigr),
  &
  G&\longrightarrow
  \underline{\operatorname{Aut}}_S
  \bigl(\underline{\operatorname{Kil}}(G)\bigr),
\end{aligned}
\]
which, in the situation of \SGARefLink{sga3:xxii:cor:5.8.3}{5.8.3}\textup{\SGARefLink{sga3:xxii:prop:5.6.9:item:iii}{(iii)}}, are identified with
the canonical actions
\[
\begin{aligned}
  G&\longrightarrow
  \underline{\operatorname{Aut}}_S
  \bigl(G/\underline{\operatorname{Norm}}_G(T)\bigr),
  &
  G&\longrightarrow
  \underline{\operatorname{Aut}}_S(G/B),
  &
  G&\longrightarrow
  \underline{\operatorname{Aut}}_S(G/T).
\end{aligned}
\]
In particular, one obtains
\[
\begin{aligned}
  \ker\!\left(
    G\longrightarrow
    \underline{\operatorname{Aut}}_S
    \bigl(\underline{\operatorname{Tor}}(G)\bigr)
  \right)
  &=
  \ker\!\left(
    G\longrightarrow
    \underline{\operatorname{Aut}}_S
    \bigl(\underline{\operatorname{Bor}}(G)\bigr)
  \right)\\
  &=
  \ker\!\left(
    G\longrightarrow
    \underline{\operatorname{Aut}}_S
    \bigl(\underline{\operatorname{Kil}}(G)\bigr)
  \right)\\
  &=
  \underline{\operatorname{Centr}}(G).
\end{aligned}
\]
Indeed, it is clear that \(\underline{\operatorname{Centr}}(G)\)
acts trivially on each of the three schemes. Conversely, the kernel
of
\[
  G\longrightarrow
  \underline{\operatorname{Aut}}_S
  \bigl(\underline{\operatorname{Kil}}(G)\bigr)
\]
is ``the intersection of the maximal tori of \(G\)'' in the sense of
\SGARefLink{sga3:xxii:cor:4.1.7}{4.1.7}, and hence equals
\(\underline{\operatorname{Centr}}(G)\) \textup{(loc.\ cit.)}. For
\(\underline{\operatorname{Bor}}(G)\), one observes that ``the
intersection of the Borel subgroups of \(G\)'' is also ``the
intersection of its maximal tori'' \textup{(see the following
section)}. For \(\underline{\operatorname{Tor}}(G)\), one uses
\SGARefLink{sga3:xii:4.11}{Exposé XII, 4.11}.

\SGARefTarget{sga3:xxii:section:5.9}
\section*{5.9. Properties Peculiar to Borel Subgroups}
\addcontentsline{toc}{section}{\SGARefLink{sga3:xxii:section:5.9}{5.9}. Properties Peculiar to Borel
Subgroups}

Most of these properties will be generalized in Exposé XXVI to
parabolic subgroups.

\SGARefTarget{sga3:xxii:def:5.9.1}
\paragraph{\SGARefLink{sga3:xxii:def:5.9.1}{Definition 5.9.1}.}
Let \(S\) be a scheme, \(G\) a reductive \(S\)-group, and \(B,B'\)
two Borel subgroups of \(G\). We say that \(B\) and \(B'\) are
\emph{in general position} \textup{(or that \(B'\) is in general
position relative to \(B\))} if \(B\cap B'\) is a torus
\textup{(necessarily maximal)} of \(G\).

If \(T\) is a maximal torus of \(G\) contained in \(B\) and \(B'\),
we say that \(B\) and \(B'\) are \emph{opposite} \textup{(relative
to \(T\))} if
\[
  B\cap B'=T.
\]

\SGARefTarget{sga3:xxii:prop:5.9.2}
\paragraph{\SGARefLink{sga3:xxii:prop:5.9.2}{Proposition 5.9.2}.}
{\itshape Let \(S\) be a scheme, \(G\) a reductive \(S\)-group, \(B\)
a Borel subgroup of \(G\), and \(T\) a maximal torus of \(B\). There
exists a unique Borel subgroup \(B'\) of \(G\), opposite to \(B\)
relative to \(T\).

If \((G,T,M,R)\) is a splitting of \(G\) relative to \(T\), and if
\(B=B_{R_+}\) \textup{(\SGARefLink{sga3:xxii:prop:5.5.1}{5.5.1})}, then
\[
  B'=B_{-R_+}.
\]}

By faithfully flat descent, it is enough to prove the proposition in
the split case, with \(B=B_{R_+}\) \textup{(\SGARefLink{sga3:xxii:cor:5.5.5:item:iv}{5.5.5(iv)})}. Then
\(B_{-R_+}\) is indeed opposite to \(B\) \textup{(\SGARefLink{sga3:xxii:prop:4.1.2}{4.1.2})}. Let us
show that it is the only Borel subgroup of \(G\) containing \(T\)
that is opposite to \(B\). If \(B'\) is a Borel subgroup of \(G\)
containing \(T\), then locally on \(S\) it has the form
\(B_{R'_+}\), where \(R'_+\) is another system of positive roots of
\(R\) \textup{(\SGARefLink{sga3:xxii:cor:5.5.5:item:iii}{5.5.5(iii)})}. If \(R'_+\ne -R_+\), there exists
\(\alpha\in R'_+\cap R_+\), and hence
\[
  U_\alpha\subset B_{R_+}\cap B_{R'_+}.
\]
\hfill\(\square\)

\SGARefTarget{sga3:xxii:prop:5.9.3}
\paragraph{\SGARefLink{sga3:xxii:prop:5.9.3}{Proposition 5.9.3}.}
{\itshape Let \(S\) be a scheme, \(G\) a reductive \(S\)-group, and
\(B\) a Borel subgroup of \(G\).

\textup{\SGARefLink{sga3:xxii:cor:5.7.6:item:i}{(i)}} If \(B'\) is a Borel subgroup of \(G\), the following
conditions are equivalent:

\begin{enumerate}
\item[\textup{(a)}] \(B'\) is in general position relative to \(B\)
\textup{(\SGARefLink{sga3:xxii:def:5.9.1}{5.9.1})}.

\item[\textup{\SGAInlineRefTarget{sga3:xxii:prop:5.9.3:item:b}(b)}] \(B'^{u}\cap B^{u}=e\).

\item[\textup{(b\(^{\prime}\))}] \(B'^{u}\cap B=e\).

\item[\textup{\SGAInlineRefTarget{sga3:xxii:prop:5.9.3:item:c}(c)}] Multiplication in \(G\) induces an open immersion
\[
  B'^{u}\times_S B\longrightarrow G.
\]

\item[\textup{(c\(^{\prime}\))}] The canonical morphism
\[
  B'^{u}\longrightarrow G/B
\]
is an open immersion.
\end{enumerate}

\textup{\SGARefLink{sga3:xxii:cor:5.7.6:item:ii}{(ii)}} The functor \(\underline{\operatorname{Opp}}(B)\),
\[
  S'\longmapsto
  \left\{
    \begin{array}{c}
      \text{Borel subgroups of \(G_{S'}\) in general position}\\
      \text{relative to \(B_{S'}\)}
    \end{array}
  \right\},
\]
is representable by an open subscheme of
\(\underline{\operatorname{Bor}}(G)\) \textup{(\SGARefLink{sga3:xxii:cor:5.8.3}{5.8.3})}. The morphism
\[
  \underline{\operatorname{Opp}}(B)
  \longrightarrow
  \underline{\operatorname{Tor}}(B)
\]
defined by \(B'\mapsto B\cap B'\) is an isomorphism. In particular
\textup{(\SGARefLink{sga3:xxii:cor:5.6.13}{5.6.13})}, the inner automorphisms of \(B^u\) equip
\(\underline{\operatorname{Opp}}(B)\) with the structure of a
principal homogeneous bundle under \(B^u\).}

Let us first examine \textup{\SGARefLink{sga3:xxii:cor:5.7.6:item:i}{(i)}}. We have
\(\textup{(a)}\Rightarrow\textup{(c)}\). Indeed,
\textup{\SGARefLink{sga3:xxii:prop:5.9.3:item:c}{(c)}} is local for the étale topology; by \SGARefLink{sga3:xxii:cor:5.5.5:item:iv}{5.5.5(iv)}, one
reduces to the case in which \(G\) is split relative to
\(B\cap B'\) and \(B\) has the form \(B_{R_+}\). By \SGARefLink{sga3:xxii:prop:5.9.2}{5.9.2},
\[
  B'^{u}=U_{-R_+},
\]
and one is reduced to \SGARefLink{sga3:xxii:prop:4.1.2}{4.1.2}.

Trivially,
\[
  \textup{(c\(^{\prime}\))}
  \Longleftrightarrow
  \textup{(c)}
  \Longrightarrow
  \textup{(b\(^{\prime}\))}
  \Longrightarrow
  \textup{(b)}.
\]
It remains to prove
\(\textup{(b)}\Rightarrow\textup{(a)}\). Let us first prove
\textup{\SGARefLink{sga3:xxii:cor:5.7.6:item:ii}{(ii)}}. The second assertion is a formal consequence of
\SGARefLink{sga3:xxii:prop:5.9.2}{5.9.2}, and the third follows at once from it by \SGARefLink{sga3:xxii:cor:5.6.13}{5.6.13}. Let us prove
the first. It is local for the étale topology, and we may therefore
suppose that \(B\) contains a maximal torus \(T\). Let \(B'_0\) be
the Borel subgroup opposite to \(B\) relative to \(T\)
\textup{(\SGARefLink{sga3:xxii:prop:5.9.2}{5.9.2})}.

By what precedes, the morphism
\[
  B^u\longrightarrow\underline{\operatorname{Bor}}(G)
\]
induced by the canonical morphism
\[
  G\longrightarrow G/B'_0
  \longrightarrow\underline{\operatorname{Bor}}(G)
\]
\textup{(\SGARefLink{sga3:xxii:cor:5.8.3}{5.8.3})} induces an isomorphism
\[
  B^u\xrightarrow{\;\sim\;}
  \underline{\operatorname{Opp}}(B).
\]
Thus there is a commutative diagram
\SGARefTarget{sga3:xxii:diag:5.9.3-opp-square}
\[
  \begin{tikzcd}[column sep=5.0em,row sep=3.3em]
  B^u
    \arrow[r]
    \arrow[d,"\sim"]
  &
  G/B'_0
    \arrow[d,"\sim"]
  \\
  \underline{\operatorname{Opp}}(B)
    \arrow[r]
  &
  \underline{\operatorname{Bor}}(G).
\end{tikzcd}

\]
Now the morphism \(B^u\to G/B'_0\) is an open immersion
\textup{(by \textup{\SGARefLink{sga3:xxii:cor:5.7.6:item:i}{(i)}}, \(\textup{(a)}\Rightarrow
\textup{(c\(^{\prime}\))}\))}, which completes the proof of
\textup{\SGARefLink{sga3:xxii:cor:5.7.6:item:ii}{(ii)}}. We record at once the following corollary.

\SGARefTarget{sga3:xxii:cor:5.9.4}
\paragraph{\SGARefLink{sga3:xxii:cor:5.9.4}{Corollary 5.9.4}.}
{\itshape Let \(G\) be a reductive \(S\)-group and let \(B,B'\) be
two Borel subgroups of \(G\). If \(s\in S\) is such that
\(B_{\bar s}\) and \(B'_{\bar s}\) are in general position, there
exists an open subset \(V\) of \(S\), containing \(s\), such that
\(B_V\) and \(B'_V\) are in general position.}

It therefore remains only to prove
\(\textup{(b)}\Rightarrow\textup{(a)}\). By the preceding corollary,
it is enough to do so when \(S\) is the spectrum of an algebraically
closed field \(k\). We may suppose that \(G\) is split relative to a
maximal torus \(T\) of \(B\). Let \(B'_0\) be the Borel subgroup
opposite to \(B\). Since the Borel subgroups of \(G\) are conjugate
under \(G(k)\), there exists \(g\in G(k)\) such that
\(\operatorname{int}(g)B'_0=B'\). By the Bruhat theorem
\textup{(\SGARefLink{sga3:xxii:thm:5.7.4}{5.7.4})}, we may write
\[
  g=bnb',
  \qquad
  b\in B(k),\quad
  b'\in B'_0(k),\quad
  n\in\underline{\operatorname{Norm}}_G(T)(k).
\]
Thus
\[
  B'=\operatorname{int}(b)\operatorname{int}(n)B'_0
\]
and
\[
  B'\cap B
  =
  \operatorname{int}(b)
  \bigl(\operatorname{int}(n)B'_0\cap B\bigr).
\]
If \(n\notin T(k)\), then
\[
  \operatorname{int}(n)(B'_0)^u\cap B^u\ne e
\]
\textup{(cf. the proof of \SGARefLink{sga3:xxii:prop:5.9.2}{5.9.2})}; it follows that
\textup{\SGARefLink{sga3:xxii:prop:5.9.3:item:b}{(b)}} implies
\[
  B'\cap B
  =
  \operatorname{int}(b)(B'_0\cap B)
  =
  \operatorname{int}(b)T.
\]
\hfill\(\square\)

\SGARefTarget{sga3:xxii:prop:5.9.5}
\paragraph{\SGARefLink{sga3:xxii:prop:5.9.5}{Proposition 5.9.5}.}
{\itshape Let \(S\) be a scheme, \(G\) a reductive \(S\)-group,
\(B\) a Borel subgroup of \(G\), and \(B^u\) its unipotent part.
There exists a sequence of subgroups of \(B\),
\[
  U_0=B^u\supset U_1\supset\cdots\supset U_n\supset\cdots,
\]
having the following properties:
\begin{enumerate}
\item[\textup{(i)}] Each \(U_i\) is smooth, has connected fibers,
and is characteristic in \(B\); the inner automorphisms of \(B^u\)
act trivially on the \textup{(sheaf)} quotients \(U_i/U_{i+1}\).

\item[\textup{(ii)}] For each \(i\geq 0\), there exists a locally
free \(\mathcal O_S\)-module \(\mathcal E_i\) and an isomorphism of
sheaves of groups on \(S\),
\[
  U_i/U_{i+1}\xrightarrow{\;\sim\;}W(\mathcal E_i).
\]

\item[\textup{\SGAInlineRefTarget{sga3:xxii:prop:5.9.5:item:iii}(iii)}] For every \(s\in S\),
\[
  (U_n)_s=e
  \qquad\text{for}\qquad
  n\geq\dim(B^u_s).
\]
\end{enumerate}}

Suppose first that there exists a splitting \((G,T,M,R)\) of \(G\)
and a system \(R_+\) of positive roots of \(R\) such that
\(B=B_{R_+}\). Denote by \(\Delta\) the set of simple roots of
\(R_+\). For each \(\alpha\in R_+\), denote by
\(\operatorname{ord}(\alpha)\) the sum of the coefficients of
\(\alpha\) in the basis \(\Delta\) of \(\Gamma_0(R)\); this is the
\emph{order} of \(\alpha\) relative to \(R_+\). We have
\[
  \operatorname{ord}(\alpha)\leq\operatorname{Card}(R_+).
\]
For every \(i>0\), let \(R^{(i)}\) be the set of roots of order
\(>i\). This is a closed set of positive roots, and we may therefore
construct \textup{(\SGARefLink{sga3:xxii:cor:5.6.5}{5.6.5})}
\[
  U_i=U_{R^{(i)}}.
\]
If \(\alpha\in R_+\) and \(\beta\in R^{(i)}\), then
\(\alpha+\beta\in R^{(i+1)}\). It follows from \SGARefLink{sga3:xxii:cor:5.5.2}{5.5.2} that each
\(U_i\) is an invariant subgroup of \(B\), and that the inner
automorphisms of \(B^u\) act trivially on \(U_i/U_{i+1}\). Moreover,
this group identifies with
\[
  \prod_{\operatorname{ord}(\alpha)=i+1}U_\alpha
\]
and is therefore endowed with a vector-group structure.

If \(B\) is of the form \(B_{R'_+}\) for another splitting
\((G,T',M',R')\) of \(G\), let us show that the groups \(U'_i\)
constructed as above using the new splitting coincide with the
\(U_i\), and that the vector-group structures on the successive
quotients likewise coincide. By \SGARefLink{sga3:xxii:cor:5.6.13}{5.6.13}, there exists
\(b\in B^u(S)\) such that \(T'=\operatorname{int}(b)T\); the
assertion to be proved is local on \(S\), and we may therefore
suppose that the isomorphism
\[
  T\xrightarrow{\;\sim\;}T'
\]
induced by \(\operatorname{int}(b)\) arises by duality from an
isomorphism of root data
\[
  h:(M',M'{}^*,R',R'{}^*)
  \xrightarrow{\;\sim\;}
  (M,M^*,R,R^*).
\]
The roots in \(R'_+\) are then the roots
\[
  \alpha\circ\operatorname{int}(b)=h(\alpha),
  \qquad \alpha\in R_+,
\]
and its simple roots are the \(h(\alpha)\), with
\(\alpha\in\Delta\). Consequently
\[
  \operatorname{ord}\bigl(h(\alpha)\bigr)
  =
  \operatorname{ord}(\alpha).
\]
By transport of structure, the vector groups
\(U'_{h(\alpha)}\) are the groups
\(\operatorname{int}(b)U_\alpha\). It follows that
\[
  \operatorname{int}(b)U_i=U'_i.
\]
But \(U_i\) is invariant in \(B\), so that \(U_i=U'_i\).
Likewise, the isomorphism
\[
  \operatorname{int}(b):
  U_i/U_{i+1}\xrightarrow{\;\sim\;}U'_i/U'_{i+1}
\]
is the identity.

We now pass to the general case. Let \(\{S_i\to S\}\) be an étale
covering family such that each
\((G_i,T_i,M_i,R_i)\) is a splitting, and choose a system
\(R_{i+}\) of positive roots such that
\[
  B\times_S S_i=B_{R_{i+}}
\]
\textup{(\SGARefLink{sga3:xxii:cor:5.5.5}{5.5.5}\textup{\SGARefLink{sga3:xxii:cor:5.5.5:item:iii}{(iii)}})}. For every \(i\), the preceding
construction gives a sequence
\[
  B_{S_i}^u=U_{i,0}
  \supset U_{i,1}\supset\cdots
  \supset U_{i,j}\supset\cdots
\]
and vector-group structures on its successive quotients.
By descent, it is enough to verify, for every \(i,i'\) and every
\(j\), that
\[
  U_{i,j}\times_{S_i}S_{ii'}
  =
  U_{i',j}\times_{S_{i'}}S_{ii'},
  \qquad
  S_{ii'}=S_i\times_S S_{i'},
\]
and that the vector-group structures on
\[
  (U_{i,j}/U_{i,j+1})\times_{S_i}S_{ii'}
  \quad\text{and}\quad
  (U_{i',j}/U_{i',j+1})\times_{S_{i'}}S_{ii'}
\]
coincide.\begingroup
\renewcommand{\thefootnote}{*}\footnote{Translator's note: in the second
base change the French source prints \(S_i\); the preceding equality
and \(S_{ii'}=S_i\times_S S_{i'}\) require \(S_{i'}\).}
\addtocounter{footnote}{-1}\endgroup
If \(S_{ii'}=\varnothing\), there is nothing to prove. Otherwise we
are in the situation considered above: \(B\times_S S_{ii'}\) is
defined by the system of positive roots \(R_{i+}\), respectively
\(R_{i'+}\), in the splittings
\[
  (G_{S_{ii'}},
   T_i\times_{S_i}S_{ii'},M_i,R_i),
  \qquad
  (G_{S_{ii'}},
   T_{i'}\times_{S_{i'}}S_{ii'},M_{i'},R_{i'}).
\]
This proves the proposition.
\hfill\(\square\)

\SGARefTarget{sga3:xxii:cor:5.9.6}
\paragraph{\SGARefLink{sga3:xxii:cor:5.9.6}{Corollary 5.9.6}.}
{\itshape If \(S\) is affine, then
\[
  H^1(S,B^u)=e;
\]
that is, every principal homogeneous bundle under \(B^u\) has a
section.}

Indeed, \(S\) is a disjoint union of subschemes on each of which
\(B^u\) has constant relative dimension. In view of
\SGARefLink{sga3:xxii:prop:5.9.5}{5.9.5}\textup{\SGARefLink{sga3:xxii:prop:5.9.5:item:iii}{(iii)}}, we may therefore suppose that \(U_n=e\).
By TDTE I, B.1.1, one has
\[
  H^1(S,U_i/U_{i+1})
  =
  H^1\bigl(S,W(\mathcal E_i)\bigr)
  =
  0.
\]
It follows by induction that \(H^1(S,B^u)=e\).

\SGARefTarget{sga3:xxii:cor:5.9.7}
\paragraph{\SGARefLink{sga3:xxii:cor:5.9.7}{Corollary 5.9.7}.}
{\itshape If \(S\) is affine, then \(B\) has maximal tori. If \(T\)
is a maximal torus of \(B\), then
\[
  H^1(S,T)=H^1(S,B).
\]}

The first assertion follows at once from \SGARefLink{sga3:xxii:cor:5.9.6}{5.9.6} and \SGARefLink{sga3:xxii:cor:5.6.13}{5.6.13}; the
second follows from it in the standard way.

\SGARefTarget{sga3:xxii:cor:5.9.8}
\paragraph{\SGARefLink{sga3:xxii:cor:5.9.8}{Corollary 5.9.8}.}
{\itshape If \(G\) is a reductive \(S\)-group, the canonical
morphism \textup{(cf.\ \SGARefLink{sga3:xxii:cor:5.8.3}{5.8.3})}
\[
  \underline{\operatorname{Kil}}(G)
  \longrightarrow
  \underline{\operatorname{Bor}}(G)
\]
has sections over every affine open subset.}

\SGARefTarget{sga3:xxii:cor:5.9.9}
\paragraph{\SGARefLink{sga3:xxii:cor:5.9.9}{Corollary 5.9.9}.}
{\itshape Under the hypotheses of \SGARefLink{sga3:xxii:prop:5.9.5}{5.9.5}, suppose that \(S\) is
affine. Then there exists a locally free \(\mathcal O_S\)-module
\(\mathcal E\) such that \(B^u\), as a scheme, is \(S\)-isomorphic
to \(W(\mathcal E)\).}

We prove by induction on \(i\) that \(B^u/U_i\) is \(S\)-isomorphic
to
\[
  W(\mathcal E_0\oplus\cdots\oplus\mathcal E_{i-1}).
\]
This is clear for \(i=0\); suppose that \(i\geq1\). Then
\(B^u/U_i\) is a principal homogeneous bundle with base
\[
  X=B^u/U_{i-1}
\]
under the group \((U_{i-1}/U_i)_X\). Since \(B^u/U_{i-1}\) is
affine by the induction hypothesis, and since
\[
  U_{i-1}/U_i=W(\mathcal E_{i-1}),
\]
this bundle is trivial. We therefore have, at least, an isomorphism
of \(S\)-schemes
\[
  B^u/U_i
  \xrightarrow{\;\sim\;}
  (B^u/U_{i-1})\times_S W(\mathcal E_{i-1})
  =
  W(\mathcal E_0\oplus\cdots\oplus\mathcal E_{i-1}).
\]
The conclusion now follows at once from condition
\SGARefLink{sga3:xxii:prop:5.9.5}{5.9.5}\textup{\SGARefLink{sga3:xxii:prop:5.9.5:item:iii}{(iii)}}.

\SGARefTarget{sga3:xxii:cor:5.9.10}
\paragraph{\SGARefLink{sga3:xxii:cor:5.9.10}{Corollary 5.9.10}.}
{\itshape Let \(S\) be a semilocal scheme, let \(\{s_i\}\) be its
closed points, and let \(B\) be a Borel subgroup of the reductive
\(S\)-group \(G\). The canonical map
\[
  B^u(S)
  \longrightarrow
  \prod_i B^u\bigl(\operatorname{Spec}\kappa(s_i)\bigr)
\]
is surjective.}

Indeed, if \(S=\operatorname{Spec}(A)\),
\(\kappa(s_i)=A/\mathfrak p_i\), and \(\mathcal E\) is defined by
the \(A\)-module \(E\), then
\[
  B^u(S)=E\otimes A,
  \qquad
  B^u\bigl(\operatorname{Spec}\kappa(s_i)\bigr)
  =
  E\otimes_A(A/\mathfrak p_i).
\]
The assertion follows at once from the surjectivity of
\[
  A\longrightarrow\prod_i A/\mathfrak p_i.
\]

\SGARefTarget{sga3:xxii:section:5.10}
\section*{5.10. Subgroups of Type \textup{(R)} with Reductive
Fibers}
\addcontentsline{toc}{section}{\SGARefLink{sga3:xxii:section:5.10}{5.10}. Subgroups of Type (R) with
Reductive Fibers}

\SGARefTarget{sga3:xxii:prop:5.10.1}
\paragraph{\SGARefLink{sga3:xxii:prop:5.10.1}{Proposition 5.10.1}.}
{\itshape Let \((G,T,M,R)\) be a split \(S\)-group, let \(R'\) be a
subset of \(R\) of type \textup{(R)} \textup{(\SGARefLink{sga3:xxii:def:5.4.2}{5.4.2})}, and let
\(H_{R'}\) be the corresponding subgroup of \(G\). The following
conditions are equivalent:
\begin{enumerate}
\item[\textup{\SGAInlineRefTarget{sga3:xxii:prop:5.10.1:item:i}(i)}] \(H_{R'}\) is reductive \textup{(that is, its
geometric fibers are reductive)}.

\item[\textup{\SGAInlineRefTarget{sga3:xxii:prop:5.10.1:item:ii}(ii)}] \(R'=-R'\), that is, \(R'\) is symmetric.
\end{enumerate}
Moreover, under these conditions, \((H_{R'},T,M,R')\) is a
splitting of \(H_{R'}\); for every system \(R_+\) of positive roots
of \(R\),
\[
  R'_+=R'\cap R_+
\]
is a system of positive roots of \(R'\), and
\[
  B_{R_+}\cap H_{R'}=H_{R'_+}.
\]
This is a Borel subgroup of \(H_{R'}\), whose unipotent part is
\[
  U_{R_+}\cap H_{R'}=U_{R'_+}.
\]}

Clearly \textup{\SGARefLink{sga3:xxii:prop:5.10.1:item:i}{(i)}} implies \textup{\SGARefLink{sga3:xxii:prop:5.10.1:item:ii}{(ii)}}: it is enough to check
this fiber by fiber, and \(R'\) is the root system of \(H_{R'}\)
relative to \(T\). To prove that \textup{\SGARefLink{sga3:xxii:prop:5.10.1:item:ii}{(ii)}} implies
\textup{\SGARefLink{sga3:xxii:prop:5.10.1:item:i}{(i)}}, observe, by \SGARefLink{sga3:xxii:lem:5.4.3}{5.4.3}, that
\[
  H_{R'}\cap Z_\alpha
  =
  \underline{\operatorname{Centr}}_{H_{R'}}(T_\alpha)
  =
  Z_\alpha
\]
for every \(\alpha\in R'\), and apply the criterion of \SGARefLink{sga3:xix:proposition:1-12}{Exposé XIX,
1.12}.

If \(R_+\) is a system of positive roots of \(R\), then
\(R'_+=R_+\cap R'\) is plainly a closed subset of \(R'\) such that
\[
  R'_+\cup(-R'_+)=R',
  \qquad
  R'_+\cap(-R'_+)=\varnothing;
\]
it is therefore a system of positive roots of \(R'\). The other two
assertions follow respectively from \SGARefLink{sga3:xxii:prop:5.6.1}{5.6.1}\textup{\SGARefLink{sga3:xxii:prop:5.6.1:item:vi}{(vi)}} and
\SGARefLink{sga3:xxii:cor:5.6.7}{5.6.7}\textup{\SGARefLink{sga3:xxii:cor:5.6.7:item:i}{(i)}}.
\hfill\(\square\)

\SGARefTarget{sga3:xxii:cor:5.10.2}
\paragraph{\SGARefLink{sga3:xxii:cor:5.10.2}{Corollary 5.10.2}.}
{\itshape Let \(S\) be a scheme, \(G\) a reductive \(S\)-group,
and \(H\) a reductive \(S\)-subgroup such that, for every
\(s\in S\), \(G_{\bar s}\) and \(H_{\bar s}\) have the same
reductive rank. Then \(H\) is closed in \(G\),
\(\underline{\operatorname{Norm}}_G(H)\) is smooth over \(S\), and
\[
  \underline{\operatorname{Norm}}_G(H)/H
\]
is representable by a finite étale \(S\)-scheme.

If \(T\) is a maximal torus of \(H\), and \(B\) a Borel subgroup of
\(G\) containing \(T\), then \(B\cap H\) is a Borel subgroup of
\(H\), whose unipotent part is
\[
  (B\cap H)^u=B^u\cap H.
\]}

The first assertions follow at once from \SGARefLink{sga3:xxii:cor:5.3.10}{5.3.10} and \SGARefLink{sga3:xxii:cor:5.3.18}{5.3.18}, using
the fact that the Weyl groups of \(G\) and \(H\) are finite
\textup{(\SGARefLink{sga3:xix:theorem:2:5}{Exposé XIX, 2.5})}. The other assertions are local for the
étale topology and reduce to the case studied in \SGARefLink{sga3:xxii:prop:5.10.1}{5.10.1}.
\hfill\(\square\)

\SGARefTarget{sga3:xxii:prop:5.10.3}
\paragraph{\SGARefLink{sga3:xxii:prop:5.10.3}{Proposition 5.10.3}.}
{\itshape Let \(S\) be a scheme and \(G\) a reductive \(S\)-group.
\begin{enumerate}
\item[\textup{\SGAInlineRefTarget{sga3:xxii:prop:5.10.3:item:a}(a)}] If \(Q\) is a torus of \(G\), then
\(\underline{\operatorname{Centr}}_G(Q)\) is a subgroup of \(G\) of
type \textup{(R)} with reductive fibers. If \(Q\subset Q'\) are two
tori of \(G\), then
\[
  \underline{\operatorname{Centr}}_G(Q)
  \supset
  \underline{\operatorname{Centr}}_G(Q').
\]

\item[\textup{\SGAInlineRefTarget{sga3:xxii:prop:5.10.3:item:b}(b)}] If \(H\) is a subgroup of \(G\) of type
\textup{(R)} with reductive fibers, then
\(\operatorname{rad}(H)\) \textup{(\SGARefLink{sga3:xxii:def:4.3.6}{4.3.6})} is a torus of \(G\). If
\(H\subset H'\) are two subgroups of \(G\) of type \textup{(R)}
with reductive fibers, then
\[
  \operatorname{rad}(H)\supset\operatorname{rad}(H').
\]

\item[\textup{\SGAInlineRefTarget{sga3:xxii:prop:5.10.3:item:c}(c)}] If \(Q\) is a torus of \(G\), then
\[
  \operatorname{rad}\bigl(
    \underline{\operatorname{Centr}}_G(Q)\bigr)
  \supset Q
  \quad\text{and}\quad
  \underline{\operatorname{Centr}}_G\!\left(
    \operatorname{rad}\bigl(
      \underline{\operatorname{Centr}}_G(Q)\bigr)\right)
  =
  \underline{\operatorname{Centr}}_G(Q).
\]

\item[\textup{\SGAInlineRefTarget{sga3:xxii:prop:5.10.3:item:d}(d)}] If \(H\) is a subgroup of \(G\) of type
\textup{(R)} with reductive fibers, then
\[
  \underline{\operatorname{Centr}}_G\bigl(\operatorname{rad}(H)\bigr)
  \supset H
  \quad\text{and}\quad
  \operatorname{rad}\!\left(
    \underline{\operatorname{Centr}}_G\bigl(
      \operatorname{rad}(H)\bigr)\right)
  =
  \operatorname{rad}(H).
\]
\end{enumerate}}

Indeed, \textup{\SGARefLink{sga3:xxii:prop:5.10.3:item:a}{(a)}} follows at once from \SGARefLink{sga3:xix:paragraph:2:8}{Exposé XIX, 2.8}. To prove
\textup{\SGARefLink{sga3:xxii:prop:5.10.3:item:b}{(b)}}, it is enough to observe that
\(\operatorname{rad}(H')\subset H\), since \(H\) contains,
fpqc-locally, a maximal torus of \(G\), hence of \(H'\). The first
assertion of \textup{\SGARefLink{sga3:xxii:prop:5.10.3:item:c}{(c)}} \textup{(respectively, of \textup{\SGARefLink{sga3:xxii:prop:5.10.3:item:d}{(d)}})}
is trivial, and the second follows by the usual argument.
\hfill\(\square\)

This proposition leads to the following definition:

\SGARefTarget{sga3:xxii:def:5.10.4}
\paragraph{\SGARefLink{sga3:xxii:def:5.10.4}{Definition 5.10.4}.}
Let \(S\) be a scheme, \(G\) a reductive \(S\)-group, \(H\) a
reductive subgroup of \(G\) of type \textup{(R)}, and \(Q\) a
subtorus of \(G\).\footnote{Editor's note: We have modified the
original by introducing the term ``\(C\)-critical torus'' in place
of ``critical torus'', in order to avoid confusion in later
references \textup{(cf.\ \SGARefLink{sga3:xxvi:item:3:9}{Exposé XXVI, 3.9})}. We have also expanded
the statement of \SGARefLink{sga3:xxii:cor:5.10.5}{5.10.5} and added \SGARefLink{sga3:xxii:rem:5.10.5.1}{Remark 5.10.5.1}.}
\begin{enumerate}
\item[\textup{1)}] \(H\) is called a \emph{critical subgroup} if it
is the centralizer of its radical.
\item[\textup{2)}] \(Q\) is called a \emph{\(C\)-critical torus} if
it is the radical of its centralizer.
\end{enumerate}

\SGARefLink{sga3:xxii:prop:5.10.3}{Proposition 5.10.3} then gives:

\SGARefTarget{sga3:xxii:cor:5.10.5}
\paragraph{\SGARefLink{sga3:xxii:cor:5.10.5}{Corollary 5.10.5}.}
\begin{enumerate}
\item[\textup{(i)}] For every subtorus \(Q\) of \(G\),
\(\underline{\operatorname{Centr}}_G(Q)\) is critical.

\item[\textup{(ii)}] For every subgroup \(H\) of \(G\) of type
\textup{(R)} with reductive fibers, \(\operatorname{rad}(H)\) is a
\(C\)-critical torus of \(G\).

\item[\textup{(iii)}] The maps
\[
  Q\longmapsto\underline{\operatorname{Centr}}_G(Q),
  \qquad
  H\longmapsto\operatorname{rad}(H)
\]
are mutually inverse bijections between the set of \(C\)-critical
tori of \(G\) and the set of its critical reductive subgroups of
type \textup{(R)}.

\item[\textup{(iv)}] If \(Q\) is a torus of \(G\), then
\[
  \operatorname{rad}\bigl(
    \underline{\operatorname{Centr}}_G(Q)\bigr)
\]
is the smallest \(C\)-critical torus of \(G\) containing \(Q\).

\item[\textup{(v)}] If \(H\) is a reductive subgroup of \(G\) of
type \textup{(R)}, then
\[
  \underline{\operatorname{Centr}}_G\bigl(\operatorname{rad}(H)\bigr)
\]
is the smallest critical reductive subgroup of \(G\) of type
\textup{(R)} containing \(H\).
\end{enumerate}

\SGARefTarget{sga3:xxii:rem:5.10.5.1}
\paragraph{\SGARefLink{sga3:xxii:rem:5.10.5.1}{Remark 5.10.5.1}.}
\textup{(58)}
\begin{enumerate}
\item[\textup{1)}] A torus \(T\) of \(G\) is a critical subgroup of
\(G\) if and only if it is a maximal torus.
\item[\textup{2)}] Henceforth, ``critical torus'' means
``\(C\)-critical torus''.
\end{enumerate}

\SGARefTarget{sga3:xxii:prop:5.10.6}
\paragraph{\SGARefLink{sga3:xxii:prop:5.10.6}{Proposition 5.10.6}.}
{\itshape Let \((G,T,M,R)\) be a split \(S\)-group, and let
\(R'\) be a subset of \(R\). The following conditions are
equivalent:
\begin{enumerate}
\item[\textup{\SGAInlineRefTarget{sga3:xxii:prop:5.10.6:item:i}(i)}] \(R'\) is of type \textup{(R)}, and \(H_{R'}\)
is reductive and critical.

\item[\textup{(ii)}] There exist a system \(\Delta\) of simple
roots of \(R\) and a subset \(\Delta'\) of \(\Delta\) such that
\(R'\) is the set of those elements of \(R\) that are linear
combinations of elements of \(\Delta'\).

\item[\textup{\SGAInlineRefTarget{sga3:xxii:prop:5.10.6:item:iii}(iii)}] \(R'\) is closed and symmetric, and every
system of simple roots of \(R'\) is the intersection with \(R'\) of
a system of simple roots of \(R\).
\end{enumerate}}

Indeed, by \SGARefLink{sga3:XXI:3.4.8}{Exposé XXI, 3.4.8}, conditions \textup{\SGARefLink{sga3:XXI:3.4.8:item:ii}{(ii)}} and
\textup{\SGARefLink{sga3:xxii:prop:5.10.6:item:iii}{(iii)}} are equivalent; they are also equivalent to the
condition that \(R'\) be the intersection of \(R\) with a
\(\mathbf Q\)-vector subspace of \(M\otimes\mathbf Q\). This last
condition follows from \textup{\SGARefLink{sga3:xxii:prop:5.10.6:item:i}{(i)}}: if
\[
  H_{R'}=\underline{\operatorname{Centr}}_G(Q),
\]
then \(R'\) is the set of elements of \(R\) that vanish on \(Q\)
\textup{(\SGARefLink{sga3:ii:theorem:5-2-3:item:ii}{Exposé II, 5.2.3(ii)})}. Conversely, this condition implies
\textup{\SGARefLink{sga3:xxii:prop:5.10.6:item:i}{(i)}}, because \(\operatorname{rad}(H_{R'})\) is the maximal
torus of
\[
  \bigcap_{\alpha\in R'}\operatorname{Ker}(\alpha).
\]
Consequently,
\[
  \underline{\operatorname{Centr}}_G\bigl(
    \operatorname{rad}(H_{R'})\bigr)
  =
  H_{R''},
\]
where \(R''\) is the intersection of \(R\) with the vector subspace
spanned by \(R'\); hence \(R''=R'\).
\hfill\(\square\)

\SGARefTarget{sga3:xxii:para:5.10.7}
\paragraph{\SGARefLink{sga3:xxii:para:5.10.7}{5.10.7}.}
Let us summarize some of the preceding results. Let
\((G,T,M,R)\) be a split \(S\)-group, let \(\Delta\) be a system
of simple roots of \(R\), and let \(R_+\) be the corresponding
system of positive roots. Choose a subset \(\Delta'\) of
\(\Delta\), denote by \(R'\) the set of those elements of \(R\)
which are linear combinations of elements of \(\Delta'\), and put
\[
  R'_+=R'\cap R_+.
\]
Let \(T_{\Delta'}\) be the maximal torus of
\[
  \bigcap_{\alpha\in\Delta'}\operatorname{Ker}(\alpha)
\]
and put
\[
  Z_{\Delta'}
  =
  \underline{\operatorname{Centr}}_G(T_{\Delta'}).
\]
Then \(Z_{\Delta'}\) is a reductive subgroup of \(G\), with radical
\(T_{\Delta'}\); \((Z_{\Delta'},T,M,R')\) is a split
\(S\)-group; \(B_{R_+}\cap Z_{\Delta'}\) is the Borel subgroup of
\(Z_{\Delta'}\) defined by the system \(R'_+\) of positive roots
\textup{(or, equivalently, by the system \(\Delta'\) of simple
roots)}, and its unipotent part is
\[
  U_{R_+}\cap Z_{\Delta'}=U_{R'_+}.
\]

\SGARefTarget{sga3:xxii:rem:5.10.8}
\paragraph{\SGARefLink{sga3:xxii:rem:5.10.8}{Remark 5.10.8}.}
Under the conditions of \SGARefLink{sga3:xxii:def:5.10.4}{5.10.4}, let \(Q\) be a critical torus of
\(G\), and let
\[
  L=\underline{\operatorname{Centr}}_G(Q)
\]
be its centralizer. Since \(Q=\operatorname{rad}(L)\), \(Q\) is a
characteristic subgroup of \(L\); it follows at once that
\[
  \underline{\operatorname{Norm}}_G(L)
  =
  \underline{\operatorname{Norm}}_G(Q),
\]
and hence also that
\[
  \underline{\operatorname{Norm}}_G(L)/L
  =
  \underline{\operatorname{Norm}}_G(Q)/
  \underline{\operatorname{Centr}}_G(Q)
  =
  W_G(Q).
\]
By \SGARefLink{sga3:xxii:cor:5.10.2}{5.10.2}, we deduce:

\SGARefTarget{sga3:xxii:prop:5.10.9}
\paragraph{\SGARefLink{sga3:xxii:prop:5.10.9}{Proposition 5.10.9}.}
{\itshape Let \(S\) be a scheme, let \(G\) be a reductive
\(S\)-group, and let \(Q\) be a critical torus of \(G\). The Weyl
group \(W_G(Q)\) is finite \textup{(étale)} over \(S\).}

\SGARefTarget{sga3:xxii:rem:5.10.10}
\paragraph{\SGARefLink{sga3:xxii:rem:5.10.10}{Remark 5.10.10}.}
Under the conditions of \SGARefLink{sga3:xxii:para:5.10.7}{5.10.7}, one may describe explicitly
\[
  W_G(T_{\Delta'})
  =
  \underline{\operatorname{Norm}}_G(Z_{\Delta'})/Z_{\Delta'}.
\]
This is the constant group associated with the quotient \(W_1/W_2\),
where \(W_1\) is the subgroup of \(W\) consisting of those elements
which normalize the subgroup of \(M\) generated by \(\Delta'\), and
\(W_2\) is the subgroup of \(W\) generated by the \(s_\alpha\), for
\(\alpha\in\Delta'\).

\SGARefTarget{sga3:xxii:section:5.11}
\section*{5.11. Subgroups of Type \textup{(RC)}}
\addcontentsline{toc}{section}{\SGARefLink{sga3:xxii:section:5.11}{5.11}. Subgroups of Type (RC)}

\SGARefTarget{sga3:xxii:def:5.11.1}
\paragraph{\SGARefLink{sga3:xxii:def:5.11.1}{Definition 5.11.1}.}
Let \(S\) be a scheme and let \(G\) be a reductive \(S\)-group. A
group subscheme \(H\) of \(G\) is said to be \emph{of type
\textup{(RC)}} if it is of type \textup{(R)}, that is, if it
satisfies the following two conditions from \SGARefLink{sga3:xxii:def:5.2.1}{5.2.1}:
\begin{enumerate}
\item[\textup{(i)}] \(H\) is smooth over \(S\), with connected
fibers;

\item[\textup{(ii)}] for every \(s\in S\), \(H_{\bar s}\) contains
a maximal torus of \(G_{\bar s}\);
\end{enumerate}
and if it also satisfies the following condition:
\begin{enumerate}
\item[\textup{\SGAInlineRefTarget{sga3:xxii:def:5.11.1:item:iii}(iii)}] for every \(s\in S\) and every maximal torus
\(T\) of \(H_{\bar s}\), the set of roots of \(H_{\bar s}\)
relative to \(T\) is a closed subset of the set of all roots of
\(G_{\bar s}\) relative to \(T\).
\end{enumerate}

\SGARefTarget{sga3:xxii:rem:5.11.2}
\paragraph{\SGARefLink{sga3:xxii:rem:5.11.2}{Remark 5.11.2}.}
As we already observed in \SGARefLink{sga3:xxii:rem:5.4.8}{5.4.8}, condition \textup{\SGARefLink{sga3:xxii:def:5.11.1:item:iii}{(iii)}} follows
from the other two when \(6\) is invertible on
\(S\).

\SGARefTarget{sga3:xxii:lem:5.11.3}
\paragraph{\SGARefLink{sga3:xxii:lem:5.11.3}{Lemma 5.11.3}.}
{\itshape Let \((G,T,M,R)\) be a split \(S\)-group, and let \(R'\)
be a closed subset of \(R\). Put
\[
  R_1=\{\alpha\in R'\mid-\alpha\in R'\},
  \qquad
  R_2=\{\alpha\in R'\mid-\alpha\notin R'\}.
\]
Then \(R_1\) and \(R_2\) are closed. Consider the groups
\(H_{R'}\), \(H_{R_1}\), and \(U_{R_2}\) \textup{(\SGARefLink{sga3:xxii:thm:5.4.7}{5.4.7} and
\SGARefLink{sga3:xxii:cor:5.6.5}{5.6.5})}, which are smooth and have connected fibers.
\begin{enumerate}
\item[\textup{\SGAInlineRefTarget{sga3:xxii:lem:5.11.3:item:i}(i)}] \(U_{R_2}\) is normal in \(H_{R'}\), and
\(H_{R'}\) is the semidirect product of \(U_{R_2}\) by
\(H_{R_1}\).

\item[\textup{\SGAInlineRefTarget{sga3:xxii:lem:5.11.3:item:ii}(ii)}] The group \(H_{R_1}\) is reductive, and
\(U_{R_2}\) has connected unipotent geometric fibers. Every normal
subgroup scheme of \(H_{R'}\) that is smooth over \(S\) and has
connected unipotent geometric fibers is contained in \(U_{R_2}\).
Every reductive subgroup of \(H_{R'}\) that contains \(T\) is
contained in \(H_{R_1}\).

\item[\textup{\SGAInlineRefTarget{sga3:xxii:lem:5.11.3:item:iii}(iii)}]
\[
  U_{R_2}\cap\underline{\operatorname{Norm}}_G(H_{R_1})=e.
\]
\end{enumerate}}

Let us first prove \textup{\SGARefLink{sga3:xxii:lem:5.11.3:item:iii}{(iii)}}, which follows from \SGARefLink{sga3:xxii:cor:5.6.7:item:i}{5.6.7(i)}. The
first assertion of \textup{\SGARefLink{sga3:xxii:lem:5.11.3:item:i}{(i)}} follows from \SGARefLink{sga3:xxii:cor:5.6.7:item:ii}{5.6.7(ii)}. Since
\[
  U_{R_2}\cap H_{R_1}=e
\]
by \textup{\SGARefLink{sga3:xxii:lem:5.11.3:item:iii}{(iii)}}, the semidirect product \(H_{R_1}\cdot U_{R_2}\)
is a subgroup of \(H_{R'}\). But these are two subgroups of \(G\)
of type \textup{(R)}, containing \(T\), and having the same Lie
algebra \(\mathfrak g_{R'}\); they coincide by \SGARefLink{sga3:xxii:cor:5.3.5}{5.3.5}, which proves
\textup{\SGARefLink{sga3:xxii:lem:5.11.3:item:i}{(i)}}.

For \textup{\SGARefLink{sga3:xxii:lem:5.11.3:item:ii}{(ii)}}, the first two assertions are merely \SGARefLink{sga3:xxii:prop:5.10.1}{5.10.1} and
\SGARefLink{sga3:xxii:cor:5.6.5}{5.6.5}. Let \(U\) be a subgroup scheme of \(H_{R'}\), smooth and of
finite presentation over \(S\), normal in \(H_{R'}\) \textup{(and
hence normalized by \(T\))}, and with connected unipotent geometric
fibers. By \SGARefLink{sga3:xxii:cor:5.6.12}{5.6.12}, locally one has \(U=U_{R''}\), where
\(R''\subset R'\) and
\[
  R''\cap(-R'')=\varnothing.
\]
If \(U\) were not contained in \(U_{R_2}\), then \(R''\) would not
be contained in \(R_2\); there would therefore exist
\(\alpha\in R''\) such that \(-\alpha\in R'\). Then
\(Z_\alpha\subset H_{R'}\) by \SGARefLink{sga3:xxii:lem:5.4.3}{5.4.3}, so \(Z_\alpha\) normalizes
\(U\). But \(U\) contains \(U_\alpha\), and \(Z_\alpha\) has a
section \(w\) such that
\[
  \operatorname{int}(w)U_\alpha=U_{-\alpha};
\]
thus \(-\alpha\in R''\), a contradiction.

Finally, if \(L\) is a reductive subgroup of \(H_{R'}\) containing
\(T\), locally one has \(L=H_{R'''}\), where \(R'''\) is symmetric
and contained in \(R'\); hence \(R'''\subset R_1\).
\hfill\(\square\)

\SGARefTarget{sga3:xxii:prop:5.11.4}
\paragraph{\SGARefLink{sga3:xxii:prop:5.11.4}{Proposition 5.11.4}.}
{\itshape Let \(S\) be a scheme, let \(G\) be a reductive
\(S\)-group, and let \(H\) be a subgroup scheme of \(G\) of type
\textup{(RC)}.
\begin{enumerate}
\item[\textup{\SGAInlineRefTarget{sga3:xxii:prop:5.11.4:item:i}(i)}] \(H\) is closed in \(G\), and
\(\underline{\operatorname{Norm}}_G(H)/H\) is representable by a
finite étale \(S\)-group scheme.

\item[\textup{(ii)}] \(H\) has a largest normal subgroup scheme
that is smooth and of finite presentation over \(S\), with
connected unipotent geometric fibers. It is called the
\emph{unipotent radical} of \(H\) and is denoted
\(\operatorname{rad}^{u}(H)\). The quotient sheaf
\[
  H/\operatorname{rad}^{u}(H)
\]
is representable by a reductive \(S\)-group.

\item[\textup{(iii)}] If \(T\) is a maximal torus of \(H\), then
\(H\) has a reductive subgroup \(L\), containing \(T\), of type
\textup{(RC)}, with the following properties:
\begin{enumerate}
\item[\textup{\SGAInlineRefTarget{sga3:xxii:prop:5.11.4:item:a}(a)}] every reductive subgroup of \(H\) containing
\(T\) is contained in \(L\);

\item[\textup{\SGAInlineRefTarget{sga3:xxii:prop:5.11.4:item:b}(b)}] \(H\) is the semidirect product
\[
  H=L\cdot\operatorname{rad}^{u}(H);
\]
in other words, the canonical morphism
\[
  L\longrightarrow H/\operatorname{rad}^{u}(H)
\]
is an isomorphism.
\end{enumerate}
\end{enumerate}

Moreover, \(L\) is the unique reductive subgroup of \(H\) that
contains \(T\) and satisfies either one of the two preceding
conditions. Finally,
\[
  \underline{\operatorname{Norm}}_H(L)=L,
  \qquad
  \underline{\operatorname{Norm}}_H(T)
    =\underline{\operatorname{Norm}}_L(T),
  \qquad
  W_H(T)=W_L(T);
\]
in particular, \(W_H(T)\) is finite over \(S\).}

\textit{Proof.}
First observe that \textup{\SGARefLink{sga3:xxii:prop:5.11.4:item:i}{(i)}} is local for the étale topology.
Thus, by \SGARefLink{sga3:xxii:cor:5.3.18}{Corollary 5.3.18}, \textup{\SGARefLink{sga3:xxii:cor:5.3.18:item:i}{(i)}} follows from the last
assertion of \textup{\SGARefLink{sga3:xxii:cor:5.3.18:item:iii}{(iii)}}.

The assertions of \textup{\SGARefLink{sga3:xxii:cor:5.3.18:item:ii}{(ii)}} are local for the étale topology.
We may therefore suppose ourselves in the situation of \SGARefLink{sga3:xxii:lem:5.11.3}{5.11.3},
where they follow immediately from \textup{\SGARefLink{sga3:xxii:lem:5.11.3:item:i}{(i)}} and \textup{\SGARefLink{sga3:xxii:lem:5.11.3:item:ii}{(ii)}}.

By virtue of the uniqueness assertions contained there,
\textup{\SGARefLink{sga3:xxii:lem:5.11.3:item:iii}{(iii)}} is likewise local for the étale topology, and we may
again reduce to the situation of \SGARefLink{sga3:xxii:lem:5.11.3}{5.11.3}, where properties
\textup{\SGARefLink{sga3:xxii:prop:5.11.4:item:a}{(a)}} and \textup{\SGARefLink{sga3:xxii:prop:5.11.4:item:b}{(b)}} have been verified. The uniqueness of
an \(L\) satisfying \textup{\SGARefLink{sga3:xxii:prop:5.11.4:item:a}{(a)}} is immediate; in view of
\textup{\SGARefLink{sga3:xxii:prop:5.11.4:item:a}{(a)}}, the uniqueness of an \(L\) satisfying \textup{\SGARefLink{sga3:xxii:prop:5.11.4:item:b}{(b)}}
is clear. The equality
\[
  \underline{\operatorname{Norm}}_H(L)=L
\]
is precisely \SGARefLink{sga3:xxii:lem:5.11.3:item:iii}{5.11.3(iii)}. If a section of \(H\) normalizes \(T\),
then it normalizes \(L\), by the uniqueness of \(L\), and hence is
a section of \(L\) by what we have just proved. This proves the
second equality; the third is then immediate.

\SGARefTarget{sga3:xxii:prop:5.11.5}
\paragraph{\SGARefLink{sga3:xxii:prop:5.11.5}{Proposition 5.11.5}.}
{\itshape Let \(S\) be a scheme, let \(G\) be a reductive
\(S\)-group, and let \(\mathcal H_c\) be the functor of subgroups
of \(G\) of type \textup{(RC)}, which is a subfunctor of the
functor \(\mathcal H\) of \SGARefLink{sga3:xxii:thm:5.8.1}{5.8.1}.
\begin{enumerate}
\item[\textup{\SGAInlineRefTarget{sga3:xxii:prop:5.11.5:item:i}(i)}] \(\mathcal H_c\) is representable by an open
subscheme of \(\mathcal H\), smooth, quasi-projective, and of finite
presentation over \(S\).

\item[\textup{\SGAInlineRefTarget{sga3:xxii:prop:5.11.5:item:ii}(ii)}] There exist a finite étale \(S\)-scheme
\(\mathcal C\ell_c\) and a morphism
\[
  \operatorname{cl}\colon\mathcal H_c\longrightarrow\mathcal C\ell_c
\]
that is smooth, quasi-projective, of finite presentation,
surjective, and has connected geometric fibers, with the following
property:

For every \(S'\to S\) and all
\(H,H'\in\mathcal H_c(S')\), one has
\(\operatorname{cl}(H)=\operatorname{cl}(H')\) if and only if
\(H\) and \(H'\) are conjugate in \(G\) locally for the étale
topology \textup{(or, equivalently by \SGARefLink{sga3:xxii:cor:5.3.11}{5.3.11}, if for every
\(s\in S\), \(H_{\bar s}\) and \(H'_{\bar s}\) are conjugate by an
element of \(G(\bar s)\))}.

\item[\textup{\SGAInlineRefTarget{sga3:xxii:prop:5.11.5:item:iii}(iii)}] \(\mathcal C\ell_c\) and
\(\operatorname{cl}\) are determined \textup{(up to a unique
isomorphism)} by the preceding conditions.

\item[\textup{\SGAInlineRefTarget{sga3:xxii:prop:5.11.5:item:iv}(iv)}] If \((G,T,M,R)\) is a splitting of \(G\), let
\(E\) be the set of conjugacy classes modulo \(W\) of closed
subsets of \(R\). Then there is an isomorphism
\[
  \mathcal C\ell_c\xrightarrow{\;\sim\;}E_S
\]
such that, for every closed subset \(R'\) of \(R\),
\(\operatorname{cl}(H_{R'})\) corresponds to the canonical image
of \(R'\) in
\[
  E_S(S)=\operatorname{Hom}_{\mathrm{loc.const.}}(S,E).
\]
\end{enumerate}}

It is first of all clear that \(\mathcal H_c\) is a sheaf for the
étale topology, and that \textup{\SGARefLink{sga3:xxii:prop:5.11.5:item:ii}{(ii)}} implies that
\(\mathcal C\ell_c\) is simply the quotient sheaf of
\(\mathcal H_c\) by the equivalence relation defined by
conjugation.

This first implies \textup{\SGARefLink{sga3:xxii:prop:5.11.5:item:iii}{(iii)}}, as well as the fact that it
suffices to verify \textup{\SGARefLink{sga3:xxii:prop:5.11.5:item:i}{(i)}} and \textup{\SGARefLink{sga3:xxii:prop:5.11.5:item:ii}{(ii)}} locally for the
étale topology. We are therefore reduced to the situation of
\textup{\SGARefLink{sga3:xxii:prop:5.11.5:item:iv}{(iv)}}; let us first construct a morphism
\[
  f\colon\mathcal H_c\longrightarrow E_S.
\]
It suffices to construct a map
\(\mathcal H_c(S)\to E_S(S)\), functorial in \(S\); let \(H\) be a
subgroup of \(G\) of type \textup{(RC)}. Since \(H\) locally for
the étale topology has maximal tori, and since the maximal tori of
\(G\) are locally conjugate for the étale topology, there exist a
covering family \((S_i\to S)\), elements \(g_i\in G(S_i)\), and
closed subsets \(R_i\) of \(R\) such that
\[
  \operatorname{int}(g_i)(H\times_S S_i)
  =H_{R_i}\times_S S_i.
\]
Each \(R_i\) defines a section \(\eta_i\) of \(E_{S_i}\), that is,
an element of \(E_S(S_i)\). It remains to prove that the family
\((\eta_i)\) comes from a section \(\eta=f(H)\) of \(E_S\) over
\(S\), and that it depends only on \(H\).

For this it is enough to prove that \(H_{R'}\) and \(H_{R''}\) are
locally conjugate for the étale topology if and only if \(R'\) and
\(R''\) are conjugate by an element of the Weyl group \(W\), which
is clear.

For every \(\eta\in E\), there exists
\(H_0\in\mathcal H_c(S)\) such that \(f(H_0)=\eta\): take
\(H_0=H_{R'}\), where \(R'\) is a closed subset whose image is
\(\eta\). If \(H\in\mathcal H_c(S')\), for \(S'\to S\), then \(H\)
is locally conjugate to \(H_0\) if and only if \(f(H)=\eta\).
Consequently,
\[
  f^{-1}(\eta)\simeq
  G\big/\underline{\operatorname{Norm}}_G(H_0),
\]
which, by \SGARefLink{sga3:xxii:cor:5.8.2}{5.8.2}, is an open subscheme of \(\mathcal H\), smooth,
quasi-projective, and of finite presentation over \(S\), with
nonempty connected fibers. Since \(E_S\) is the disjoint union of
the open subschemes that are the images of the sections
corresponding to the elements \(\eta\in E\), \(\mathcal H_c\)
identifies with the disjoint union of the \(f^{-1}(\eta)\). This
proves \textup{\SGARefLink{sga3:xxii:prop:5.11.5:item:i}{(i)}} and \textup{\SGARefLink{sga3:xxii:prop:5.11.5:item:ii}{(ii)}}; assertion \textup{\SGARefLink{sga3:xxii:prop:5.11.5:item:iv}{(iv)}}
follows from the construction.

\SGARefTarget{sga3:xxii:cor:5.11.6}
\paragraph{\SGARefLink{sga3:xxii:cor:5.11.6}{Corollary 5.11.6}.}
{\itshape Let \(u\in\mathcal C\ell_c(S')\), where \(S'\to S\).
Then \(\operatorname{cl}^{-1}(u)\) is a smooth,
quasi-projective \(S'\)-scheme of finite presentation, with
nonempty connected fibers. It is an open subscheme of
\(\mathcal H_c\) and a ``homogeneous'' scheme under \(G_{S'}\)
\textup{(acting by inner automorphisms)}.

In particular, if
\(H\in\operatorname{cl}^{-1}(u)(S')\), the morphism
\[
  G_{S'}\longrightarrow(\mathcal H_c)_{S'},\qquad
  g\longmapsto\operatorname{int}(g)H
\]
identifies
\[
  G_{S'}\big/
  \underline{\operatorname{Norm}}_{G_{S'}}(H)
  \quad\text{with}\quad
  \operatorname{cl}^{-1}(u).
\]}

\SGARefTarget{sga3:xxii:examples:5.11.7}
\paragraph{\SGARefLink{sga3:xxii:examples:5.11.7}{Examples 5.11.7}.}
There are two canonical sections \(u_t\) and \(u_b\) of
\(\mathcal C\ell_c\). The section \(u_t\) corresponds to maximal
tori, for which \(R'=\varnothing\). The section \(u_b\) corresponds
to Borel subgroups, for which \(R'\) is a system of positive roots.
The schemes
\(\operatorname{cl}^{-1}(u_t)\) and
\(\operatorname{cl}^{-1}(u_b)\) are respectively
\(\underline{\operatorname{Tor}}(G)\) and
\(\underline{\operatorname{Bor}}(G)\), introduced in \SGARefLink{sga3:xxii:cor:5.8.3}{5.8.3}.
Other examples will be given in Exposé XXVI.

\SGARefTarget{sga3:xxii:rem:5.11.8}
\paragraph{\SGARefLink{sga3:xxii:rem:5.11.8}{Remark 5.11.8}.}
One can construct an \(S\)-scheme \(\mathcal C\ell\), of finite
presentation and unramified over \(S\), and a smooth surjective
morphism
\[
  \mathcal H\longrightarrow\mathcal C\ell
\]
with connected geometric fibers, having properties analogous to
those of \SGARefLink{sga3:xxii:prop:5.11.5}{5.11.5}\textup{\SGARefLink{sga3:xxii:prop:5.11.5:item:ii}{(ii)}} and \textup{\SGARefLink{sga3:xxii:prop:5.11.5:item:iii}{(iii)}}.

\SGARefTarget{sga3:xxii:section:6}
\section*{6. The Derived Group}
\addcontentsline{toc}{section}{6. The Derived Group}

\SGARefTarget{sga3:xxii:section:6.1}
\subsection*{\SGAInlineRefTarget{sga3:xxii:subsection:6-1}6.1. Preliminaries}
\addcontentsline{toc}{subsection}{\SGARefLink{sga3:xxii:subsection:6-1}{6.1}. Preliminaries}

In this section, we fix a scheme \(S\), a split \(S\)-group
\((G,T,M,R)\), and a system \(R_+\) of positive roots of \(R\).
We write
\[
  B=B_{R_+},\qquad B^-=B_{R_-},\qquad
  U=B^u,\qquad U^-=(B^-)^u,
\]
and
\[
  \Omega=\Omega_{R_+}=U^-\cdot T\cdot U.
\]

\SGARefTarget{sga3:xxii:par:6.1.1}
\paragraph{\SGARefLink{sga3:xxii:par:6.1.1}{6.1.1}.}
Let \(T'\) denote the subtorus of \(T\) that is the ``image of the
family \(\alpha^*\), \(\alpha\in R\)''; in other words, \(T'\) is
the image of the homomorphism
\[
  (\mathbb G_{m,S})^R\longrightarrow T
\]
defined by
\[
  (z_\alpha)_{\alpha\in R}
  \longmapsto
  \prod_{\alpha\in R}\alpha^*(z_\alpha).
\]
It follows at once that, if \(\Delta\) denotes the set of simple
roots of \(R_+\), the homomorphism
\[
  (\mathbb G_{m,S})^\Delta\longrightarrow T'
\]
defined in the same way is surjective and has finite kernel. If
we identify \(T\) with \(D_S(M)\), then \(T'\) identifies with
\(D_S(M/N)\), where
\[
  N=M\cap V(R^*)^\perp
\]
\textup{(we write \(V(R^*)^\perp\) for the orthogonal complement
of \(V(R^*)\) under the duality between \(V\) and \(V^*\))}.

\SGARefTarget{sga3:xxii:lemma:6.1.2}
\paragraph{\SGARefLink{sga3:xxii:lemma:6.1.2}{Lemma 6.1.2}.}
{\itshape The homomorphism defined by multiplication in \(T\),
\[
  \operatorname{rad}(G)\times_S T'\longrightarrow T,
\]
is an isogeny \textup{(cf.\ \SGARefLink{sga3:xxii:def:4.2.9}{4.2.9})}.}

Indeed, the canonical homomorphism
\[
  \operatorname{rad}(G)\longrightarrow T/T'
\]
comes by duality from the homomorphism of commutative groups
\[
  M\cap V(R^*)^\perp
  \longrightarrow
  M\big/\bigl(M\cap V(R)\bigr),
\]
which is immediately seen to be injective with finite cokernel
\textup{(cf.\ \SGARefLink{sga3:XXI:6.3}{Exposé XXI, 6.3})}.

\SGARefTarget{sga3:xxii:def:6.1.3}
\paragraph{\SGARefLink{sga3:xxii:def:6.1.3}{Definition 6.1.3}.}
We put
\[
  \Omega'=U^-\cdot T'\cdot U;
\]
this is a closed subscheme of
\(\Omega=U^-\cdot T\cdot U\).

\SGARefTarget{sga3:xxii:lemma:6.1.4}
\paragraph{\SGARefLink{sga3:xxii:lemma:6.1.4}{Lemma 6.1.4}.}
{\itshape Let \(\alpha\) be a simple root and let
\(w_\alpha\in\underline{\operatorname{Norm}}_G(T)(S)\) lift
\(s_\alpha\). Then
\[
  \operatorname{int}(w_\alpha)\Omega'\cap\Omega\subset\Omega'.
\]}

It is enough to prove that, if \(g\in\Omega'(S)\) and
\(\operatorname{int}(w_\alpha)g\in\Omega(S)\), then
\(\operatorname{int}(w_\alpha)g\in\Omega'(S)\). By \SGARefLink{sga3:xxii:cor:5.6.8}{5.6.8}, write
\[
  g=a\exp_{-\alpha}(Y)t\exp_\alpha(X)b,
\]
where
\[
  a\in U_{-\widehat\alpha}(S),\qquad
  Y\in\Gamma(S,\mathfrak g^{-\alpha}),\qquad
  t\in T'(S),\qquad
  X\in\Gamma(S,\mathfrak g^\alpha),\qquad
  b\in U_{\widehat\alpha}(S).
\]
We then have
\[
  \operatorname{int}(w_\alpha)g
  =\operatorname{int}(w_\alpha)a\cdot
   \operatorname{int}(w_\alpha)
   \bigl(\exp_{-\alpha}(Y)t\exp_\alpha(X)\bigr)\cdot
   \operatorname{int}(w_\alpha)b.
\]
By \SGARefLink{sga3:xxii:cor:5.6.8}{5.6.8}\textup{\SGARefLink{sga3:xxii:cor:5.6.8:item:iv}{(iv)}},
\[
  \operatorname{int}(w_\alpha)a\in U_{-\widehat\alpha}(S),
  \qquad
  \operatorname{int}(w_\alpha)b\in U_{\widehat\alpha}(S).
\]
It follows that, on putting
\[
  h=\exp_{-\alpha}(Y)t\exp_\alpha(X),
\]
we have the equivalences
\[
  \operatorname{int}(w_\alpha)g\in\Omega(S)
  \quad\Longleftrightarrow\quad
  \operatorname{int}(w_\alpha)h\in\Omega(S),
\]
and
\[
  \operatorname{int}(w_\alpha)g\in\Omega'(S)
  \quad\Longleftrightarrow\quad
  \operatorname{int}(w_\alpha)h\in\Omega'(S).
\]
We are therefore reduced to the case \(g=h\). Since, by
\textup{(\SGARefLink{sga3:xxii:rem:4.1.12}{4.1.12})},
\[
  Z_\alpha\cap\Omega
  =U_{-\alpha}\cdot T\cdot U_\alpha,
  \qquad
  Z_\alpha\cap\Omega'
  =U_{-\alpha}\cdot T'\cdot U_\alpha,
\]
we are reduced to proving the following assertion:
\[
  \operatorname{int}(w_\alpha)h
  \in\bigl(U_{-\alpha}\cdot T\cdot U_\alpha\bigr)(S)
  \quad\Longrightarrow\quad
  \operatorname{int}(w_\alpha)h
  \in\bigl(U_{-\alpha}\cdot T'\cdot U_\alpha\bigr)(S).
\]
But this follows at once from \SGARefLink{sga3:xx:3.12}{Exposé XX, 3.12}, which shows that the
\(T\)-component of \(\operatorname{int}(w_\alpha)h\) is of the form
\[
  t\cdot\alpha^*(z)\in T'(S).
\]

\SGARefTarget{sga3:xxii:lemma:6.1.5}
\paragraph{\SGARefLink{sga3:xxii:lemma:6.1.5}{Lemma 6.1.5}.}
{\itshape For every
\(w\in\underline{\operatorname{Norm}}_G(T)(S)\), there exists an
open subscheme \(V_w\) of \(G\), containing the unit section, such
that
\[
  \operatorname{int}(w)\Omega'\cap V_w\subset\Omega'.
\]}

For each simple root \(\alpha\), choose an element
\(n_\alpha\in\underline{\operatorname{Norm}}_G(T)(S)\) lifting
\(s_\alpha\). For every point \(s\in S\), there exist an open
subscheme \(V\) of \(S\) containing \(s\), an element \(t\in T(V)\),
and, over \(V\), a relation
\[
  w=n_{\alpha_1}\cdots n_{\alpha_p}t,
\]
where the \(\alpha_i\) are simple. We may plainly restrict ourselves
to proving the result when \(V=S\); we argue by induction on \(p\).
If \(p=0\), then \(w\in T(S)\), and we take \(V_w=G\). Suppose,
therefore, that \(w=n_\alpha w'\), where \(w'\) satisfies the
conclusion of the lemma. There is then an open subscheme \(V_{w'}\)
of \(G\), containing the unit section, such that
\[
  \operatorname{int}(w')\Omega'\cap V_{w'}\subset\Omega'.
\]
We may then write
\[
\begin{aligned}
  \operatorname{int}(w)\Omega'
    \cap\operatorname{int}(n_\alpha)V_{w'}\cap\Omega
  &=
  \operatorname{int}(n_\alpha)
  \bigl(\operatorname{int}(w')\Omega'\cap V_{w'}\bigr)\cap\Omega\\
  &\subset
  \operatorname{int}(n_\alpha)\Omega'\cap\Omega
  \subset\Omega',
\end{aligned}
\]
by \SGARefLink{sga3:xxii:lemma:6.1.4}{6.1.4}. We now take
\[
  V_w=\operatorname{int}(n_\alpha)V_{w'}\cap\Omega,
\]
and the proof is complete.

\SGARefTarget{sga3:xxii:lemma:6.1.6}
\paragraph{\SGARefLink{sga3:xxii:lemma:6.1.6}{Lemma 6.1.6}.}
{\itshape There exists an open subscheme \(V_0\) of \(G\),
containing the unit section, such that for every \(S'\to S\),
\[
  U(S')U^-(S')\cap V_0(S')\subset\Omega'(S').
\]}

Let \(n_0\) be an element of
\(\underline{\operatorname{Norm}}_G(T)(S)\) lifting the symmetry
\(w_0\) of the Weyl group, that is, such that
\[
  \operatorname{int}(n_0)U=U^-
  \qquad\text{\textup{(cf.\ \SGARefLink{sga3:XXI:3.6.14}{Exposé XXI, 3.6.14})}}.
\]
Then \(n_0^2\in T(S)\). We show that the open subscheme
\(V_0=V_{n_0}\) of \SGARefLink{sga3:xxii:lemma:6.1.5}{6.1.5} has the required property. Indeed,
\[
\begin{aligned}
  U(S')U^-(S')
  &=
  \operatorname{int}(n_0)
  \bigl(
    \operatorname{int}(n_0)^{-1}U(S')
    \cdot\operatorname{int}(n_0)^{-1}U^-(S')
  \bigr)\\
  &=
  \operatorname{int}(n_0)\bigl(U^-(S')\cdot U(S')\bigr)
  \subset\operatorname{int}(n_0)\Omega'(S').
\end{aligned}
\]
Consequently,
\[
  U(S')U^-(S')\cap V_0(S')
  \subset
  \operatorname{int}(n_0)\Omega'(S')\cap V_0(S')
  \subset\Omega'(S').
\]

\SGARefTarget{sga3:xxii:lemma:6.1.7}
\paragraph{\SGARefLink{sga3:xxii:lemma:6.1.7}{Lemma 6.1.7}.}
{\itshape Consider the morphism
\[
  f:\Omega=U^-\cdot T\cdot U\longrightarrow T/T'
\]
obtained by composing the second projection with the canonical
morphism from \(T\) to \(T/T'\). Then \(f\) is
``generically multiplicative'': there exists an open subscheme
\(V\) of \(\Omega\times_S\Omega\), containing the unit section
\textup{(and hence relatively schematically dense,
\SGARefLink{sga3:xviii:cor:1.3}{Exposé XVIII, 1.3})}, such that for every \(S'\to S\) and every
\((x,y)\in V(S')\), one has \(xy\in\Omega(S')\) and
\[
  f(xy)=f(x)f(y).
\]}

Let \(x,y\) be sections of \(\Omega\) over \(S'\). Write
\[
  x=utv,\qquad y=u't'v',
\]
where
\[
  u,u'\in U^-(S'),\qquad t,t'\in T(S'),\qquad v,v'\in U(S').
\]
Let \(V_0\) be the open subscheme of \SGARefLink{sga3:xxii:lemma:6.1.6}{6.1.6}, and let \(V\) be the
open subscheme of \(\Omega\times_S\Omega\) defined by the condition
\[
  vu'\in V_0(S').
\]
This open subscheme has the required property. Indeed, if
\((x,y)\in V(S')\), then
\[
  xy=(utv)(u't'v')=(ut)(vu')(t'v').
\]
Since \(vu'\in\Omega'(S')\), we obtain
\[
  xy
  \in U^-(S')\,t\Omega'(S')t'\,U(S')
  \subset U^-(S')\,tt'T'(S')\,U(S'),
\]
and therefore \(xy\in\Omega(S')\) and
\[
  f(xy)=f(tt')=f(t)f(t')=f(x)f(y).
\]

\SGARefTarget{sga3:xxii:prop:6.1.8}
\paragraph{\SGARefLink{sga3:xxii:prop:6.1.8}{Proposition 6.1.8}.}
\textit{There exists a group morphism}
\[
  f:G\longrightarrow T/T'
\]
\textit{inducing the canonical projection on \(T\). The kernel
\(\operatorname{Ker}(f)\) of \(f\) is a closed subgroup scheme of
\(G\), smooth over \(S\), with connected fibers. Every group
morphism from \(G\) to an \(S\)-presheaf of commutative groups
that is separated for the fppf topology vanishes on
\(\operatorname{Ker}(f)\).}

The first assertion follows at once from \SGARefLink{sga3:xxii:prop:4.1.11}{4.1.11}. We immediately
have
\[
  \operatorname{Ker}(f)\cap\Omega=\Omega',
\]
which proves that \(\operatorname{Ker}(f)\) is smooth over \(S\)
at every point of the unit section.
By \SGARefLink{sga3:xxii:prop:5.6.9}{5.6.9}\textup{\SGARefLink{sga3:xxii:prop:5.6.9:item:ii}{(ii)}}, every morphism
\(\phi\) from \(G\) to an \(S\)-presheaf of commutative groups
that is separated for the fppf topology vanishes on \(U\) and on
\(U^-\). By \SGARefLink{sga3:xx:2.7}{Exposé XX, 2.7}, \(\phi\) also vanishes on \(T'\), and
hence on \(\Omega'\). Using the notation of \SGARefLink{sga3:xxii:prop:5.7.10}{5.7.10}, we see that the
monoid \(U_1\) is contained in \(\operatorname{Ker}(f)(S)\), which
shows that
\[
  \operatorname{Ker}(f)
  =
  \bigcup_{u\in U_1}u\Omega'.
\]
It follows, on the one hand, that every \(\phi\) as above vanishes
on \(\operatorname{Ker}(f)\), and, on the other hand, that
\(\operatorname{Ker}(f)\) has connected fibers and is therefore
smooth over \(S\), by Exposé \(\mathrm{VI_B}\), 3.10.

\SGARefTarget{sga3:xxii:section:6.2}
\subsection*{\SGAInlineRefTarget{sga3:xxii:subsection:6-2}6.2. The Derived Group of a Reductive Group}
\addcontentsline{toc}{subsection}{\SGARefLink{sga3:xxii:subsection:6-2}{6.2}. The Derived Group of a
Reductive Group}

\SGARefTarget{sga3:xxii:thm:6.2.1}
\paragraph{\SGARefLink{sga3:xxii:thm:6.2.1}{Theorem 6.2.1}.}
\textit{Let \(S\) be a scheme and \(G\) a reductive \(S\)-group.}
\begin{enumerate}
\item[\textup{\SGAInlineRefTarget{sga3:xxii:thm:6.2.1:item:i}(i)}]
\textit{The group}
\[
  \mathrm D_S(G)
  =
  \underline{\operatorname{Hom}}_{S\text{-gr.}}
    (G,\mathbb G_{m,S})
\]
\textit{is represented by a twisted constant \(S\)-group whose
type at \(s\in S\) is}
\[
  \mathbb Z^{
    \operatorname{rgred}(G_s)-\operatorname{rgss}(G_s)}.
\]

\item[\textup{\SGAInlineRefTarget{sga3:xxii:thm:6.2.1:item:ii}(ii)}]
\textit{Write}
\[
  \operatorname{corad}(G)
  =
  \mathrm D_S\bigl(\mathrm D_S(G)\bigr),
\]
\textit{which is therefore an \(S\)-torus. The biduality morphism
\textup{(cf.\ Exposé VIII, \S1)}}
\[
  f_0:G\longrightarrow\operatorname{corad}(G)
\]
\textit{is smooth and surjective.}

\item[\textup{\SGAInlineRefTarget{sga3:xxii:thm:6.2.1:item:iii}(iii)}]
\textit{The composite morphism}
\[
  \operatorname{rad}(G)\longrightarrow G
  \longrightarrow\operatorname{corad}(G)
\]
\textit{is an isogeny \textup{(cf.\ \SGARefLink{sga3:xxii:def:4.2.9}{4.2.9})}.}

\item[\textup{(iv)}]
\textit{The kernel of \(f_0\), denoted}
\[
  \operatorname{der}(G)=\operatorname{Ker}(f_0),
\]
\textit{is a closed subgroup scheme of \(G\), semisimple over
\(S\), called the derived group of \(G\). If \(G\) is semisimple,
then \(\operatorname{der}(G)=G\).}

\item[\textup{\SGAInlineRefTarget{sga3:xxii:thm:6.2.1:item:v}(v)}]
\textit{Every group morphism from \(G\) to an \(S\)-presheaf of
commutative groups that is separated for the fppf topology
vanishes on \(\operatorname{der}(G)\), and therefore factors
through \(f_0\).}
\end{enumerate}

\textit{Proof.}
All the assertions of the theorem are local for the étale
topology; we may therefore reduce to the case in which \(G\) is
split over \(S\). Consider the morphism \(f\) of \SGARefLink{sga3:xxii:prop:6.1.8}{6.1.8}. By the
last assertion of \SGARefLink{sga3:xxii:prop:6.1.8}{6.1.8}, we immediately have an isomorphism
\[
  \underline{\operatorname{Hom}}_{S\text{-gr.}}
    (G,\mathbb G_{m,S})
  \xrightarrow{\ \sim\ }
  \underline{\operatorname{Hom}}_{S\text{-gr.}}
    (T/T',\mathbb G_{m,S}),
\]
which proves \textup{\SGARefLink{sga3:xxii:thm:6.2.1:item:i}{(i)}}, then \textup{\SGARefLink{sga3:xxii:thm:6.2.1:item:ii}{(ii)}}, and yields a
commutative diagram
\[
  \begin{tikzcd}[column sep=5.0em,row sep=3.2em]
  G
    \arrow[r,"f_0"]
    \arrow[dr,"f"']
  &
  \operatorname{corad}(G)
    \arrow[d,"\sim"]
  \\
  &
  T/T'.
\end{tikzcd}

\]
Assertion \textup{\SGARefLink{sga3:xxii:thm:6.2.1:item:v}{(v)}} now follows from \SGARefLink{sga3:xxii:prop:6.1.8}{6.1.8}, and
\textup{\SGARefLink{sga3:xxii:thm:6.2.1:item:iii}{(iii)}} from \SGARefLink{sga3:xxii:lemma:6.1.2}{6.1.2}. We also have
\(\operatorname{Ker}(f)=\operatorname{Ker}(f_0)\), which, by
\SGARefLink{sga3:xxii:prop:6.1.8}{6.1.8}, implies that \(\operatorname{der}(G)\) is smooth over
\(S\) with connected fibers. It remains to verify that its fibers
are semisimple. They are reductive by \SGARefLink{sga3:xix:paragraph:1-7}{Exposé XIX, 1.7}, as invariant
subgroups of reductive groups. By \textup{\SGARefLink{sga3:xxii:thm:6.2.1:item:iii}{(iii)}},
\(\operatorname{rad}(G)\cap\operatorname{der}(G)\) is finite,
which indeed implies that the fibers of \(\operatorname{der}(G)\)
are semisimple.

\SGARefTarget{sga3:xxii:rem:6.2.2}
\paragraph{\SGARefLink{sga3:xxii:rem:6.2.2}{Remark 6.2.2}.}
\SGAInlineRefTarget{sga3:xxii:rem:6.2.2:item:a}\textup{a)} By construction, when \(G\) is split,
\(\operatorname{der}(G)\) is the fppf subsheaf of \(G\) generated
by the \(U_\alpha\), \(\alpha\in R\). It even suffices to take the
\(U_\alpha\), \(\alpha\in\Delta\cup-\Delta\), where \(\Delta\) is
a base of \(R\).

\textup{b)}\footnote{Editor's note: In what follows, we have
expanded the original and removed the assertion that
``\(\operatorname{der}(G)\) is the fppf-separated presheaf of
commutators of \(G\).''}
Let \(C\) be the presheaf of commutators of \(G\), that is, the
\(S\)-functor in groups that assigns to every \(S'\to S\) the
commutator subgroup of \(G(S')\) (that is, the subgroup of
\(G(S')\) generated by the elements
\(xyx^{-1}y^{-1}\), for \(x,y\in G(S')\)), and let
\(\widetilde C\) be the associated fppf sheaf. Since the quotient
\(G/\operatorname{der}(G)=T/T'\) is commutative,
\(\operatorname{der}(G)\) contains \(C\), and hence
\(\widetilde C\) \textup{(cf.\ \SGARefLink{sga3:iv:corollary:4:3:12}{Exposé IV, 4.3.12})}.

On the other hand, the quotient presheaf \(G/\widetilde C\) is
separated \textup{(\SGARefLink{sga3:iv:lemma:4:4:8:1}{Exposé IV, 4.4.8.1})}, and therefore, by
\textup{(v)}, we have
\(\operatorname{der}(G)\subset\widetilde C\). Thus
\(\operatorname{der}(G)=\widetilde C\); in other words,
\(\operatorname{der}(G)\) is the fppf sheaf of commutators of
\(G\).

Let us finally note that, since \(C\) is a subpresheaf of \(G\),
it is separated, but it is not equal to
\(\operatorname{der}(G)\) in general: for example,
\(\operatorname{der}(\mathrm{SL}_2)=\mathrm{SL}_2\), whereas
\(\mathrm{SL}_2(\mathbb F_2)\simeq\mathfrak S_3\) is not equal
to its commutator subgroup.

\textup{c)} When \(S\) is the spectrum of an algebraically closed
field \(k\), \(\operatorname{der}(G)(k)\) is the commutator
subgroup of \(G(k)\) \textup{(Exposé \(\mathrm{VI_B}\), 7.10)}.

\SGARefTarget{sga3:xxii:par:6.2.3}
\paragraph{\SGARefLink{sga3:xxii:par:6.2.3}{6.2.3}.}
Consider now the two exact sequences
\[
\begin{gathered}
  1\longrightarrow\operatorname{rad}(G)
  \longrightarrow G
  \longrightarrow\operatorname{ss}(G)
  \longrightarrow 1,\\
  1\longrightarrow\operatorname{der}(G)
  \longrightarrow G
  \longrightarrow\operatorname{corad}(G)
  \longrightarrow 1.
\end{gathered}
\]
Since \(\operatorname{rad}(G)\) is central in \(G\),
multiplication in \(G\) defines a group morphism
\[
  u:\operatorname{rad}(G)\times_S\operatorname{der}(G)
  \longrightarrow G,
\]
which is covering by \SGARefLink{sga3:xxii:thm:6.2.1}{6.2.1}\textup{\SGARefLink{sga3:xxii:thm:6.2.1:item:iii}{(iii)}}, and hence surjective and
flat \textup{(Exposé \(\mathrm{VI_B}\), 9.2\textup{(xi)})}.
Its kernel is isomorphic to
\(\operatorname{rad}(G)\cap\operatorname{der}(G)\), which is also
the kernel of
\[
  \operatorname{rad}(G)\longrightarrow\operatorname{corad}(G);
\]
it is a finite subgroup of multiplicative type of
\(\operatorname{rad}(G)\). Similarly, the kernel of the morphism
\[
  G\longrightarrow
  \operatorname{corad}(G)\times_S\operatorname{ss}(G)
\]
is isomorphic to
\(\operatorname{der}(G)\cap\operatorname{rad}(G)\).

\SGARefTarget{sga3:xxii:prop:6.2.4}
\paragraph{\SGARefLink{sga3:xxii:prop:6.2.4}{Proposition 6.2.4}.}
\textit{Let \(G\) be a reductive \(S\)-group. The three morphisms}
\[
\begin{gathered}
  \operatorname{rad}(G)\times_S\operatorname{der}(G)
    \longrightarrow G,\\
  G\longrightarrow
    \operatorname{corad}(G)\times_S\operatorname{ss}(G),\\
  \operatorname{rad}(G)\longrightarrow\operatorname{corad}(G)
\end{gathered}
\]
\textit{are central isogenies, and their kernels are isomorphic.}

\SGARefTarget{sga3:xxii:cor:6.2.5}
\paragraph{\SGARefLink{sga3:xxii:cor:6.2.5}{Corollary 6.2.5}.}
\textit{The following conditions are equivalent:}
\begin{enumerate}
\item[\textup{(i)}]
\textit{\(G\) is the product of a semisimple group and a torus.}

\item[\textup{(ii)}]
\textit{The morphism}
\[
  \operatorname{rad}(G)\times_S\operatorname{der}(G)
  \xrightarrow{\ \sim\ }G
\]
\textit{is an isomorphism.}

\item[\textup{(iii)}]
\textit{The morphism}
\[
  G\xrightarrow{\ \sim\ }
  \operatorname{corad}(G)\times_S\operatorname{ss}(G)
\]
\textit{is an isomorphism.}

\item[\textup{(iv)}]
\textit{\(\operatorname{rad}(G)\cap\operatorname{der}(G)=e\).}
\end{enumerate}

\SGARefTarget{sga3:xxii:par:6.2.6}
\paragraph{\SGARefLink{sga3:xxii:par:6.2.6}{6.2.6}.}
Let us return temporarily to the case of a split group. We retain
the notation of \SGARefLink{sga3:xxii:subsection:6-1}{6.1}. Put
\[
  N=M\cap V(R^*)^\perp,
  \qquad\text{so that}\qquad
  T'=\mathrm D_S(M/N).
\]
Then \(U^-\cdot T'\cdot U\) is an open neighborhood of the unit
section in \(\operatorname{der}(G)\). Consequently,
\[
  \operatorname{Lie}\bigl(\operatorname{der}(G)/S\bigr)
  =
  \mathfrak t'\oplus
  \coprod_{\alpha\in R}\mathfrak g^\alpha.
\]
Since the characters induced on \(T'\) by the
\(\alpha\in R\) are nonzero and distinct
\textup{(\SGARefLink{sga3:XXI:1.2.5}{Exposé XXI, 1.2.5}---which we already used in \SGARefLink{sga3:xxii:lemma:6.1.2}{6.1.2})},
\(R\) is a root system of \(\operatorname{der}(G)\) relative to
\(T'\). It is then immediate \textup{(since
\(U_\alpha\subset\operatorname{der}(G)\))} that the
\(\exp\)-morphisms of \(\operatorname{der}(G)\) ``are'' those of
\(G\), and likewise for the coroots.

It follows that:

\SGARefTarget{sga3:xxii:prop:6.2.7}
\paragraph{\SGARefLink{sga3:xxii:prop:6.2.7}{Proposition 6.2.7}.}
\textit{With the preceding notation,}
\[
  \bigl(\operatorname{der}(G),T',M/N,R\bigr)
\]
\textit{is a split group with root datum
\(\operatorname{der}\bigl(\mathcal R(G)\bigr)\). The canonical
morphism \(\operatorname{der}(G)\to G\) induces, by functoriality,
the canonical morphism of root data}
\[
  \mathcal R(G)\longrightarrow
  \operatorname{der}\bigl(\mathcal R(G)\bigr)
\]
\textit{of \SGARefLink{sga3:XXI:6.5}{Exposé XXI, 6.5}.}

\textit{N.B.} As an exercise, the reader may construct the diagram
of split groups corresponding to the three left-hand columns of
the diagram of root data in \SGARefLink{sga3:XXI:6.5.7}{Exposé XXI, 6.5.7}.

\SGARefTarget{sga3:xxii:prop:6.2.8}
\paragraph{\SGARefLink{sga3:xxii:prop:6.2.8}{Proposition 6.2.8}.}
\textit{Let \(S\) be a scheme, \(G\) a reductive \(S\)-group, and
\(\operatorname{der}(G)\) its derived group.}
\begin{enumerate}
\item[\textup{\SGAInlineRefTarget{sga3:xxii:prop:6.2.8:item:i}(i)}]
\textit{For every maximal torus \(T\) of \(G\),
\(T\cap\operatorname{der}(G)\) is a maximal torus of
\(\operatorname{der}(G)\). For every maximal torus \(T'\) of
\(\operatorname{der}(G)\),}
\[
  \underline{\operatorname{Centr}}_G(T')
  =
  \operatorname{rad}(G)\cdot T'
\]
\textit{is a maximal torus of \(G\). These two constructions are
inverse to one another and establish a bijective correspondence
between the maximal tori of \(G\) and those of
\(\operatorname{der}(G)\).}

\item[\textup{\SGAInlineRefTarget{sga3:xxii:prop:6.2.8:item:ii}(ii)}]
\textit{For every Borel subgroup \(B\) of \(G\),
\(B\cap\operatorname{der}(G)\) is a Borel subgroup \(B'\) of
\(\operatorname{der}(G)\). We have \(B'{}^u=B^u\). For every
Borel subgroup \(B'\) of \(\operatorname{der}(G)\),}
\[
  \underline{\operatorname{Norm}}_G(B')
  =
  \operatorname{rad}(G)\cdot B'
\]
\textit{is a Borel subgroup of \(G\). These constructions are
inverse to one another and establish a bijective correspondence
between the Borel subgroups of \(G\) and those of
\(\operatorname{der}(G)\).}
\end{enumerate}

By the local conjugacy theorem for maximal tori and the
construction of the derived group, the only assertion that
remains to be proved in \textup{\SGARefLink{sga3:xxii:prop:6.2.8:item:i}{(i)}} is the following: if \(T\)
is a maximal torus of \(G\), then
\[
  T
  =
  \bigl(T\cap\operatorname{der}(G)\bigr)\cdot
    \operatorname{rad}(G)
  =
  \underline{\operatorname{Centr}}_G
    \bigl(T\cap\operatorname{der}(G)\bigr).
\]
The first equality is immediate \textup{(since we may reduce to
the split case)}; the second follows at once, because
\(\operatorname{rad}(G)\) is central in \(G\), and hence
\[
  T
  =
  \underline{\operatorname{Centr}}_G(T)
  =
  \underline{\operatorname{Centr}}_G
    \bigl(T\cap\operatorname{der}(G)\bigr).
\]
The argument for \textup{\SGARefLink{sga3:xxii:prop:6.2.8:item:ii}{(ii)}} is the same.

\SGARefTarget{sga3:xxii:section:6.3}
\subsection*{\SGAInlineRefTarget{sga3:xxii:subsection:6-3}6.3. Subgroups with Commutative Quotients}
\addcontentsline{toc}{subsection}{\SGARefLink{sga3:xxii:subsection:6-3}{6.3}. Subgroups with Commutative
Quotients}

\SGARefTarget{sga3:xxii:par:6.3.1}
\paragraph{\SGARefLink{sga3:xxii:par:6.3.1}{6.3.1}.}
Let \(G\) be a reductive \(S\)-group. If \(H\) is a subsheaf in
groups of \(G\), the following conditions are equivalent:
\begin{itemize}
\item \(H\) contains \(\operatorname{der}(G)\);
\item \(H\) is normal and \(G/H\) is commutative.
\end{itemize}
In this case, the canonical morphism
\[
  f_0:G\longrightarrow\operatorname{corad}(G)
\]
maps \(H\) onto a subsheaf \(f_0(H)\) of
\(\operatorname{corad}(G)\), and we have
\[
\begin{gathered}
  G/H\simeq\operatorname{corad}(G)/f_0(H),
  \qquad
  H/\operatorname{der}(H)\simeq f_0(H),\\
  \operatorname{der}(G)=\operatorname{der}(H),
  \qquad
  H=f_0^{-1}\bigl(f_0(H)\bigr).
\end{gathered}
\]

Since \(\operatorname{der}(G)\) is smooth over \(S\) with
connected fibers,\footnote{Editor's note:
We have expanded the original in what follows. In particular, we
have added the conclusion, implicit in the original, that every
subgroup of \(G\) that is smooth with connected fibers and contains
\(\operatorname{der}(G)\) is closed in \(G\).} it follows from
\SGARefLink{sga3:iv:proposition:5:3:1}{Exposé IV, 5.3.1} and \SGARefLink{sga3:xxii:par:6.3.1}{6.3.1}, and Exposé \(\mathrm{IV_B}\), 9.2,
that the assignment
\[
  H\longmapsto f_0(H)
\]
establishes a bijective correspondence between the group
subschemes \textup{(respectively, the closed group subschemes)} of
\(G\) that contain \(\operatorname{der}(G)\), are smooth over
\(S\), and have connected fibers, and the group subschemes
\textup{(respectively, the closed group subschemes)} of
\(\operatorname{corad}(G)\) that are smooth over \(S\) and have
connected fibers.

Now, if \(H'\) is a group subscheme of
\(\operatorname{corad}(G)\), smooth over \(S\) with connected
fibers, then \(H'\) is of finite presentation over \(S\)
\textup{(Exposé \(\mathrm{VI_B}\), \SGARefLink{sga3:xxii:subsection:5-5}{5.5})}, and its fibers are tori
\textup{(because those of \(\operatorname{corad}(G)\) are)}.
Hence, by \SGARefLink{sga3:x:corollary:8-2}{Exposé X, 8.2}, \(H'\) is a subtorus of
\(\operatorname{corad}(G)\), and is therefore closed in
\(\operatorname{corad}(G)\) \textup{(\SGARefLink{sga3:ix:remark:2:6}{Exposé IX, 2.6})}.

Consequently, every subgroup of \(G\) that is smooth with
connected fibers and contains \(\operatorname{der}(G)\) is closed
in \(G\).

\SGARefTarget{sga3:xxii:par:6.3.2}
\paragraph{\SGARefLink{sga3:xxii:par:6.3.2}{6.3.2}.}
If \(H\) is a closed group subscheme of \(G\), smooth over \(S\),
with connected fibers, and normal in \(G\), then \(H\) is
reductive. If, moreover, \(H\supset\operatorname{der}(G)\), then
\[
  \operatorname{der}(H)=\operatorname{der}(G)
  \qquad\text{and}\qquad
  f_0(H)\simeq\operatorname{corad}(H).
\]
We have therefore proved the following proposition.

\SGARefTarget{sga3:xxii:prop:6.3.3}
\paragraph{\SGARefLink{sga3:xxii:prop:6.3.3}{Proposition 6.3.3}.}
\textit{Let \(G\) be a reductive \(S\)-group. Every group
subscheme \(H\) of \(G\) that is normal in \(G\), has commutative
quotient \textup{(that is, contains
\(\operatorname{der}(G)\))}, is smooth over \(S\), and has
connected fibers}
\textit{is closed and reductive. We have}
\[
  \operatorname{der}(H)=\operatorname{der}(G),
  \qquad
  f_0(H)\simeq\operatorname{corad}(H),
\]
\textit{and}
\[
  G/H\simeq
  \operatorname{corad}(G)/\operatorname{corad}(H),
  \qquad
  H=
  \bigl(H\cap\operatorname{rad}(G)\bigr)
  \cdot\operatorname{der}(G).
\]
\textit{Moreover, the assignment \(H\mapsto f_0(H)\) defines a
bijection between the set of subgroups \(H\) of \(G\) having the
preceding properties and the set of subtori of
\(\operatorname{corad}(G)\).}\footnote{Translator's note: The French
source prints \(\operatorname{corad}(H)\) here; the codomain of
\(f_0\) and the preceding correspondence require
\(\operatorname{corad}(G)\).}

By another application of Noether's theorem \textup{(\SGARefLink{sga3:iv:proposition:5:3:1}{Exposé IV,
5.3.1} and \SGARefLink{sga3:xxii:par:6.3.1}{6.3.1})}, we deduce the following.

\SGARefTarget{sga3:xxii:prop:6.3.4}
\paragraph{\SGARefLink{sga3:xxii:prop:6.3.4}{Proposition 6.3.4}.}
\textit{Let \(S\) be a scheme, \(G\) a reductive \(S\)-group, and
\(T\) a maximal torus of \(G\). For every subgroup \(H\) of \(G\)
as above, \(T\cap H\) is a maximal torus of
\(H\),}\footnote{Translator's note: The French source prints \(G\)
here; since \(T\cap H\subset H\), and by the formulas that follow,
the intended group is \(H\).}
\textit{and}
\[
  G/H\simeq T/(T\cap H),
  \qquad
  H=(T\cap H)\cdot\operatorname{der}(G).
\]
\textit{Moreover, the assignment \(H\mapsto T\cap H\) is a
bijection between the set of subgroups \(H\) of \(G\) as above
and the set of subtori of \(T\) containing
\(T\cap\operatorname{der}(G)\).}

\clearpage
\SGARefTarget{sga3:xxii:bibliography}
\section*{Bibliography}
\addcontentsline{toc}{section}{Bibliography}

\begin{description}
\SGARefTarget{sga3:xxii:bib:BLie}
\item[{[BLie]}]
N.~Bourbaki, \emph{Groupes et algèbres de Lie}, Chapters
IV--VI, Hermann, 1968.

\SGARefTarget{sga3:xxii:bib:Ch05}
\item[{[Ch05]}]
C.~Chevalley, \emph{Classification des groupes algébriques
semi-simples} \textup{(with the collaboration of P.~Cartier,
A.~Grothendieck, and M.~Lazard)}, \emph{Collected Works},
vol.~3, Springer, 2005.

\SGARefTarget{sga3:xxii:bib:DG70}
\item[{[DG70]}]
M.~Demazure and P.~Gabriel, \emph{Groupes algébriques},
Masson \& North-Holland, 1970.

\SGARefTarget{sga3:xxii:bib:Gi71}
\item[{[Gi71]}]
J.~Giraud, \emph{Cohomologie non abélienne},
Springer-Verlag, 1971.

\SGARefTarget{sga3:xxii:bib:Se64}
\item[{[Se64]}]
J.-P.~Serre, \emph{Cohomologie galoisienne},
Springer-Verlag, 1964; 5th ed., 1994.
\end{description}
